
%%%%%%%%%%%%%%%%%%%%%%%%%%%%%%%% SCRIPT FOR E-RARA AND MDZ FILES     %%%%%%%%%%%%%%%%%%%%%%%%%%%%%%%%%%%%%%%%%%%%%%%%
%%%%%%%%%%%%%%%%%%%%%%%%%%% fini le 30.04.2025 par F. GOY            %%%%%%%%%%%%%%%%%%%%%%%%%%%%%%%%%%%%%%%%%%%%%%%%
% !TeX TS-program = lualatex
\documentclass{article}
\usepackage[T1]{fontenc}
\usepackage{microtype}
\usepackage[pdfusetitle,hidelinks]{hyperref}

\usepackage{polyglossia}
\setmainlanguage{english}
\setotherlanguages{latin,greek}
\usepackage[series={},nocritical,noend,noeledsec,nofamiliar,noledgroup]{reledmac}
\usepackage{reledpar}

\usepackage{fontspec}
\setmainfont{TeX Gyre Termes}

\usepackage{sectsty}
\usepackage{xcolor}

\usepackage{fancyhdr}
\pagestyle{fancy}
\fancyhf{}
\fancyhead[LE,RO]{\nouppercase{\leftmark}}  
\cfoot{\thepage}
\renewcommand{\headrulewidth}{0.4pt}

% Redefine \section to remove numbering
\usepackage{titlesec}
\titleformat{\section}[block]{\normalfont\scshape\color{gray}}{}{0pt}{} % no number in heading
\titleformat{\subsection}[hang]{\normalfont}{}{0pt}{} % also remove subsection number
\titleformat{\subsubsection}[hang]{\normalfont\footnotesize\color{black}}{}{0pt}{}

% Modify how section marks are stored to exclude numbers
\makeatletter
\renewcommand{\sectionmark}[1]{%
	\markboth{#1}{}} % Only store the section title, without number
\renewcommand{\subsectionmark}[1]{%
	\markright{#1}} % Only store the subsection title, without number
\renewcommand{\numberline}[1]{} % Hide the section number in TOC
\makeatother

\begin{document}

\date{}
        \title{Commentarius in priorem Timothei epistolam : [Anonymus], [1533]}
\maketitle
\tableofcontents
\clearpage
\begin{pages} 
\beginnumbering
        \pstart COMMENTARIVS IN PRIOREM TIMOTHEI EPISTOLAM a uiro summae pietatis studio conscriptus.  \pend\pstart BASILEAE EXCVDEBAT HENRICVS PETRVS.  \pend
\marginpar{[ p.3 ]}
\phantomsection
\addcontentsline{toc}{subsection}{\textit{COMMENTARIVS IN PRIOREM TIMOTHEI EPISTOLAM a uiro summae pietatis studio conscriptus. }}
\subsection*{\textit{COMMENTARIVS IN PRIOREM TIMOTHEI EPISTOLAM a uiro summae pietatis studio conscriptus. }}
\section{CAPVT I.}
\phantomsection
\addcontentsline{toc}{subsection}{\textit{\huge\textbf{P}\normalsize Aulus apostolus Iesu christi. }}
\subsection*{\textit{\huge\textbf{P}\normalsize Aulus apostolus Iesu christi. }}\pstart Regale et laudatum apud Iudaeos no men deposuit, inuidiam Christianis anne xam habens, quod Saul Dauid regem sa crum et electum perpetuo persecutus scribitur, ne Pau lus Christi persecutoris nomine ipso, sicut re, et ope re no torius habitus, fuerat, perpetuo uocaretur, mu tauit Sault nomem et humile delegit Pauli nomem, si la tine interpreteris, uel operatorium si Hebraice, qui tantum de suis laboribus gloriari solitus est. Reuera autem Paulus est, aliis literis Hebraicis mirabilis in uer bis et operibus. Nostri pontifices electi, humile nomem ponunt, et imperialia assumunt, ex Ioanne Leo. ex Gibriele Alexander, ex Iuliano Iulius, uastator ex ua statore. Curanda nomina et cognomina, ne uel ridicula uel obscoena infantibus obueniant unde uiri ali quando erubelc ant, maxime doctores honesta et ad uirtutes commonitoria nonina habere non est inutile, quod et ethnicis olim curae fuit, et sanctis patriar\pend
\section*{COMMEN. IN PRIOREM }
\marginpar{[ p.4 ]}\pstart chis, non parum refert honesto uel ridiculo nomine uocari apud uulgum.  \pend\pstart Apostolus IA Christo iam nomen aoeperat, uas electi onis erat, reuera doctor gentium in fide et ueritate erat, ecclesiarum pastor. Sed seipsum, et officio legatum fatetur, ni quid altius uelit nominari, ne se ho norem querere. sedofficium, legationem agere teste tur, plenum laboribus, periculis, cura, non otiosum, gloriosum saeculare, quale nunc sectantur multi qui apostolici uocari diligunt, reuera regibus mundani ores, et a Christo alienissima conuersatione, aposta tic potius iudicandi.  \pend\pstart Iesu Christil Cuius nunc nomen celebrat, in omnibus epistolis incessanter, cuius inuocatores iam an te persequutus fuerat, ad mortem, nunc ipsum nomen portat coram gentibus, regibus et filiis Israë in quo solo seruamur, nomen nouum quod os domi ni nominat. Saluator et seruator est omnium sibi credentium, qui ut Iosue, inducit ad terram promissi onis, quod non poterat Moses, qui de captiuitate li berat Babilonica, qui semper orat pro populo dei. Qui unctus est rex regum, et dominus dominantium qui sacerdos magnus semper astat pro nobis, qui plenus gratia et ueritate, filios habet omnets unctos, a quo nomen nostrum consequimur, ut Chrisii  \pend
\section*{TIMOTHEI EPISTOLAM. }
\marginpar{[ p.5 ]}\pstart ni nominemur et simus.  \pend
\phantomsection
\addcontentsline{toc}{subsection}{\textit{Secundum imperium dei et saluatoris nostri, et Iesu Christi spei nostrae. }}
\subsection*{\textit{Secundum imperium dei et saluatoris nostri, et Iesu Christi spei nostrae. }}\pstart Erasmus: Secundum delegationem dei seruatoris nostri, et domini Iesu Christi, qui est spes nostra. Misit me recalcitrantem contrastimulum et legatum ordinauit, nihil tale cogtantem, omnipotens deus, iam mundum seruare uolens, iam conciliatus per Christi mysterium, qui nunc seruatorem se exhi bet piissimum, olim seuerus, et iudex durus, populo etiam suo peculiari. Is me misit, eius imperio subtr ahe re me non possum, eidem seruio, tibi, et prote Chri stianis scribens, non meo nomine, sed eius praecepto et uerbis, qui in me loquitur, et scribit, ne humanum putetur et mendax aut fabula inutilis, sed uerbum dei uiuum.  \pend\pstart Et domini nostri Iesu Christiq Per excellentiam a nobis dominus agnoscendus, qui nos eripuit de seruitio et iugo Sathanae, et error um et peaato rum. Ei serui datisumus, a deo patre, et per eum fi£ lii dei facti et cohaeredes secum, per suum sanguinem et mortem. Non solum omnipotentis dei et creatoris, sed et Christi dominium agnoscimus, ne uel Iudaeis uel paganis scrupulum ingeramus, si uel domini caeli uel Christum filium non conftemur.  \pend
\section*{COMMEN. IN PRIOREM }
\marginpar{[ p.6 ]}
\phantomsection
\addcontentsline{toc}{subsection}{\textit{Qui est spes nostra. }}
\subsection*{\textit{Qui est spes nostra. }}\pstart Vnica fiduciae nostrae anchora, fidelis in promissis qui pro nobis satisfecit, qui gratis nos saluos faciet credentes. Sine quo mera desperatio, sicut ex Eua mors et infernus nos manent. Ipse enim solus et nemo alius. Est uia, ueritas, et uita nostra, sicut et spes nostra. Maledictus omnis qui sperat in alium a Christo. Ipse est spes nostra. Si quis iactat spem suam in Christi dignissima matre et dei genitrice Maria, is uiderit, quid intelligat, credat, uel confiteatur, ne Christi solius gloriam alteri tradet, ne creatori hominem praeponat, ne matrem ex gratia Christi Christo praeferat, natura filio dei, id est, genuit nobis Christum, spem nostram, in quem et ipsa sperans et credens saluari debuit, cuius benignitati nihil addi potest, qui fons est bonitatis, et plenus gratia, a quo et mater eius saluata est, et omnis qui saluatur. Vel bene intelligat et bene loquatur et uelit, uel taceas a blasphemia Christi, quam unicam, mater, odit, et auersabitur semper. Christus solus clamat. Venite ad me qui laboratis et onerati estis, ego reficiam uos. Non hoc matri detulit, sicut nec minora quidem. Quid enim mihi et tibi est (inquit) mulier etc.  \pend
\section*{TIMOTHEI EPISTOLAM. }
\marginpar{[ p.7 ]}
\phantomsection
\addcontentsline{toc}{subsection}{\textit{Timotheo dilecto filio infide. }}
\subsection*{\textit{Timotheo dilecto filio infide. }}\pstart Timotheus interpretatur lingua sua (Graecus enim ex patre, licet matre Iudaea et Christana) ho norans deum, Τιμόςetet Eoee honorabilis uirihonorificum nomen quoque. Nam et Tiret honorabilem notat. Habebat bonum testimonium a fratribus familiaribus, erat eruditus in lege et discipli nis saecularibus adhuc iuuenis, Paulo gratus et sodalis, ideo eum circuncidi uoluit, ut pro tempore Iudaeis superstitiosis gratificaret, et sandalum ampu£ taret. Factus est episcopus in Asia EPhesi, taleis decet esse episcopos, bono testimonio, continentes, doctos, fideleis, fama et uita uenerabiles, ne domini no men blasphemetur.  \pend\pstart Dilecto IErat enim uirtuosus, et fidelis, et amo re Aposteli dignus. Dilectio in minus se habet quam amor et homines respcit. Amor et charitas in plus et ad deum, sicut lieb et minn quod charitatem explicat, elegantissimum nomen, quod foeda cupido et abusus lubricae linguae uertit in pudendum.  \pend\pstart Filiol Qui paternas dotes, omnium uirtutum ut germanus filius referebat, ideo bene Eras. retulit Germano filio, cuius Paulum non solum non puderet, sed insignem sectatorem, sed sodalem et coapostolum fecit. Sic gaudent patres de filiis sensatis, sic doctores  \pend
\section*{COMMEN. IN PRIOREM }
\marginpar{[ p.8 ]}\pstart de discipulis studiosis sibi gratulantur. Sic dulciter al lo quendi sunt filii et discipuli et auditores. Contra ludi magistros ferocissimos.  \pend\pstart In fide NNe carne genitum aestimares, quae et nobi liorfiliatio, sicut fides naturae praefertur ad dominum Sic pater generando similem, naturam communicat Sic Paulus huic filio fidem in Christum contradidit et omnia, non sterilem, sed foecundissimam quae po£ pulos Christo plurimos genuit, et nutriuit uerbo ui tae et cuangelica doctrina. Fidem prae dotibus omni bus commendat in discipulo Apostolus, non circun cisionem, non scientiam legis, non parentis fides facit Christianum, et Christianorum antistitem et dei filium adoptiuum.  \pend
\phantomsection
\addcontentsline{toc}{subsection}{\textit{Gratia, misericordia et pax. }}
\subsection*{\textit{Gratia, misericordia et pax. }}\pstart Optat non diuitias, non sanitatem, non honores, quae hodie sectantur, apostolici, sed dei optimi fauorem et gratiam, qua nobis beneficus est, intra nos bona opera operatur, amat ut filios dilectissimos, immeritis munera spiritualia largitur, et tempora libus non destituit. Gratia haec non est habitus aliquis qualitatis in animam, influxus, stertens, et otiosus, do nec aetatis profectum suscitet libero arbitrio ad bonum operandum, et ad imitatorios actus, sed est dei fauor et insinctus, ad opera dei in nobis, quem ut  \pend
\section*{TIMOTHEI EPISTOLAM. }
\marginpar{[ p.9 ]}\pstart summum bonum in hac uita nobis appetere debemus quem et Paulus imprecatur primum. Si enim deus pro nobis, quis contra nos.  \pend\pstart Misericordia) Pius dei animus, quo nostris ma lis non offenditur, tollerat, propitiatur blaspheman tibus, benefacit, ignorantibus parcit, errantes redu cit, hanc nobis innotescere cupit semper, et pro hoc sperare docet.  \pend\pstart PaxI Conscientiae certo confidentis, de gratia dei et misericordia, exhilarans animum, similiter et externam pacem pro dei misericordia optat, ut libere deo uacari possit, hoc primum impensissime optabatur a Christi discipulis, ut lupos periclitantibus, de uita, et praedicatione ueritatis oportuna sa cro osculo finis synaxcos imprecabatur haec pax.  \pend
\phantomsection
\addcontentsline{toc}{subsection}{\textit{A deo patre et Christo, domino nostro. }}
\subsection*{\textit{A deo patre et Christo, domino nostro. }}\pstart Erasinus: A deo patre nostro, et domino Iesu Christo domino uostro.  \pend\pstart Gracia et misericordia et pax, ut optima nobis non habentur nisi desuper descendentia a patre luminum, a quo omne bonum, ea precibus nostris et sanctorum fratrum obtinere debemus.  \pend\pstart Iam nunc in nouo testamento deus iugiter pater no minari uoluit, inueteri dominus et zelotes: fratres  \pend
\section*{COMMEN. IN PRIOREM }
\marginpar{[ p.10 ]}\pstart enim sumus Iesu Christi filii sui. Est autem pater uni uersae creaturae sed specialiter recipientium filium suum, et eum audientium. Quis cunctaiur adire ta lem patrem, misericordiarum, et totius consolatio nis, nisi degener et nothus, male et impie de eo sen tiens. Nihil dulcius nominari potest patre et flio.  \pend\pstart Et domino Iesu Christoj Christus dominus, non solum fidelium, sed uniuersae creaturae, cuius honorem agnoscunt, caelestia, terrestria, et inferna, data est ei omnis denique potestas in calo et terra, etiam angeli et dominationes ei seruiunt. Subiugata habet reges et regna, uelint nolint, tribunali suo comparebunt iudicandi ab eo. Doruinus uniuersorum ipse est, in cuius ueste seriptum: Rex regum etc.  \pend\pstart Domino nostro I Secundo adiicitur, quod fide lium et redemptorum proprie dominus sit, qui per fi dem recognoscunt dominum suum, qui eiusdem mandatis obediunt, doctrina seruant, promissis credunt, et uitae ueffigia imitantur. Dominum eum non agnoscunt, qui euangelium suum contemnunt, et uerbo suo non credunt prohibent autem uel non adiuuant pro uiribus,  \pend
\phantomsection
\addcontentsline{toc}{subsection}{\textit{Sicut rogaui te ut remaneres Ephesi cum irem in Macedoniam, ita facito. }}
\subsection*{\textit{Sicut rogaui te ut remaneres Ephesi cum irem in Macedoniam, ita facito. }}\pstart Quemadmodum rogaui te I Vide modestiam tan ii apostoli, non mandat, non urget, non compellit, sed  \pend
\section*{TIMOTHEI EPISTOLAM. }
\marginpar{[ p.11 ]}\pstart sincere, mansuete, paterne, Christiano pectore, rogat apostolico more, quem a Christo didicerunt. Smo re duci uolunt Christiani, non minis, non fulminibus mon anathematis ad ea quae conuen unt.  \pend\pstart Vt remaneres Ephesil Egregie diuelli poterat, fidelis et docilis discipulus a uerbis uitae, quae Chri stus ex ore eius loquebatur iugiter, sicut testatur Petrus et populus Christum sequens, quem uix aliquando nauigio declinari poterat, et desertis locis Sic Paulum aemulabatur Timotheus cupiens per la bores, pericula et itinera mogistrum sequi in patriam alienam Macedoniam, sed rogatur a Paulo ut  \pend\pstart ialarepae, Ephesus urbs est Asiae minoris, tunc cultrix simu lachri Dianae ab Ioue de caelo delapsi, quo curreba tur, non solum ex Asia uniuersa, sed ab omni orbe, sicut nunc ad Aquisgranum et Loretum, quaestum mul tum import abat ciuitati. Illic Paulus biennio praedicans profecit multum in Christianis illius urbis et a gentilibus multa passus, lege decimumnonum Actuum. Ei ciuitati praefecit doctorem Timotheum, quem antea ad Macedoniam miserat, et cum quo Ephesum redierat.  \pend\pstart Ita facitoJ Erasmus addit ex Graecis. Non sensu Pprio, niti eum uult, sed rogat ut faciat ca et loquetur  \pend
\section*{COMMEN. IN PRIOREM }
\marginpar{[ p.12 ]}\pstart iuxta doctrinam et admonitionem et infitutum Pault.  \pend
\phantomsection
\addcontentsline{toc}{subsection}{\textit{Denunciaui quibusdam ne aliter docerent. }}
\subsection*{\textit{Denunciaui quibusdam ne aliter docerent. }}
\phantomsection
\addcontentsline{toc}{subsection}{\textit{Erasmus: Vt denuncies quibusdam ne diuersam sequantur doctrinam. }}
\subsection*{\textit{Erasmus: Vt denuncies quibusdam ne diuersam sequantur doctrinam. }}\pstart Ece non inutilem nouam translationem et Graece studium. Vt denunties, hoc uolui, faceres Ephesi, de nuntiares scilicet quibusdam iudaizantibus et aliam ab euangelica doctrina subintroducentibus. Pri mum officium praelati et doctoris ostendere, euangelio solo credendum, et sacris litteris, eidem soli in sistendum, certo credendum, sufficienti, ad animae sa lutem, et deo placendum.  \pend\pstart Quicquid apostolicae doctrinae aduersatur, hoc ana thema iudicetur. Aduersatur autem philosophia multis, ideo obmittenda, maxime iunioribus, profrctis in fide non nocuit forsan, breuiter et neglectim uisa. Et quia hactenus eidem toti institimus, ideo hodie euangelium prae dicantibus, nouitatem improperamus, cum nihil sit antiquius uere Christiano.  \pend
\phantomsection
\addcontentsline{toc}{subsection}{\textit{Neque intenderent fabulis et genealo gnjs interminatis. Erasmus: nunquam finiendis. }}
\subsection*{\textit{Neque intenderent fabulis et genealo gnjs interminatis. Erasmus: nunquam finiendis. }}\pstart Subintrauerant Iudaei curiosorum dogmatum ob  \pend
\section*{TIMOTHEI EPISTOLAM. }
\marginpar{[ p.13 ]}\pstart scuri doctores, quorum auditor fuisse uidetur Apel les, valentinus, et Sileus, qui de propagntionibus aeonum, et quasi scotistarum formalitatum garriebant multa, et solum admirationem uenabantur, et quae stum, de quibus Tertullianus et Ireneus in primo contra haereticos. Hi docebant genealogias infinitas et miras fabulas de deo et angelis. Etiam Aposteli temporibus et unus semper addebat alterum, et magna diligentia et registra describebant maiorum nomina, uolebatque quilibet nobiliori familia natus uideri  \pend\pstart Fabulas denique Paulus appellat, quicquid extra dei uerbum et sacras litteras docetur, ut est totum Tal mud Iudaicum, et Cabula obscurissima, de quibus omnibus sequitur.  \pend
\phantomsection
\addcontentsline{toc}{subsection}{\textit{Quae quaestiones magis praestant quam aedificationem dei, quae est in fide. }}
\subsection*{\textit{Quae quaestiones magis praestant quam aedificationem dei, quae est in fide. }}\pstart Questiones I Cauendum sumopere Paulus do cet ab eis dogmatibus, quae causam praestant quaesionum et non dicit quarum aut qualium quaestionum Certe uanas, curiosas, gentilicias, philosophicas, metaphysicas, quaestiones omnium scholaflicorum sententionariorum, omnium pessime tractasset et taxasset diuinus Paulus, quibus nihil discitur, nisi de omnibus dubitare, contradictoria asseuerare, fallendo uin  \pend
\section*{COMMEN. IN PRIOREM }
\marginpar{[ p.14 ]}\pstart cere, uincendo uideri. Quis hodie numerat libros quaestionum, cum ex centum unus non habeatur, qui sacram tractauerit scripturam, et eo modo, ut prae stiteirt non tractasse, quasi Paulus iusserit quaestones mouere, sic studuerunt quaestionibus, ideo euanu erunt cogitationibus. Vtilia spiritui neglexerunt. Praestaat non soluunt, semper enim uani habent quod quaeratur, et stulti quaerere multum quam docere. Chri stus autem docuit, pietatem, simplicitatem, fidem. Aedificatio dei ) Aedificare in scripturis docere est, et dei aedificatio a deo docerisonat, cum spiritus suggerit ueram pietatem qua philosophi et quaestionarii ignorant.  \pend\pstart Que est per fidem ) Catholica documenta non ha bentur nisi per fidem, nec aedificatur ecclesia nisi sa cris litteris et fide, quae sola deus credi et fieri prae aepit. Estque profectus Christiani tantum per fidem, ut quo fide magis proficiat eo sit sanctior et opera fide iustificantur et placent deo. Sine fide nulla aedifi catio sed subuersio.  \pend
\phantomsection
\addcontentsline{toc}{subsection}{\textit{Finis autem praecepti est charitas de corde puro, et conscientia bona, et fide non ficta. }}
\subsection*{\textit{Finis autem praecepti est charitas de corde puro, et conscientia bona, et fide non ficta. }}\pstart Praeceptum dei id est, dei uoluntas est, credere in Christum saluatorem nostrum, et reco naliatorem pa tri, a quo missus, finis autem fidei est charitas in quam  \pend
\section*{TIMOTHEI EPISTOLAM. }
\marginpar{[ p.15 ]}\pstart fructificat. Opera enim charitatis sic sunt fidei fructus, ut sine eis nec sit, nec apareat. Charitas autem duplex: Dei qui fidem exigit et prae stat et fide placatur, et colitur. Proximi denique charitas exfide operatur. Si non operatur, nec charitas est, nec fidem ostendit. Noli iactare fidem et Christianismum, nisi operibus charitatis. Non satisfacies deo uero et confessione fidei. viuax est et actione augescit, et sine operibus moritur, et mortua ostenditur.  \pend\pstart Nec ficta charitatis opera, sed corde puro, fide sa licet purificato, ut quae bona fiunt, propter dei gloriam tantum fiant, sine respectu praemii uel mercedis, aliter enim est deo uti, et hypocritam agere, duploque deo displiaere.  \pend\pstart Conscientia bona, et fide non ficta, explicandi gra cia haec addunt. Primo cum idem sit, puro corde, bo£ na conscientia, et fide germana, agere dei uoluntatem, et benefacere proximo.  \pend
\phantomsection
\addcontentsline{toc}{subsection}{\textit{A quibus quidam aberrantes conuersi sunt in uaniloquium, uolentes esse legis doctores, non intelligentes, neque quae loquuntur, neque de quibus affirmant. }}
\subsection*{\textit{A quibus quidam aberrantes conuersi sunt in uaniloquium, uolentes esse legis doctores, non intelligentes, neque quae loquuntur, neque de quibus affirmant. }}\pstart Necesse est eos errare qui fabulis et questionibus insistentes dei uerbum nobis tantum proposi\pend
\section*{COMMEN. IN PRIOREM }
\marginpar{[ p.16 ]}\pstart tum, negligunt uiam scilicet ueritatis, scripturae et Christum, quem omnis scriptura tantum docet, uetus et noua, cuius scopus est unicus, et finis totius legis charitas fraterna. Otiosi enim a cura proximo rum et Christum ignorantes, et minime uiuere cupientes, quid agerent, nisi sensu humano, curiosa su perflua, uana, nihil ad pietatem facientia cogitantes, inuoluerunt se talibus quaestionibus, quaerentes seipsos, luxum, quaestum, et fastum.  \pend\pstart Doctores legis se uocant, soli qui dei uerba ignorant quod humiliat, et deprimit, sensum, nec sinit tu mescere. Solum enim Christum et spiritum Christi doctorem agnoscunt, qui uere sapiunt, humiliantur et ignorantiam fatentur. Quanto magis proficiunt tanto plura se ignorare deprehendunt, insciciamque ́ agnoscunt, nec doctiores uideri uolunt, quam sint agnoscunt curtam suppellectilem, non ambiunt titu los, non iactant, sed laudari erubescunt.  \pend\pstart Non intelligentes 7 Liquet legentibus Valentinianorum, Apellis et aliorum, illius temporis uanorum scripta doctorum, quod plura scire se uideri uo lebant, et uerbis inusitatis sua inuoluerunt, et praeter scripturam plurima somniabant, sic quod loque bantur, ut nec ipsi intelligerent, nec hi quibus pro mer aede et quaestu loquebantur.  \pend
\section*{TIMOTHEI EPISTOLAM. }
\marginpar{[ p.17 ]}\pstart Quanta autem obscuritas sit quaestionum in schola fticis uel ex hoc liquet, quod perpetuo consequitur, quod semper discipuli a magistris recedunt, et quae stionibus sese fatigant, et argumentis inuoluunt qui bus sese explicare nequeant. Regula enim ueritatis canonica scriptura perdita, non mirum si in Christiana sinceritate def cimus et gentilium ac Iudaeorum sapientia et superstitione degenerauimus. Lumen uero animarum per didimus, intellectum legis dei et fidem et amorem.  \pend
\phantomsection
\addcontentsline{toc}{subsection}{\textit{Scimus autem quia bona est lex, siquis ea legitime utatur. }}
\subsection*{\textit{Scimus autem quia bona est lex, siquis ea legitime utatur. }}\pstart Dubitare nemo potest legem, et quicqui d a deo est, bonam esse. Et cum lex doctrinam sonet, et uitae bonam institutionem, et a malo dehortationens, non potest non esse bona. Cunque ab finem bonum ordinetur, ad charitatem scilicet dei et proximi, necessario bona est. Verum nihil tam bonum quo non possumes abuti ad malum. Nulla legis obseruatio bona, nsi qua tenus seruiat charitati. Finis enim praecepti charitas, quam siquis negligat, nihil profuerit obseruatio legis. Abusus potius est si militet contra charitatem quam solam scopum habere debet, quicquid constituitur. Quin et deus suam legem uult charitati ser uire, ut patet de sabbato in euangeliis et Dauid pa  \pend
\section*{COMMEN. IN PRIOREM }
\marginpar{[ p.18 ]}\pstart nes prohibitus comedente.  \pend
\phantomsection
\addcontentsline{toc}{subsection}{\textit{Scientes hoc, quia lex iusto non est posita. }}
\subsection*{\textit{Scientes hoc, quia lex iusto non est posita. }}
\phantomsection
\addcontentsline{toc}{subsection}{\textit{Erasmus: Sciens illud, ꝙ iusto lex non sit posita. }}
\subsection*{\textit{Erasmus: Sciens illud, ꝙ iusto lex non sit posita. }}\pstart Non abutitur lege, qui scit finem legis charitate esse, quam, qui cum fide habet, sine qua haberi non po test, iustus est. Iusticia enim ex fide est, quae per chari tatem operatur. Eidem lex non imponitur ut onus, sed regula actionum, quam non ut legem obseruat, et non quasi ex necessitate, sed in spiritus libertate, ex amore et cum gaudio, sciens hoc ipso deo placere et proximo gratificari, nihil minus facturus etiam sine necessitate legis impositae. Non enim impedit in sic iusto lex libertatis, spiritum, sed dirigit et certum facit, deo placere quod agitur, a quo legem acepit, ut studiosius et alacrius operetur.  \pend
\phantomsection
\addcontentsline{toc}{subsection}{\textit{Sed iniustis et non subditis, impiis et peccatoribus, sceleratis, et contaminatis, patricidis, et matricidis, homicidis, et fornicariis, masculorum concubitoribus, plagiariis, mendacibus, periuris. Et siquid aliud sanae doctrinae aduersatur. }}
\subsection*{\textit{Sed iniustis et non subditis, impiis et peccatoribus, sceleratis, et contaminatis, patricidis, et matricidis, homicidis, et fornicariis, masculorum concubitoribus, plagiariis, mendacibus, periuris. Et siquid aliud sanae doctrinae aduersatur. }}\pstart Iniustis I Qui enim sine fide dei uiuunt, dominum non timent, iuditium non uerentur, uitam fu\pend
\section*{TIMOTHEI EPISTOLAM. }
\marginpar{[ p.19 ]}\pstart turi saeculi non credunt, legem dei non aduertunt, hos necesse est a malis studiis, legibus non solum dei, manutenendis, sed ab hominibus statuendis, retineri, oeroeri, et puniri, non tam in suum bonum, quam eorum quibus cum conuersatur.  \pend\pstart Non subditis) Erasmus: inobsequentibus.  \pend\pstart Inobedientibus, ut qui sine lege, sua sponte, sanctis monitis a malis retrahi non uolunt, legis timore, et poenarum incommodo, in suum bonum a ueti tis arceantur et malis.  \pend\pstart Impiis I qui deum non credunt uel contemnunt uel non curare nostra suspicantur. Ideoque sine fide quicquid licet agere non uerentur, horum malitiis lege obsistendum est, ut si non dei iram, saltem ho minum poenas legibus ferendas timeant, et iuxta eas uiuere compellantur. Pii autem homines non expectant legum mandata, sed suapte probitate et sinoe ritate agunt sancte, charitatisque ; dulci uinculo, uina untur, praecurruntque potius, quam ducantur, tantum abest, ut cogi uelintuel debeant.  \pend\pstart Pecatoribus Pecutores cum omnes simus quantum libet religiosi ac pii,ut nec Apost. Ioannes diffiteri possit nec aliquis sanctorum, sequitur omnibus esse le gem positam hominibus, quia pecetores, sed iusti ne hilomnussunt et exleges, qui infide et charitate  \pend
\section*{COMMEN. IN PRIOREM }
\marginpar{[ p.20 ]}\pstart uiuunt et sinelege, hoc est sine legis coactione quae legis sunt libere, sponte, etalacriter agunt. Qui uero ad legis necessitatem, uiuunt, et ut fit, legem ostendentem pecatum odiunt, hii uere peccatores et iniusti sunt, poenitentia et conuersioni indigent, donec tales fuerint non saluantur.  \pend\pstart Sceleratis I Erasm, irreuerentibus, qui non sua sponte reuerentiam maioribus exhibent. Hii ne ordo rerum turbetur legibus disciplinandi sunt, ut cui honorem, honorem etc. Honorem maioribus per con temptum negantes pacem, et concordia rumpunt. Hii lege coërcendi et humiliandi per magistratum.  \pend\pstart Contaminatis I uel prophanis, impuris. Soli ergo puri, mundi, sacri, quorum mentem, dei spiritus inha bitat ex sine lege necessitate, non tamen stne lege ui uunt, habent enim dei legem et beneplacitum, cui sponte obediunt. Contaminati qui omnium spurcitiarum genere se foedos ostentant.  \pend\pstart Patrici lis matricidis homicidis I Hii legibus opus habent, et ut formidine poenae oderint peccare, qui amore uirtutis carent, qui naturam superant malitia sua et contumeliam ingerunt creatori, bestias superant crudelitate.  \pend\pstart Scortatoribus I Hic uidemus fornicationem et scor tationem inter crimina a Paulo computari. Ne quis lia  \pend
\section*{TIMOTHEI EPISTOLAM. }
\marginpar{[ p.21 ]}\pstart tam mentiatur, hominum generationi et educationi et conuiciui puuuto mamme eumrariam.  \pend\pstart Masculorum concubitores ꝗ Crimen infandum se cundum species suas non paucas, deo naturae et ho minibus abhominabile et coelesti ultione dignum, acer bis et duris legibus uindicandum, infidelium peculiare hominum et det contemptorum, et humanam sapienti am diuiuae praeponentium. Roma.I.Vtinam legibus matrimonii utraque reprimerentur, quibus coelibatus maxime sufgragatur.  \pend\pstart Plagiariis l’Homines furto rapientibus et uendentibus, abutentibus ad crudelitatis obsequia omnis humanitas obliuiscendo pietate.  \pend\pstart Mendacibus 7 Verbis circumueniendo, decipiendo nocendo, irridendo, mentiendo, ueritatem negando, impugnando, conculcando, omnia legibus opus habent. Periuris ) In dei contumeliam, diuint nominis de decus, blasphemiam, mendatium, infidelitatem, domino impingendo mala quae fiunt, imprecando mala quae sunt in eius dedecus, omnibus his leges dantur.  \pend
\phantomsection
\addcontentsline{toc}{subsection}{\textit{Etsi quid aliud sanae doctrinae aduersatur, quae est secundum Euangeliumgloriae beati dei, quod creditum est mihi. Erasmus: Et si quid aliud est, quod }}
\subsection*{\textit{Etsi quid aliud sanae doctrinae aduersatur, quae est secundum Euangeliumgloriae beati dei, quod creditum est mihi. Erasmus: Et si quid aliud est, quod }}
\section*{COMMEN. IN PRIOREM }
\marginpar{[ p.22 ]}
\phantomsection
\addcontentsline{toc}{subsection}{\textit{sanae doctrinae aduersetur secundum euangelium gloriae beati dei, quod concreditum est mihi. }}
\subsection*{\textit{sanae doctrinae aduersetur secundum euangelium gloriae beati dei, quod concreditum est mihi. }}\pstart Seruantibus euangelicam doctrinam, a Christo et Apostolis ac Paulo praedicatam, quae sola sana est, et credentibis necessaria, utilis et certa, non est neces saria lex quaecunque humana, quandoquidem nec Mo saica quae dei est. Iusti enim sunt, qui euangelio credunt, et Christum amplectuntur, ut suum et saluatorem et fratrem. Quicquid autem huic doctrina aduersatur, hoc lege quoque prohibitum est, et legis uirtute opus habet, quod reuincat, deterreat, puniat peaatores, et transgressores et ministros, ut qui charitatis uinculo dulci deo connecti nolunt, legum laqueis innectantur, alioqui iustis non est lex posita.  \pend\pstart Euangelium ) Promissorum bonorum et maximorum ac spiritualium, certa et laeta denunciatio, quod iam nunc instent et credentibus reddantur. Et quod deus perditum mundum inChristo sibiiperfecte re conciliauerit, et omnium peautorum satisfactionem re ceperit, ut non tam de peaatorum remissione simus sol liciti, quam ut post hac pientissimi patris beneplacitis insistamus, et ut grati filii in fide et charitate to ti inhaereamus.  \pend
\section*{TIMOTHEI EPISTOLAM. }
\marginpar{[ p.23 ]}\pstart Gloriae ) Euangelium gloriae nominauit. Quid enim gloriosius deo quam bonitas sua per Christum et euangelium nobis reuelata, docet enim nos, ut non tam gloriosum sit, Iquod in principio sua omnipotentia creatus est mundus, qui quod in fine temporum pro nobis immolatus est Christus. Glorificatur euangelio deus toto orbe, deprimit et subiicit omnem gloriam mundi. Glorificat obseruatores et creden tes facit uere sapientes et sanctos, cui quicquid contrariatur uel non concordat insanum est.  \pend\pstart Beati deil Defecit Apostelus in uerbis, deum epitheto honoris nominare uolens, quicquid enim de deo dicitur, minus est a dignitate eius. Beatum di xit, quia omnis beatitudinis fons, autor, et scopus, quem scire et amare consummata felicitas, et abso tuta perfectio, cui nihil bonorum abest, sed a quo omne bonum qui ad beatitudinem capessendam nos creauit.  \pend\pstart Creditum ICreditum sibi euangelium Paulusprae dicat, et gloriatur, ut in eo deo seruiat, et fideliter incumbat, licet non sit de duodecim Apostolis, et persecutor ecclesiae fuerit. Ex gratia tamen uocatus et missus ad gloriam euangelii disseminandam, tanquam fidelis et non adulterator et quaestuosus ac ua nus uel deses etfi igdus. Creditum dicitur quasi fide  \pend
\section*{COMMEN. IN PRIOREM }
\marginpar{[ p.24 ]}\pstart litate praecipuo nuntio, tanquam thesaurus inexhau ribilis in orbem spargendus, semen centesimum fru ctum ablaturi.  \pend
\phantomsection
\addcontentsline{toc}{subsection}{\textit{Gracias ago ei, qui me confortauit in christo Iesu domino nostro, quia fidelem me iudicauit ponens in minsterio. Erasm: Et gratiam habeo qui me potentem reddidit, Christo Iesu domino nostro, quia fidelem me iudicauit ponendo in ministerium. }}
\subsection*{\textit{Gracias ago ei, qui me confortauit in christo Iesu domino nostro, quia fidelem me iudicauit ponens in minsterio. Erasm: Et gratiam habeo qui me potentem reddidit, Christo Iesu domino nostro, quia fidelem me iudicauit ponendo in ministerium. }}\pstart Licet etiam fuerim megister legis et in uaniloquio in schola Pharisaeorum exercitatus, et humanis stuaiis non indoctus, sed gracias ago deo, qui per infusam mihi fidem in dominum nostrum Iesum Christum, me liet maxime indignum sua mera gracia iudicauit dignum, qui in suum ministerium poneretur, ut praedicem dominum Iesum Christum sine omni respectu ad priuatum uel pu£ blicum bonum meum, sed gloriam dei patris et in laudem gra ciae bonitatis suae annunciandae per me uniuerso orbi.  \pend\pstart Gracias ago eil Vocatio ad spiritualia bona et ad uerbi dei praedicationem soli deo acepta referrt debet, non huiusmodi merert possumus, praeuenimur ad id det bonitate, ut non unquam in aliquo de nobis gloriemur. Sed honorificemus deum, qui uocat et uoationi adaptat, et bona quae fiunt per nos ipse agt.  \pend
\section*{TIMOTHEI EPISTOLAM. }
\marginpar{[ p.25 ]}\pstart Et eum in suis muneribus, misericordissimum agnoscamus, et per fidem maiora recipiamus.  \pend\pstart Qui me confortauit J Cum exuiribus meis nihil possem et tamen me fortissimum legis dei propugna torem super coaetaneos meos efferrem, iustcia legali inflatus, ipse me humiliatum et exinanitum uiribus naturae, et graciam suam in bonum confortauit. Spiri tus nouo feruore imbuit et inflammauit, ut non solum arserim, sed etiam inflammauerim alios.  \pend\pstart In Christo Iesu domino nostro I In gloriam et lau dem Iesu Christi, quem deus mundo lumen, creatorem, et redemptorem et doctorem instituit. Cuius uerbi et euangelii me praedicatorem instituit, cuius gracia et meritis confortatus a deo gracias ago.  \pend\pstart Quia fidelem me iudicauit ) Non inuenit fidele sed per gratiam suam sum id quod sum, decreuit me facere fidelem sibi, quem praesumptuose fidelem re perit in traditiones Iudaicas, ut tanto magis eluce ret bonitatis diuinae eminentia, quanto indigniorem deligeret in opus suum.  \pend\pstart Ponens in ministerio ) Vt sim non proprie doctrinae inuentor, sed uerbi sui tantum minister, idoneus non piger, et propriae utilitatis captator, sed dei minister in opus sanctum, non in gloriam praelationis, sed ministerium subiectionis ut eidem mi\pend
\section*{COMMEN. IN PRIOREM }
\marginpar{[ p.26 ]}\pstart nistrare nunquam cessem, sed ministerium meum indefesso conatu implem. Omnis, non solum Apostoli, sed et Christiani uocatio haberi debet ut quoddam ministerium, ut mutuum semper famulemur, ideo uoluit inter nos esse multos, imbecilleis et infirmos. Et paruulos in Christo nutriendos uerbo et exemplo. Vnum egere alterius opera, et dona ministerii singula singulis impartiri.  \pend
\phantomsection
\addcontentsline{toc}{subsection}{\textit{Qui prius eram blasphemus et persecutor, uiolentus, sed misericordiam adeptus sum, quod ignorans fecerim per incredulitatem, }}
\subsection*{\textit{Qui prius eram blasphemus et persecutor, uiolentus, sed misericordiam adeptus sum, quod ignorans fecerim per incredulitatem, }}\pstart Christum dei filium, eius caelestem doctrinam, dei patris exhibitam charitatem in Christo, stupenda miracula blasphemo ore criminatus sum, dixi ut caeteri pharisaicae scholae sectatores, demoniacum, haereti cum, seductorem, plusquam pecatorem, et amicum pecatorum, et contemptorem diuinorum pariter et humanorum, docui haec eadem alios, seduxi quos po tui, et quos non uerbis ualui, factis inftiti.  \pend\pstart Persecutor I Non suffrcit oris blasphemia et clamor sacrilegus, quin et ubique persequebar, primum quidem 25 discipulos, Actuum, et Stephanum, sed et epistolas acepi, ad omneis synagogas, ut si quos inuenissem etc. Non neg impictatis meae in\pend
\section*{TIMOTHEI EPISTOLAM. }
\marginpar{[ p.27 ]}\pstart felicltatem. Nec solum personas persequutus sum, sed usque ad mortem perstiti, ubique traxi, tradidi, accusaui, incitaui, damnaui, lapidaui, interfeci.  \pend\pstart Sed misericordiam adeptus sum) Vbisuperabundauit crimen, ibi superabundat et gracia, prae uenit me sine meritis, impediuit in malo, corripuit clamore de caelo, indicauit errorem, suppeditauit doctorem, non solum de Damascena ecclesia, sed etiam de caelo, et de tertio caelo, et ad supereffluen tiam, ut plura eorum eloqui non oporteat, addidit et animum publice confitendi Christi gloriam et ueritatem euangelicam, utque posterius defen sarem quos prius uastator persequebar agerrime semper, et ubique. Quomodo sic inimicessurus conuerti potuissem, nisi praeuentus gracia dei si prior dilexisset me, nisi me elegisset, nisi celebriorem suam gloriam facere uoluisset, quanto malum meum fuisset impotentius.  \pend\pstart Quod ignorans feoerim per incredulitatemIRe pletus etiam saentia traditionum humanarum, quae a magistris nostris didiceram, promissiones diuinas de modo mittendi Messia et Christum suum nemo docuerat, legis iusticiam sufficere mihi putabam, praeter quam nullam saebam etc. multa sacerdotii Leuitici cere mas maximi duoebam. Templum crebro ut religosior  \pend
\section*{COMMEN. IN PRIOREM }
\marginpar{[ p.28 ]}\pstart giosior uisitabam. Cum itaq uiderem a Christi doctrina et Apostolorum alienari fratres meos ab humanis et Iudaicis illis traditionibus, quas omnes parui faciebant, Christi dei filium asserebant, et deo aequalem, Sabbatum infringentem, manus non lauantes, pau cas uel nullas oblationes offerentes, orationes solen nes priuatis et clanculariis ad deum desideriis, postponentes. Sacerdotum et scribarum autoritatem la befactantes etc. multa, quasi Lutherizantes. Existimaui falsa sententia excecatus, et ignorans omnes tales cum Christo illo suo iustissime eliminandos, occidendos, et eorum memoriam exterminandam, ne dum personas, hoc ipso arbitratus me obsequium prae stare deo etc.  \pend\pstart Facile misericordiam dei consequamur, si ignoranter peccamus nisi affectemus ignoranti am, quae est peccatum in spiritum sanctum qui inspirat, cuius in spirationem ad bonum detestari, est in spiritum sanctum peccare. Sed per incredulitatem peccator potest habere excusationem, si studiosus quispiam et pius, ut erat Paulus, non intelligit aliquid de necessi tate credendum ob defectum doctorum. Paulus autem Christum non uiderat, miracula ex auditu tantum didicerat. Et utriusque rumoris autores audierat hinc pontifices et scribas et Phariseos Christianorum  \pend
\section*{TIMOTHEI EPISTOLAM. }
\marginpar{[ p.29 ]}\pstart et Christi aduersarios suspitiebat tanquam primores et ueritatis indagatores credebat simpliciter contra Christum. Hinc simplices, idiotas, humiles et doctrinam crucis despiciebat, et miraculis parum credebat, homo natura prudens et suspitiosus. Ideo incredulitas fere ei inuincibilis erat, et idco dei mi sericordiam suscepit ob ignorantiam et incredulitatem. Sed non sic de Pharisaeis aliis et pontificibus, qui Christum audierant, uiderant miracula, saebant quia a deo uenisset, quia nemo haec signa facre poterat, nisi deus fuisset cum eo.  \pend\pstart Pecabat tamen, quia non apostolos audire curabat prophetas et Christi uerba et opera non componebat scrutando scripturas. Magnatorum iuditio magis quam simplicium sinceritati eredebat.  \pend
\phantomsection
\addcontentsline{toc}{subsection}{\textit{Superabundauit autem gracia domini cum fide et dilectione, quae est per Chri}}
\subsection*{\textit{Superabundauit autem gracia domini cum fide et dilectione, quae est per Chri}}
\phantomsection
\addcontentsline{toc}{subsection}{\textit{stum lesum. Erasmus: Exuberauit autem supramodum gracia domini nostri cum fide et dilectione, quae est per Christum Iesum. }}
\subsection*{\textit{stum lesum. Erasmus: Exuberauit autem supramodum gracia domini nostri cum fide et dilectione, quae est per Christum Iesum. }}\pstart Ex mera autem dei misericordia, quae me praeor dinauit ad suum obsequium, quam, quia peaator senper eram mereri non poteram, alioqui etiam gracia dici non possit. Sed elegerat me ad portandum  \pend
\section*{COMMEN. IN PRIOREM }
\marginpar{[ p.30 ]}\pstart nomen suum coram gentibus et eum sibi seruum de legarat: Contulit mihi eam fidem ac eam dilectionem, quam nullis meis studiis obtinere, sed per graciam Christi Iesu suscipere poteram. Et ea in gradu magno, quia multum pecaueram, zelo nimio sine sa entia dei uera et cum ignorantia excusabili penes Christum.  \pend\pstart Supra modum autem gratia collata dicitur, quia sine praedicatione, et iam tunc spirante minas et cae des in discipulos. Et reuelatio tertii caeli et doctore Christo instituebatur. Et subito: Et imbuebatur spiri tu sancto et replebatur, et ab illo omnia bona suscipiebat, tunc cum iam persequeretur Christum, ut Chri sti domini benignitas apareret magis.  \pend\pstart Fide autem uera pariter implebatur per gradam quae sine dilectione esse non poterat. Operatur enim per dilectionem dumuiuit, et ideo magnafides in Faulo, magnam charitatem ostendit, et maxima cha£ ritas fecit eum plus omnibus laborare, et sollicitum esse de omnibus, sicut prius omneis uoluerat profligasse.  \pend\pstart Gracia autem Christi non stat in aliis qua in fide et dilectione, non in operibus legis Mosaicae, nisi in fide et dilectione complerentur. Nullum euim opus praesente fide et dilectione displicet, sed non iusti\pend
\section*{TIMOTHEI EPISTOLAM. }
\marginpar{[ p.31 ]}\pstart ficat, sed fides sola, et ca quae est in Chriftum qui est nostra iusticia et remissio peautorum et legis  \pend
\phantomsection
\addcontentsline{toc}{subsection}{\textit{obseruatio. Fidelis sermo et omni acceptione dignus, quia christus Iesus uenit in mudum ut peccatores saluos faceret, quorum primus ego sum. }}
\subsection*{\textit{obseruatio. Fidelis sermo et omni acceptione dignus, quia christus Iesus uenit in mudum ut peccatores saluos faceret, quorum primus ego sum. }}
\phantomsection
\addcontentsline{toc}{subsection}{\textit{Erasmus: Certus sermo et dignus quem modis omnibus amplectamur, ꝙ Christus Iesus uenit in mundum, ut peccato res saluos faceret, quorum primus sum ego. }}
\subsection*{\textit{Erasmus: Certus sermo et dignus quem modis omnibus amplectamur, ꝙ Christus Iesus uenit in mundum, ut peccato res saluos faceret, quorum primus sum ego. }}\pstart Fidelus sermo I Certissima traditio fidei nostrae et indubitabilis promissio dei, cuius fides Christianum facit et fidelem. Sermo quo nihil quidem cer tius dici potest, et credi utilius, et amplecti suauius et utilius.  \pend\pstart Et omni aceptione dignus ) Nihil optabilius mundo dici potest, euangelium plausibilibus annun ciari non potuit, cuius prae dicatio omnium antiquissi ma statim primis parentibus et omnibus patriarchis promissa et credita et sperata fuit, cuius fide saluati sunt omnes, qui fidelissime crediderunt, et desiderantissime praestolabantur implendam quod nos dei misericordia gaudemus impletum.  \pend
\section*{COMMEN. IN PRIOREM }
\marginpar{[ p.32 ]}\pstart Christus Iesus uenit in hunc mundum pecatores saluos ficere Z Qui dudum in patriarchis et prophetis loquebatur, nunc dixit, assum, de sinu pa tris sponte obediens, et hunc mundum perditum dolens, misericordia motus redemptor aduenit in hunc tam miserabilem mundum errore, ignorantia infirmitate, malitiaque repletum per omnia. Fuit qui dem semper omnia in omnibus, abfuit nulli creaturae, gubernauit et impleuit diuinitate sua mundum. Sed in hominis forma et ueritate nunquam inmun do prius aparuerat, nunc uterum uirginis egressus in hunc mundum, uerus homo et deus pro homine certauit cum daemone quem deus uicit per hominem.  \pend\pstart Peautores saluos facere I Omnes peautores re perit, ideo omnes saluauit. Non iustos, qui uel sint uel se esse putent nemo saluotur, nisi peaati sui consci us, et humiliter redemptorem desiderans, de seipso desperans, inscitiam suam ut est nihilifaciens, miseri ricordiam petens et indulgentiam promittente credens.  \pend\pstart Saluos facere IEst de peguto liberare, non quidem ne insit, donec uixerit, sed per fidem et charita tem deo unire, amicum deo facere, peaata omnia non ad damnationem imputari, conscientiae consolationem aertam tribueret de dei misericordia oertum  \pend
\section*{TIMOTHEI EPISTOLAM. }
\marginpar{[ p.33 ]}\pstart reddi, et pro totius mundi delictis deo patri satisfac re, infernum claudere credentibus, et deo semp uiuere Euangelii haec breuissima summa, alta mente reponenda, et inspiritus dulcedine semper ruminanda, simulque peaati magnitudo, quod tanti redempto ris necessitatem induxit, ut iugiter magnificemus gra ciam dei, et culpam nostram.  \pend\pstart Quorum primus ego sum I Ego sicut primus Christum dominum persecutus sum, et in suis membris ocadi uoto et facto, et ecclesiam suam agerrime afflixi, et primum martyrem ocadi, sic primo omnium tantus inimicus Christi misericordiam et graciam expertus sum. Ego Christum blasphemaui, ipse me uas electionis, gentibus prae dicatorem transmisit. Ego profligare cupiebam, ipse sectatorem fecit. Ocudere cu piebam, ipse me propugnatorem fidelium ordinauit Nomen odiebam quod nunc praedico et magnifico. Iudaeos tuebar, quos nunc confundo et impugno. Omnia secus quam uolebam, a deo miserante oesserunt, ut miraculum conuersionis meae clarius in dei gloriam eluoesoeret.  \pend
\phantomsection
\addcontentsline{toc}{subsection}{\textit{Sed ideo misericordiam consequutus uel adeptus sum, ut in me primo ostenderet Christus Iesus omnem patientiam et clementiam ad informationem eorum (ad }}
\subsection*{\textit{Sed ideo misericordiam consequutus uel adeptus sum, ut in me primo ostenderet Christus Iesus omnem patientiam et clementiam ad informationem eorum (ad }}
\section*{COMMEN. IN PRIOREM }
\marginpar{[ p.34 ]}
\phantomsection
\addcontentsline{toc}{subsection}{\textit{exprimendum exemplar his, qui credituri sunt illi in uitam aeternam. }}
\subsection*{\textit{exprimendum exemplar his, qui credituri sunt illi in uitam aeternam. }}\pstart Sanctos ideo sinit enormiter peccare, ut sanctiores, feruentiores, humiliores, fide et fidutia firmiores exurgant, et de se nihil praesumant, omnia det misericordiae tribuant, in qua sola sperent, affectuosi us pro pecato orent sibiipsis displioeant. Petrum et Paulum doctores Christianorum pecare disposuit, ut cle mentiae suae speculum et exemplar haberent omnes, non ad peaundum (ad quod natura impellit nimium) sed non desperandum, ad fidendum, ad informationem nostram qui credituri eramus. Si enun tanto perseca tort tanta benignitate indulsit crimina tanta, quid non miseri pecatores speremus? Neminem itaque iud oemus de nullo desperemus, neminem abiiciamus pro deplo rato. Non exigamus etiam in sanctis innocentiam per petuam, humilis peaator gratior deo est, quam superbus sanctulus.  \pend\pstart Consequimur misericor diam, non nostris uiribus, poeutentiis, operibus, sed gratuito dei dono nos praeueniente sua gracia  \pend\pstart Omnem patientiam I Id est, perfectam et absolutam utmaior em ex gere non possimus. Non gra ueis iniungit poenitenetas, non dure uindict, sed ab soluit gratis, prioris uitae conuersionem intungit, e\pend
\section*{TIMOTHEI EPISTOLAM. }
\marginpar{[ p.35 ]}\pstart uendationemque ut culpam agnoscamus, et pro ea humiliter ueniam postulemus, et firmiter obtinendam credamus.  \pend\pstart In uitam aeternam J Referendum puto ad dei misericordiam, quae credentes quantumlibet pecutores per clementiam informat, et adducit in uitam aeternam, quam nullis nostris meritis promeremur, sed dei misericor dia credentes assequimur.  \pend
\phantomsection
\addcontentsline{toc}{subsection}{\textit{Regi autem saeculorum immortali, inuisibili, soli deo, honor et gloria, in saecula saeculorum Amen. }}
\subsection*{\textit{Regi autem saeculorum immortali, inuisibili, soli deo, honor et gloria, in saecula saeculorum Amen. }}\pstart Crebro memoriam habet Paulus maximae gratiae dei, sibi a Christo factae, qui post tam horribilem per secutionem Christoillatam, tantam rependit misericordiam, longanimitatem, patientiam, et affiuentiam donorum spiritualium sibi, cum inimicus esset delatam, uocationemque mirabilem ad Apostolatum, et se exemplar datum pecatoribus fide in Christum saluandis, et qua memoria, iam nunc promissa permotus, gratitudine et admiratione tantae gratiae a deo susaeptae, exclamat, nec gratiam reddens, et beneficium agnoscens, dicit: Regi autem saeculorum immortali etc.  \pend\pstart Rex saeculorum omnipotens dicitur qui caelum et con tenta condidit oia, ac sapientiasua gubernatmoderatur  \pend
\section*{COMMEN. IN PRIOREM 36 }
\marginpar{[ p.36 ]}\pstart et regit, et in saecula omnia semper immutabilis perseuerat, ne solem, stellas, et caelos reges nostros credamus, sed qui haec omnia fecit in gloriam suam, quae suo nutu subduntur, et obediunt, in cuius manu sunt non solum corda, sed animae uoluntates et uita regum, qui sine eo nec momentum uiuunt nec possunt uel minimum, et hoc per saecula omnia.  \pend\pstart Saeculum ) Saeculum spatium humanae uitae, communiter centum anni et saepe quinquaginta. Item anni restantes ad quinquagesimum Iubilei, etiam si pauci sint. Hic saeculum ampiendum est longaeuitas omnium uiuentium.  \pend\pstart Creator itaque rex noster est, semper uiuens quan tumlibet s’equantur se uicissim hominum generationes et oentenarii, nos ut aqua dilabimur, ipse fluuius iugiter perseuerans. Nos temporarii, ipse aeternus, sine initio et fine, immobilis, perseueraus, dat cunsta moueri.  \pend\pstart Immortali I Immortalis qui dat omnibus uitam in quo uiuimus mouemur et sumus, non ut dii falsi, qui homines mortales fuerunt. Qui nihil perdere potest, semper immutabilis, semper perfectus, sum mum bonum, et ideo sine summo malo salioet morte uel nihilo corruptione. Neque in Christo diuinitas mor tua sed tantum caro, de patre ac diuinitate loquitur.  \pend
\section*{TIMOTHEI EPISTOLAM. }
\marginpar{[ p.37 ]}\pstart Inuisibili ) Patrem nemo uidit unquam, nemo ul det et uiuet, nedum oculis, sed nec intellectum qualis sit et quantus aliqua intelligit creatura, etiam an glica, sed filius qui est in sinu patris, ipse uidet, nouit, comprehendit, reuelat. Si uult uidetur, agnosci tur, non naturale speculum, sed uoluntarium, tamen a nemine perfecte uidetur, imo ne unquam alicui uifus est pater. Omnia agit pater per filium, per quem fe cit et saecula, deus semper uisibilis, uerbo suo omnia administrat, potestas Christo data est omnis, in caelo et in terra, dei filio.  \pend\pstart Soli deo honor et gloria ) A quo omne bonum ideo ei honor, bonitas, testimonium. Si regi honorem dependimus, non regi, sed deo cuius ipse minster est et dona regia acepit. Sic pro sancti tate, uirtute, sapientia, honor omnis deo debetur, qui solus haec omnia donat, ut non tam illa quam deus in illis honoretur. Alioqui creaturae tribuentes, quod dei solius est idololatriam committimus, ut hodie uidemus ho norari sanctos, a quibus petimus, quae solus deus donare potest, et deo propriumillis communicamus. Soli deo omnis honor, a quo omne bonum.  \pend\pstart Soli deo ) Non excluditur Filius et Spiritus sanctus, quin et illi sunt unus cum patre, immortalis, et inuisibilis, sed pater omnium perfectionum  \pend
\section*{COMMEN. IN PRIOREM }
\marginpar{[ p.38 ]}\pstart diuinarum, fons est et orig emanationum, et ipsius denique deitatis, quae lioet eadem sit in tribus, sed non eodem modo habenda, quod omnia a patre sint, et ipse ab alio nihil habeat uelsit.  \pend\pstart Gloria I Recognitio et attributio omnis boni, deo gloria est, et quod dominus uniuersorum ipse sit cuius bonitas relucet in omnibus rebus, qui ab omni creatura colendus sit, et omnia optime moderetur cui omnia seruiunt, haec omnia dei gloria sunt, clara infinitae bonitatis agnitio.  \pend\pstart In saeculasaec. J Id est sine fine et initio aeternaliter AmenJ Attestatio fidei omnium quae dicta sunt, ma xime proxime dictorum de diuinitatis excellentia.  \pend
\phantomsection
\addcontentsline{toc}{subsection}{\textit{Hoc praeceptum commendo tibi fili Timothee, iuxta prophetias quae de te praecesserunt, ut milites in eis bonam militiam, habens fidem et bonam conscientiam. }}
\subsection*{\textit{Hoc praeceptum commendo tibi fili Timothee, iuxta prophetias quae de te praecesserunt, ut milites in eis bonam militiam, habens fidem et bonam conscientiam. }}\pstart Vt intelligas fili Timothee, quid maxime te facere uelim, et quid te faoere deoeat ob prophetias quae de te ab aliquibus praeuisae sunt, quia futurus esses deo dignus, episcopus et sanctissimus huiusmo di admoni tione id offitii iniunctum tibi est a me, nunc tibi commendo pr aeceptum, non quod ego impono, sed a spiritu Christi et dei uerbo tibi iniumgitur, commendo inquam tibi et con monefacio issius praecepti, ut memineris sollicite mili\pend
\section*{TIMOTHEI EPISTOLAM. }
\marginpar{[ p.39 ]}\pstart a’e bonam miliitiam contra infideleis, pseudoaposto los, criminosos Christianos, et falsos si atres, eos ma grifice superando et uincendo uera fide Christi et bona conscientia. Intelligmus Timotheum fuis c opti mae famae et sanctae conuersationis, et eruditum in sa cris literis, de quo sancti Christani prophetauerint eun futurum idoneum episcopum Et id forsan Paulus in piritu praecognouit et dei reuelatione commonitus cum LP cuangelicum episcopum ordinauit. Erant in Ephesomul ti aduersarii euangelicae doctrinae uariae conditionis gentiles scilicet, pro Diana sua ob quaestum zelantes, iu daei Paulo et eius sociis et dogmatibus exitium medi tantes, pseud oapstoli legem Mosi eum euangelio seruan dum docentes, falsi fratres, qui infirma fidelium traducebant, et persecutiones eis suscitabant et inordinate uiuebant et scandalose. Hos non poterat Timothe us aliis debellare armis quamfide et conscientia: de bebat enim fortissime uerbis doctrinae promisis et praeceptis constanterque ; adhaerere, et se contra aduer sarios uerbo ueritatis et animi fortitudine, et fide firma fiducialiter defendere sine trepidationc. Consa entia sua sibi testante immaculatam et religos a conuersati onem coram oi ecclesia, eum audacem reddere deberet ad animose et fiducialiter ogurrendum autoritate apostotatus sui, et Spiritus san. quo perfusus erat, il ue illis aduersariis potestatibus, doctrina et moribus  \pend
\section*{COMMEN. IN PRIOREM }
\marginpar{[ p.40 ]}\pstart commissam sibi ab Apostolo ecclesiam peruertere, et corrumpere, et inquietare uolentibus, uel pesti lenti exemplo, uel falsa doctrina, uel insidiosa delatione.  \pend
\phantomsection
\addcontentsline{toc}{subsection}{\textit{Quam quidam repellentes circa fidem naufragium fecerunt. Ex quibus est Hymeneus et Alexander quos tradidi Sathanae, ut discant non maledicere. }}
\subsection*{\textit{Quam quidam repellentes circa fidem naufragium fecerunt. Ex quibus est Hymeneus et Alexander quos tradidi Sathanae, ut discant non maledicere. }}
\phantomsection
\addcontentsline{toc}{subsection}{\textit{Erasmus: Qua repulsa, nonnulli circa fidem nauftagium fecerunt, quorum de numero est Hyemenus et Alexander, quos tradidi satanae, ut discant non maledicere. }}
\subsection*{\textit{Erasmus: Qua repulsa, nonnulli circa fidem nauftagium fecerunt, quorum de numero est Hyemenus et Alexander, quos tradidi satanae, ut discant non maledicere. }}\pstart Erant qui contra conscientiam et agnitam ueritatem Apostolum traduoebant, obloquebantur, mendacia in eum fingebant, et fidem semel susceptam abiidebant Hymeneus et Alexander, et qui uerbis Apostoli pertinaciter reftitit susatauit seditiones con tra Paulum. Iudaeus erat genere, et a fide apostata, et persecutor Pauli subdolus, non tamen semper ma nifestus a quo et Paulus commonet Timotheum ut sibi aueat ne nihil fidat, in Secunda epistola et Lucas lus Actuum.  \pend\pstart Quos tradidi Sathanae) Excommunicatur, et cieat a coelu fidelium, ut sibi boni ab eo auerent, et  \pend
\section*{TIMOTHEI EPISTOLAM. }
\marginpar{[ p.41 ]}\pstart non inficerentur, ideo prodidit eius malitias et ecclesiae patefecit, et is resipisceret, et illi sibi prospicerent. Sic olim excommunicatio in usu erat, sic publicorum facinorum propalatio et exhomologesis in aedificationem et pacem ecclesiae, utinam denuo hae consuetudines resumerentur.  \pend
\endnumbering\beginnumbering\section{CAPVT II.}
\phantomsection
\addcontentsline{toc}{subsection}{\textit{\huge\textbf{A}\normalsize Dhortor igitur uos, ut ante omnia fiant deprecationes, obsecrationes, interpellationes, gratiarum actiones pro omnibus hominibus. }}
\subsection*{\textit{\huge\textbf{A}\normalsize Dhortor igitur uos, ut ante omnia fiant deprecationes, obsecrationes, interpellationes, gratiarum actiones pro omnibus hominibus. }}\pstart Igitur cum necessitate habere militare militiam periculosam cum fide et conscientia opus sit, ut malis imminentibus undique fortiter occurrere possis, et agere officium Spiritu sancto te uocante commissum, quae sine dei gracia et uirtute, nec a te, nec ab aliquo praelato haberi possunt.  \pend\pstart Adhortor ) Non impero, non urgeo, importunus ut dominus, sed ut fidelis pater cohortor filium quem scto pro affectu sincero, et filii amore non minus complexaturum hortamentum, quam feruus imperium. Non enim praeceptum patris filius expectare debet, sed nutui solo insinuasse beneplacitum filio tantum suffecerit, ut protinus sedulo perficiat, quod patrem optare crediderit. Sic non mandatis,  \pend
\section*{COMMEN. IN PRIOREM }
\marginpar{[ p.42 ]}\pstart imperiis, iussionibus in rebus dei agendis, nec Christus, nec Apostoli id egerunt, sed admonitionibus, orationibus, consiliis, adhortationibus suis profecerunt, ut nihil mandauerint, nisi quae pro charitate seruanda proponebant.  \pend\pstart Vt ante omnia fiant orationes ) Cum sanctum et deo aoeeptum nihil prorsus ualeamus nostris fieri uiribus sine auxilio dei, nec uerba, nec opera, nec exempla nec sollicitudo aliqua nostra prosit quicquam, necesse erit ante omnia nostra conamina deum implorare, eius ad iutorium petere, et asistentiam, ut ipse bonum inspiret, desiderari faciat, et perfici aliunde enim non obtinebimus.  \pend\pstart Deprecationes ) Orationes sunt ad deum, quibus aduersa quae meruimus peaatis nostris per dei misericordiam amoueri precamur, iram scilicet dei et iustam poenam, et angustias con scientiae, spiritus sciliet et corporis incommoda, tentationes periculosas, et casus peccatorum.  \pend\pstart Obsecrationes ) Quibus per omnia sacra et deo grata consequi desideramus dona dei, nobis necesaria, spiritualia imprimis et corporalia, ut quod nostris meritis non possumus obtinere, eius nobis gratiam et Christi passione et meritis interoedentibus a domino postulemus concedi in gloriam dei.  \pend\pstart Interpellationes ) Cum unus pro alterius necessitate  \pend
\section*{TIMOTHEI EPISTOLAM. }
\marginpar{[ p.43 ]}\pstart ex charitate sollicitus intercedit et pro fratre orat, ut uel aduersa amoueantur, uel optabilia consequam tur, cum pro inuicem oramus ut saluemur. Graciarum ac. ) Cum pro aoeptis beneficiis deum laudamus, et omnia ei acepta reddimus, et a nobisipsis nihil nos habuisse uel potuisse fatemur, nos uicissim suae uoluntati offerimus et beneplacito suo per omnia subiici desideramus, hoc est enim eucharistiae uerum et christianiss, sacrificium laudis, sc. et confessionis et gratia Pro om. INon solum amicis fidelibus, (rum actionis. iustis, familiaribus, nobis ipsis, benefactoribus, sed omnium maxime ꝗ magis indigent dei magna misericordia, ut sunt inimici nostri, ut caritatem acapiant christianam, infideles, Iudaei, haeretici, ut respectu dei, lumine uerae fidei et ueritatis illustrentur et credant deo, iniusti et impii ut resipiscant. Ignoti nobis ut domino miserante quilibet suis pecatis liberetur. Nihil unquam interorandum discriminis occurrat, quo minus affectuose pro uno genere hominum oremus quam pro alio uel excludere aliquos praesumamus uelfrigdius intendere. Non etiam ob beneficia tantum oremus, nostrum commodum, sed ex sola fraterna charitate ad orandum pmoueamur. Et quam id domino placere credimus, similia ab aliis deside rantes. Non dicit pro animabus, sed hominibus: animae enim nobis proximae non sunt, sed deo: uiui nobis con mendantur, mortui deo. Nos contra ob quaestum cleri.  \pend
\section*{COMMEN. IN PRIOREM }
\marginpar{[ p.44 ]}
\phantomsection
\addcontentsline{toc}{subsection}{\textit{Pro regibus et omnibus in eminentia constitutis, ut placidam, ac quietam uitam degamus, in omni pietate et honestate. }}
\subsection*{\textit{Pro regibus et omnibus in eminentia constitutis, ut placidam, ac quietam uitam degamus, in omni pietate et honestate. }}\pstart Reges Apostoli tempore infideles erant et fidelibus aduersarii sed tamen pro eis orari uoluit: imo tanto magis, quanto respectus dei magis eis pro regimine nostro necessarius ut boni fiant, ut tractabiles, ut mansuete ac benigne regant, ut sustinere possimus, ut cor eorum e domno regatur, ut tranquillius regamur ab eis, ut sapientia praestetur illis et clementia, et ante omnia timor dei et amor iusticiae.  \pend\pstart Et omnibus qui in sublimitate consuluti sunt ) Et coniunctio crebro expositiue accipienda. Regibus, id est, sublimibus: reges enim omnes sunt, qui in magistratu aliquo consuntuuntur ut praesint et regant, moderentur. Constituti autem non ab alio, quam a deo. Omnis enim potestas a deo. Etiam mala propter peccata populi, et dei imminens iuditium, etiam mali confutuuntur, ut per cos sua deus iudicia exequitur. Dabo tibi regem in furore meo. Effoeminati erunt principes eorum, et c. Maximam poenam de precabimur, si bonum principem orationibus promerebimur propter pecata populi.  \pend\pstart Sublimitas est regere subiectos iuxta stilum Spiritus sancti et a deo ordinata, non ut honorentur, sed ut  \pend
\section*{TIMOTHEI EPISTOLAM. }
\marginpar{[ p.45 ]}\pstart onerentur sollicitudine maiori prodesse cupientes eis quibus ut seruiant a deo dati sunt. Reges gentium dominantur et benefici dicuntur. Non sic reges fide lium, hi sciunt sese tanto modestores esse debere, quanto sunt multorum commodis destinati et obsequio charitatis deuincti. Debetur eis honor, reuerentia, timor et obedientia, et sustentari a subiectis. Ipsi uicissim debent amorem, curam, tutelam, instructionem, consolationem, regmen, coërtionem a malis ac poenam discolis.  \pend\pstart Constituti ) Non negligenter ol mittendum. Non sibi ipsis, prece precio, et inconstantia, sublimitas, et regimen ambienda, sed constitui debent ab aliis, nihil nisi sub ditorum bonum respicientibus, et non sine consen su eorum, ut non tam seruitus, quam administratio regmninis habeatur.  \pend\pstart Vt placidam et quietam uitam agamus) Anniti omni studio debent reges, ut populus beneuolenter et cum gaudio, non ex trifticia obediant, sed placite optemperet ex sincero affectu, et sine murmure et dispicientia, si posset fieri sine uiolentia, ut sint omnia sine tumultu, uiolentia, odio, seditione, rebellione, sed Christana pace conuersentur, ut fratres, et domestici et discipuli Christi pacifici, modesti, mensueti et humiles. Haec est enim uita filicissima, planeque bea\pend
\section*{COMMEN. IN PRIOREM }
\marginpar{[ p.46 ]}\pstart ta et Chriftiana.  \pend\pstart Cum omni pietate et castitate ) Melius honestate, faciunt haec duo perfectam quietem si deum nostrum colamus fide, spe, et charitate, ut totos nos ei committamus, ut patri et omnia quae ocurrunt patris prouidentia et pietate euenire certo credamus, tunc enim murmm aomms minquam, temere non iudicabimus, sed ut sunt omnia, optima cordis paoe ut patris beneficia amplectamur.  \pend\pstart Honestas denique in conuersatione ad Christum et morum omnium castitas et temperantia, nos commendabit fratribus, cauebit offensam et inimicitias et perpetuam tranquillitatem fouebit et amicitiam quae inter inhonestos et leues constare non potest diu.  \pend
\phantomsection
\addcontentsline{toc}{subsection}{\textit{Hoc enim bonum est et acceptum coram saluatore nostro deo, qui omnes homines uult saluos fieri, et ad agnitionem ueritatis uenire. }}
\subsection*{\textit{Hoc enim bonum est et acceptum coram saluatore nostro deo, qui omnes homines uult saluos fieri, et ad agnitionem ueritatis uenire. }}\pstart Quieta conuersatio nostra et tranquillitas simplicis obedientiae regibus etiam ethnicis exhibenda, bonum, est nobis, ut careamus dominantium molestia, ut non grassentur, non magis opprimant, non irritentur contra nos, et contra euangelicam doctrinam, non schandalizentur in moribus, ne ansam rapiant mouendi persecutiones, blasphemandi  \pend
\section*{TIMOTHEI EPISTOLAM. }
\marginpar{[ p.47 ]}\pstart nomen Christi, tyrannidisque ́.  \pend\pstart Cedet in gloriam magnam religionis nostrae, imo et domini dei nostri, qui per filium salualiit nos, si pacifice, mansuete et sedulo dominis obsequemur, subiecti omnibus in timore dei, monstrantes optimi optimas leges scriptas in cordibus nostris. Et ex bonis operibus nos considerantes, glorifiemat deum dum aliquando respectu misericordiae a deo fuerint uisitati, et ad cor et pium religiosu perducti, dominuque ; deum nostrum ex fidelum religione cognouerint, nobiscumque uenerari ceperint, ora ionibus nostris domino commendari  \pend\pstart Omne is hom. ) Est enim dominus deus summae bonus et nihil odit eorum quae condidit, neminemque uult perire, adest gratia sua omnibus, ut ueniant ad agnitionem uerit atis, quam per propheticas et euangelicas ac apostolicas fecit communem omnibus. videtur hic denuo coniunctio et, expositiue sumi. Saluos uult fieri et ad agnitionem ueritatis uenire, ut nihil sit aliud agnoscere ueritatem quam saluum fieri. Etiam uult omnes saluos ficri quia uocauit omnes, anunciatur omnibus Christus, pro omnibus mostuus est, omnes ad se trahit, pro oibus tradidit semetipsum. contemptores autem dei et negatores non computentur inter omnia, sed redigentur in nihilum, si sermoincidat quare non omnes beatitudinem assequantur qui sunt redempti, del allegabitur uo\pend
\section*{COMMEN. IN PRIOREM }
\marginpar{[ p.48 ]}\pstart luntas, quae regula est iusticiae, et cuius non est quaerenda ratio, homines sumus, ipse deus nihil nobis debet, uoluntate signi et antecedenti uult omnes saluari dicunt scholastici.  \pend\pstart Agnitio ueritatis ) Id est, Christi qui ueritas est, qui euangelium praedicauit, qui fidem salutis portum statuit et iusticiae, lux mundi qui illuminat omnem homanem mamanae seienttae muilles et mendatium. Christum scire sanctorum est, stultus fiat in hoc mundo sit ut sapiens et saluus.  \pend
\phantomsection
\addcontentsline{toc}{subsection}{\textit{Vnus enim deus, unus etiam conciliator uel mediator dei et hominum homo Christus Iesus, qui dedit semetipsum redemptionem pro omnibus. }}
\subsection*{\textit{Vnus enim deus, unus etiam conciliator uel mediator dei et hominum homo Christus Iesus, qui dedit semetipsum redemptionem pro omnibus. }}\pstart Hic Apostolus summam Christianae fidei et ueritatis proponit, cuius solius agurit jiluos facit, omneis homines, qui scilicet agnoscunt, non fimplici notitia, sed potius sapientia, qui sic nouerint, ut sapiant etiam quam suauis sit dominus, hunc ergo tenorem fidei et regulam discutiamus.  \pend\pstart Vnus esse deus agnoscitur, cum praeter eum nihil aliud summum bonum desideratur, colitur, amatur, antefertur, confertur, nihil. Non uenter, non caro, et sanguis, aurum, honos, laus, et totus hic mundus Omnia haec postponuntur, non uerbo et lingua, sed  \pend
\section*{TIMOTHEI EPISTOLAM. }
\marginpar{[ p.49 ]}\pstart corde et sensu, dum unum illud semper intendimus et attendimus, et propter eum omnia agimus Non circa plurima ocupamur et distrahimur, hoc est uere agnoscere unum deum, praeter eum nihil desiderare quod non propter eum.  \pend\pstart Vnus concliator) Secundus articulus fidei, credere et sare Christum esse unicum mediatorem, quem statuit pater noster inter se et nos, ut cum essemus pecutores et a deo proiecti et inimia per Christum reconciliamur. Hic includitur esse omnes homines peautores dei aduersarios et inimicos, deum etiam offensum omnibus, et abiecsse damnasseque ; in primo Adam et sola Christi gratia conciliatos. Et non id solum, sed et filiationis et haereditationis aoeepisse et in Christo omnia bona, ut nostra sint, quae fratris nostri sunt Vnus autem non multi mediatores uel conciliatores, unus enim Christus, unus saluator, non ergo praeter eum interponendi sunt alii, quantumlibet sancti omnes enim pecatores sunt, et eguerunt, imo egent gracia dei. Solus Christus suapte dignitate et suis meritis, non ex gracia concliat nos patri. Omnia enim pater dedit in manus filio. Et eum currendum quoties urgemur, fons est pictatis, non estut eum promereamur, sanctorum interoessionibus, benignior nobis est, quam omnes sancti. Solus pro nobis  \pend
\section*{COMMEN. IN PRIOREM }
\marginpar{[ p.50 ]}\pstart mortuus est. Et nos toties inuitat, pulsat, expectat hortatur allicit, rogat et cogit aliquando. Venite ait ad me omneis qui laboratis et onerati estis, etc. Venite edite absque argento et caet. Audit nos senper, adest, intendit singulis, quasi non sui essent omnes, non distrahitur multitudine saluandorum, omnia nouit melius nobis, scit quibus indigeamus, praecepit ut in solo suo nomine precaremur, honorem hunc suum alteri non defert, nec patitur communem.  \pend\pstart Homo Christus Iesus ) Mirum quomodo dubitauerint hominem Christum fuisse quida haereticorum qui hic simpliciter homo dicitur: caro enim et frater noster temptatus per omnia absque peccato Et pro nobis hominem nouissimum se exhibuit et uermem, ut hominem perditum redimeret. Et primogenitus ex multis fratribus chirographum peccati deleret per crucem mortuus uere ut homo esset, qui pro homine satisfecit Christus Iesus unctione sua consecrans gentes et passione saluans Iudaeos, faciens utraque unum redemptorem.  \pend\pstart Qui dedit semetipsum ) Obtulit in remissionem omnium pecatorum seipsum sacrificium perfectum et maximum, et omnium praecedentium sacrificiorum scopum et finem pro captiuis nobis commutationem sufficientissimam. Non raptus est inuitus, sed sua sponte obediuit Semetipsum obtulit precium deo patri pro salute hu\pend
\section*{TIMOTHEI EPISTOLAM. }
\marginpar{[ p.51 ]}\pstart mani generis et reconciliatione perfecta cum inimici essemus, et ita inimici, ut etiam eum nos ipsi crucifixerimus, homines scilicet pro quibus se precium obtulit. sicut et nos in Adam peaumus. Nec dicamus non esse nos posteros Iudaeorum qui hodie quoq simile facimus non refugeremus simili occasione permoti, qua Iudaeorum primores, quibus meliores paucos habemus quin et iterum aoerbius eum crucifigimus uitiis nostris.  \pend\pstart Pro omnibus ) Dedit se pro omnibus. Nemo nisi nolens saluari relinquitur. Suffecit, satisfecit deo patri pro omnibus, nihil obmisit de suo, quo minus omnes redimi possint. Excusatio nulla restat nobis qui scimus, quid uel facere, uel desiderare, uel obmissa dolere, uel non indolentia deprecari, et humiliter agnoscere debeamus, et pro fide habemus et charitate deum orare, quis haec se non posse uero fatebitur? Ad huius ueritatis agnitionem siquis illuminante deo peruenerit ut credat et agnoscat certa fide unum suum deum, et conciliatorem unicum dei et hominum, hominem Christum Iesum, et qui ipse seipsum dederit ex maxima gracia redemptionis precium pro omnibus credentibus sine dubio saluabitur et perfectus fidelis et Christia nus coram deo et Christo suo reputabitur, quae ei fides reputabitur ad iusticiam ut Abraae, et est totius noui testam, et euangelicae doctrinae conclusio et summa.  \pend
\section*{COMMEN. IN PRIOREM }
\marginpar{[ p.52 ]}
\marginpar{[ p.53 ]}
\phantomsection
\addcontentsline{toc}{subsection}{\textit{Cuius testimonium temporibus suis confirmatum est, in quod positus sum ego praedicator et apostolus. Veritatem dico non mentior. Doctor gentium in fide et ueritate. }}
\subsection*{\textit{Cuius testimonium temporibus suis confirmatum est, in quod positus sum ego praedicator et apostolus. Veritatem dico non mentior. Doctor gentium in fide et ueritate. }}\pstart Erasmus Graece legit: Testimonium temporibus suis in quod positus sum ego praedicator etc. quod intelligo sic: Adhoc Christus ueniens se tradidit pro nobis patri, ut id suo tempore testaretur toti orbi, omnibus hominibus redemptis, ut id omnes credant, et credendo saluentur. Noluit enim otcultum quod tot sacrificiorum figuris in ueteri testamento toties uoluit esse testatum, et in figuris representatum. Quin et ego Paulus ad hoc nostrae redemptionis euangelium praedicandum, ab eodem Christo domino positus sum segregatus et electus post alios, ut toto orbe gentibus haec praedicem annunciarem hanc dei stupendam graclam. Nuncius nupus a Christo et apostolus, licet cum Christo non uixerim, tamen ex misericordia sua, ex persecutore blasphemo me conuertit in Apostolum et ueritatis doctorem ad omneis genteis.  \pend\pstart Veritatem dico non mentior ) Celebris est fama insantae meae cumspirans minarum et caedis in disapulos, subito in multorum comitiua e caelo uoca\pend
\section*{TIMOTHEI EPISTOLAM. }
\marginpar{[ p.54 ]}\pstart tus sum, et a Christo conuersus, postea denique cum Barnaba a spiritu sancto segregatus in genteis, attestatione non modica miraculorum, et pros ctu praedicationis maximo a Iudaea per Asiam et usque ad Illyricum repleui omnia euangelio Christi. Non potest esse dubia apostolatus mei uocatio, et quod Christus me gentium doctorem dederit, testantur signa apostolatus mei, qualem et me gessi fideliter in gloriam Chrifti, non meum compendium aliquod temporale, nihilque admiscendo doctrinae quam in me Christus non fuerit loquutus. Et ueritate, sine omni hypocrisi et figmento, dextre et plana ueraciter, ut acepi a domino omnes tradidi.  \pend\pstart Praedicator fuit Paulus euangelii, nullius alterius doctrinae. Iurat modestissime, quod notorium esset, et non poterat negari. Doctorem se nominat, non in uniuersitate promotus: nonenin utuiceran Aristotelem et parua logicalia, sed spiritu sancto permotus docebat sine intermissione, uerbo et exemplo, non tam nomine doctor, quam officio et assiduo studio docendi, sicut Origines, Augustinus, et Gregorius.  \pend
\phantomsection
\addcontentsline{toc}{subsection}{\textit{Volo autem uiros orare in omni loco, leuantes puras manus, sine ira et disceptatione. }}
\subsection*{\textit{Volo autem uiros orare in omni loco, leuantes puras manus, sine ira et disceptatione. }}\pstart Ante omnia orandum docuerat superius pro o\pend
\section*{COMMEN. IN PRIOREM }
\marginpar{[ p.55 ]}\pstart mnibus, quatenus tranquille uiuere possent, et ut ad agnitionem ueritatis omnes peruenirent saluandi. Nunc uiros et mulieres instruit, quomo do orandum de aliis necessariis incidenter, foeminarum mores instituendo, uiros primum.  \pend\pstart Volo uiros orare in omni loco ) Nondum amplissima templa erant sufficiebat locus pro familia, longo tempore in cryptis praedicatum, ubicumque ecclesia erat, duo aut tres ubi locus orationis erat et praedicationis. In cubilibus, in agro, in naui, in montibus, syluis, ubique orationis locus, quia ubique dominus, ubique  templum dei quod sumus nos, ubique cor fidele et deuotum Id Christus docuit: Cum inquit oraueris intra cubiculum et ora patrem in abscondito. Hodie asueuerunt multi, ut non nisi in ecclesia orare ualeant, et exaudiri credant, si tamen orant uere, et non potius multa sensuum et mentis distractione numerent potius uerba, quam corde, affectuque et feruore leuentur in deum, et in frugilitatum suarum defectus et miserias, quarum eonsideratione oratio promoueri debet. In omni inquit loco. Consecratus est homo et mundus totus orationi non templo Hierosolymis, sed in spiritu ct ueritate orandum, id quod docuit dominus, et quae se exauditurum promisit. Caetera sunt dei prouidentiae et benignitati committenda.  \pend
\section*{TIMOTHEI EPISTOLAM. }
\marginpar{[ p.56 ]}\pstart Leuantes puras manus ) Non lotas Pharisaeorum more sed puras a sanguine proximorum, a foedis, libidinosisque contractibus, et a munere et usura, ab omni malo opere. Non uult Paulus more Pharisaico manus extendi, sed more desiderantium leuare ex affectu petitiones. Non obsunt omnes cerimonae, malta projant, jeu non mate, et quae uel mentis distractionem, uel hypocrisim mouere possint.  \pend\pstart Sine ira et disceptatione) Docet nihil a deo postulandum obtine ri, nisi prius proximo conciliatus Christiana charitate fueris. Solet animum uirorum ira ocupare sicut superbia mulieres, et incontinentia. Ira autem charitati maxime contraria et inimica, oratioui locum praecludit omnino. Et nist proata reconciuatuos paupiui se oraturum nemo ignoret. Non uult ignoscere, nec pietatis dona largri, nisi peaantibus in nos, ante indulserimus, nisi fratri offe nsam remiscrimus et ut nos dilexerimus.  \pend
\phantomsection
\addcontentsline{toc}{subsection}{\textit{Similiter et mulieres in habitu ornato (amictu modesto) cum uerecundia et sobrietate (castitate) ornantes se non in tortis crinibus, aut auro, aut margaritis, uel ueste preciosa, sed quod }}
\subsection*{\textit{Similiter et mulieres in habitu ornato (amictu modesto) cum uerecundia et sobrietate (castitate) ornantes se non in tortis crinibus, aut auro, aut margaritis, uel ueste preciosa, sed quod }}
\section*{COMMEN. IN PRIOREM }
\marginpar{[ p.57 ]}
\phantomsection
\addcontentsline{toc}{subsection}{\textit{decet mulieres, promittentes (profitentes) pietatem per bona opera. }}
\subsection*{\textit{decet mulieres, promittentes (profitentes) pietatem per bona opera. }}\pstart Mulieres post uiros docet Apostolus ne superbiant de forma, amictu, moribus et ornatu, id quod agnatum solet mulieribus uitium, etiam alioqui optime licet sint castae et prudentes, sed anumni superbia et speciei corporis et ornatui praeampere uix possunt. Admittit ut habitu ornato sint, non foedo amictu et immundo, quo abhominatio maritis sint et adulterii ocasio. Sed modesto amictu, ut decet castas, uere cundas et minime irritatrices, mundae, cleganter, et honeste uestiantur, hoc est, ne nuditate et superfluo ornatu uestium, et exquisita subtilitate et curiositate meretricibus assimilentur, aparantque ; affectu placendi, in plateis uestiri id enim Christianas non decet, nec maritis placere potest, nec debet. Ornentur autem legittimis monilibus, uerecundia scilicet, et sobrietate, ex uultu, oculis, incessu, uerbis, ostendendo inuitas se offerri uirorum aspectibus, et ab aliquo desiderari in malum. Nullae semper contegi, et egre uidere uiros, multo molestius uideri. Etiam ad mari tum uerecundiam non postponant semper et ubique  Gratiosissimum enim monile uerecundiam habet, quisque sapiens maritus extiterit.  \pend\pstart Sobrietate ) Ne uinolentia et crapula notetur  \pend
\section*{TIMOTHEI EPISTOLAM. }
\marginpar{[ p.58 ]}\pstart ali quatenus quae lubricitatis et libidinis testes certi£ sunt, uenter et genitalial proxime cito spumant et caetera impedimenta praebent foetibus optatissimis generandis. Insipientiae et garrulitatis uitia semenafferunt, mendatii et litium matres, et nihil secreti seruare possint, seminaria malorum.  \pend\pstart Castitate ) Castitas sobrietate nascitur, et sunt soronres, ut Icliltet, jini foenane comtime e castae, temperantes libidines et desideria sua inhonesta ne effluant.  \pend\pstart Ornantes se non in torttis crimbus ) Non compositlone tapiuorant, O morpueaua, auri et argenti, ac adulterino colore et cerussa, fucoque peregrino. Iniuriam non faciant formae naturali, sed animi pulchritudine uirtuteque ornentur.  \pend\pstart Nec auro aut margaritis I Non onerent maritos impensis, colligant potius pro haeredibus, uel subueniant indigentibus in opera misericordia. Super a naturali non addant fomenta, non certare mutuum mulieres, nisi modestia et maturitate et sapientia. Stultum est margaritis uti, aurum expendendum potius in pauperes.  \pend\pstart Vel ueste preciosa ) Veste non uestibus, quod superfluae Christianos dedecant, et excusari non posfint. Non ad luxum sint, sed ad necessitatem. Vnica  \pend
\section*{COMMEN. IN PRIOREM }
\marginpar{[ p.59 ]}\pstart quidem non sufficit, sed superfluitas non licet. Pretiositas intollerabilis, scandalizat, irritat, zelum et scquelam susatat, opprimit mareum, difficile praecauetur tineis, et non sine sollicitudine et negotio.  \pend\pstart Sed quod decet mulieres, profitentes pictatem )  Propensiore natura ipsas fecit ad dei cultum, ad cercmonias multas et superfluas. Hanc discunt non esse pietatem, frequentare ecclesias, gestare oratoria, uel mussitare preculas, cantus mirari et organa, aspicere picturas et inclinari, cereos incendere, aquas spargere: haec qui dem pietatis signa esse deberent, sed minime sunt pietas, sed prorsus in abusum uersa sunt omnia. Pietatem profiteantur recordatione peautorum, suspiriis pro gracia et indulgentia et emendatione et sancto proposito, audire attentissime uerba dei, meditari, et ruminare audita, ut liberis postea domi recenseant. Orationes breues et feruidas. Non multas, non toties repetitas, et ex more pro numero praefixo, non quantitatem precum, sed qualitatem deus exigit. Monetam probatam charitate, non loquacitate. Domi oret, familiae prospiciat, marito ob sequatur et auxilietur. Caueat a girouago aspectu, lumen seruet operi et tenebris. Fides sine aquarum aspersione suffecerit, et pro istis omnibus exerceat.  \pend
\section*{TIMOTHEI EPISTOLAM. }
\marginpar{[ p.60 ]}\pstart Bonaopera ) Ministret liberis, marito, familiae uicinis, pauperibus, pregnantibus, puerperibus, infirmis, domui, lauando, ornando, purgando. Et omnia haec in fide agendo, certae sint, quia deo plaoeant. Et sint ipsissimae charitatas officia et opera a deo impofita, et oueuiemiae janctae optimi fructus. Si uacet lectioni incumbant, et pueris itidem doceant, moneant, corrigant et caetera, haec uera pietas.  \pend
\phantomsection
\addcontentsline{toc}{subsection}{\textit{Mulier in silentio discat, cum omni subiectione. Docere autem mulieri non permitto, neque dominari in uirum, sed esse in silentio. }}
\subsection*{\textit{Mulier in silentio discat, cum omni subiectione. Docere autem mulieri non permitto, neque dominari in uirum, sed esse in silentio. }}\pstart Viris quidem permittitur ut in coetu quae de auditis non intelligit, a presbyteris inquirat, et si quid ei spiritus reuelet loquatur cum modestia, et sine disoeptatione. Non prohibetur a uerbo ueritatis proponendo uir scripturas intelligens, qualescumque . Voluitut passim quini et seni pariter prophetarent, etiam in ecclesia cantum et organa et psallentium mores et assiduum studium non docuit, in cordibus psallere et spiritualibus hymnis domino canere mente non spiritu Christianos omneis docet. Mulieres uero in ecclesia sileant, domi uiros interrogent, siquid discendum fuerit, et in silentio, ne sensu abundent, et litigent uerbis et sensu proprio.  \pend
\section*{COMMEN. IN PRIOREM }
\marginpar{[ p.61 ]}\pstart Dooere autem mulieri non permitto ) Non audeo adiicere, in coetu uirorum, sic ut liceat eis in ecclesia foeminarum: sensu enim infirmae sunt et instabiles. Viros presbyteros, senes habeant, a quibus foeminarum congregatio instruatur publice. Nemo uero negauerit aliquot pariter conferre posse de auditis et lectis, ac legere scripta sancta in publico, sic, ut intelligant. Verum sine intellectu, orare, legere, canere, psallere, opus est sine omni utilitate et maxima temporis iactura, solaignorantia et pigritia monachorum introduxit haec omnia. Olim doctae erant moniales et seruire hoc poterant, deuotioni. Hodie non tam obedientia excusat, quam dementia culpat. Quid enim inutilius potuerit institui, quid stultius ex cogitari, quam ut tantum temporis expendant in uerbis nihil significantibus, et reliquum tempus ab operi illius fatigatione, otio, et remissione, deperdi, ut omne carum tempus longe quam somnium sit eis in fructuosius etc. Vt autem docere non permittuntur ita doceri iubentur, ut Christianae sint, et dei uerbum ignorent minime. Hic maior querela locum habet.  \pend\pstart Nec dominari in uirum ) Foedum esset marito superbiae fomentum, contemptus uiri. Monere potest, rogare blandiri, sed esse in silentio, parce loqui et breuiter, et cum reuerentia et timore, illico admo\pend
\section*{TIMOTHEI EPISTOLAM. 61 }\pstart nita tacere, taceat. Irato marito non respondeat, tempus captet oportunius. Semel et non crebro repetat. Lingua mulierum diligentius moderanda, ꝗ natura propensiores sint ad cauillationes. Non etiam Circumcursari in aedes alienas, non morari in plateis Non foris plurimos miscere sermones, domi agant in silentio.  \pend
\phantomsection
\addcontentsline{toc}{subsection}{\textit{Adam enim primus formatus est, deinde Eua. Adam non rust deceptus, sed mulier seducta, obnoxia facta est praeuaricationi. }}
\subsection*{\textit{Adam enim primus formatus est, deinde Eua. Adam non rust deceptus, sed mulier seducta, obnoxia facta est praeuaricationi. }}\pstart Ne quo pacto superbiendi ansa mulier sibi uendict donanandi in uirum, admonet Apostolus meminerit originis suae, quam a uiro sumpta aoeepit, qui primo conditus erat a deo. Natura denique perfectior et absolutior et sapientia clarior, et ad seductio nem cautior ac fortior, quo sopore merso, e costa formata est Eua, ut unum corpus et os ue uno corpore amoris uinculo arctius constringerentur, primus tempere et uirtute, Eua posterior utroque . Ideo duci, instrui, parere, sequi debet et se agnoscere dei ordinatione subiectam.  \pend\pstart Adam non fuit deceptus ) A serpente non audente eum aggredi, sine dubio uincendo. Non fefellisset mendatio tam impudentissimo. Non ambiuisset  \pend
\section*{COMMEN. IN PRIOREM }
\marginpar{[ p.62 ]}\pstart excellentiam, non colloquium tulisset, nisi in gratiam mulieris, consortis, amicae, cui cum haec placere uideret, Adam, delitias suas contristare non poterat, sed acquieuit roganti mulieri attendens, praecepti oblitus et parum curius.  \pend\pstart Mulier seducta, obnoxia facta est praeuaricationi )  Seducta mendaciis, inuidia, dolo, blasphemia diaboli et permissione, dispensationeque diuina. Nondum iusta, sed iusticiae originem nacta, qua proficere debuisset et progredi, gracia dei praesente, et praecepto acepto retrocessit, transgressa mandatum et uerbis dei non credens, sed primo infidelitatis peccato superata, superbia, auaritia, gu’a et luxuria decepta fuit, et uiro insuper insidiata blanditiis et instantia traxit in praeuaricationem, ut pariter cum ea coeperit dubitare si uero haec deus praedixerit, mortem scilicet imminere, qua post gustum uideret superuiuentem, ut ipse peauret, cui semel seductrix facta, post hac directricis et doctoris munus par nun quam fuerit, uel aequum praestare, uel praesumere. Semel tam pernitiosae suasioni consentiens Adam, merito suspectam habuit omnem mulieris informationem. Et ipsa experimento sumptae fragilitatis et insipientiae cum perfecta esset, nunc bis misera, et admonita sui praesumere doctrinam non audeat solidae ueritatis  \pend
\section*{TIMOTHEI EPISTOLAM. }
\marginpar{[ p.63 ]}
\phantomsection
\addcontentsline{toc}{subsection}{\textit{Saluabitur autem per filiorum generationem, si permanserint in fide et dilectione, et sanctificatione cum sobrietate. }}
\subsection*{\textit{Saluabitur autem per filiorum generationem, si permanserint in fide et dilectione, et sanctificatione cum sobrietate. }}\pstart Sub dita quidem sit uiro mulier, et humiliter segerat inter uiros et in ecclesia, agnoscatque mentis im becillitatem et inconstantiam, ut sibi docendi munus in ecclesia non arroget:uerum non a uiris contemnenda est, quam per generationem liberorum beatam fecit dei munificentia, et insigni munere illustrauit. Quid enim admirabilius in rerum natura generatur, quam inutero foeminae, non tantum corpus humanum, sed et anima creando infunditur, cuius generationis, nutritionis et educationis pariter et instructionis offitium contulit deus in mulierem, et in cunabula pietatis tradendae et pfectus commisit matribus ut et lacte pascant. So lido cibo a uiris nutriantur, et in fide et sacris dogmatibus perfecto: audiant et sequantur. Curet igitur liberorum mater, ut in ea quam a teneris unguiculis a se fidem et dilectionem dei imbiberunt filii, in eis permaneant, iugiter inculcando fidem et timorem dei, ut sancti sint moribus ab omni contagione criminum sibi perpetuo cauentes, carni et sanguini non indulgentes et dei mandata quibus eos mater pueros edocuit semper memoria cogitantes: sancti sint et permaneant corpore et spiritu. Sobrietatem in primis curantes, ut cibo et potu  \pend
\section*{COMMEN. IN PRIOREM }
\marginpar{[ p.64 ]}\pstart somno, otio, lusibus, lenociniis, indulgentia, sodalita te, caueant incumbere, sed potius inculcare studeant matres filiis, quos deo genuerint et proximis, fidem dilectionem, sanctimoniam, pudicitiam, sobrietatem.  \pend
\endnumbering\beginnumbering\section{CAPVT III.}
\phantomsection
\addcontentsline{toc}{subsection}{\textit{\huge\textbf{F}\normalsize Idelis sermo. }}
\subsection*{\textit{\huge\textbf{F}\normalsize Idelis sermo. }}\pstart Quae tibi filio meo dilecto ducenda ad populum scribo, ut Christi aposte lus et tibi pater, scribo quae certissima habeas, et tam quae nunc dixi de uiris et mulieribus, quam nunc, scribenda, scias a domino reuelata, et fidelia, et Christi doctrinas esse qui in me loquitur. Non sunt hominum uerba, mendatia, nihil ex me loqui in ecclesiam audeo, nec opiniones et fabulas. Fidelis est omnis sermo dei et essicax.  \pend
\phantomsection
\addcontentsline{toc}{subsection}{\textit{Siquis episcopatum desiderat, bonum opus desiderat. }}
\subsection*{\textit{Siquis episcopatum desiderat, bonum opus desiderat. }}\pstart Episcopus Graece obseruator, speculator, super intendens, et custos dicitur. Offitium episcopale est, bona procurare, et a malis custodire, curam proximi salutis gestare. Nihil autem melius quam pascere, et id uerbo dei, pascuis salutis. Qui ergo desiderat huiusmodi exercitia conferre in proximum, opus bonum desiderat ex charitate propter deum. ut benefaciat  \pend
\section*{TIMOTHEI EPISTOLAM. }
\marginpar{[ p.65 ]}\pstart desiderare autem honorem, eminentiam, census, pom pam, dominia, luxum, otium, castra, uenationes, regna, non est opus desiderare, multo minus bonum opus. Tales hodie sunt qui defiderant episcopatus, ambiunt, uenantur, emunt, rapiunt, tempore Pauli non sic. Nihil hi minus quam episcopi, principes sunt et dicuntur, et haberi uolunt, et huiusmodi mundani satrapae imperii. Omnes episcopi esse debemus, uisitare, curare proximos, unicuique de proximo man dauit deus, qui communitatem curat, custodit, nutrit docet, fouet, regit fideliter. Hic uero episcopus omni no nomen functionis et operis et officii est, non ho noris et eminentiae, sicut antistes et praepositus ac prior Επίςηοzet Graece.  \pend
\phantomsection
\addcontentsline{toc}{subsection}{\textit{Oportet igitur episcopum irreprehensibilem esse. }}
\subsection*{\textit{Oportet igitur episcopum irreprehensibilem esse. }}\pstart Oportet dicit, necessitatis est et praecepti Apostolici, imo domini. Quid hodic facerent, qui sua tueri uolunt, exuerbis dei, si hic uel alibi inuenirent Oportet episcopum habere decimas? Mandatum dicerent sub poena pecati mortalis uel damnationis aeternae. Nunc autem oportet episcopum esse irrepre hensibilem, non sine peauto, pro quo nemo ab alio reprehensibilis sit qui omnes pecatores sumus et offendimus omnes. Sed ut nihil criminis ei imping  \pend
\section*{COMMEN. IN PRIOREM }
\marginpar{[ p.66 ]}\pstart possit, nulla saecularis macula, non aliqua notatus infa mia. quo modo minus aptus fit. Obsecrare, increpare, arguere libera fronte. Hic necessitas uidetur tol lendi concubitus, quo alioqui optimi parrochi, uiri do cti, bengni, solliciti, deumque timentes, hoc solo uitio reprehensibiles. Episcopi sunt omnes ecclesiarum praelati, non solum magnarum, sed etiam uel minimarum, ha beantur irreprehensibiles, innoxii, integri famae, sed haec clarius sequuntur.  \pend\pstart Vnius ixoris uirum ) Sancti patres plures simul ha bebant uxores, filiorum multitudo bonorum domini bene dictio est. Christus hoc nunquam improbauit, christianis diuinae literae multitudinem nusquam interdicunt uxorum, tantum hic episcopis et diaconis prohibetur uxorum pluralitas huiusoemodi. Et cum statim praemitta tur, oportet, non sine causa Chryfosto. periclitatur, ne Paulus uideatnr exigere matrimonium, oerte ut non ten tantur a satana, ne infamiam facile inadant, ut discant fa miliam regere, ut hospitales esse possint, ut studioscri pturae domestida cura non distrahantur, ut aliae mulieres non sit necessaria familiaritas, sunt quae in domo sine foemina expedire nequeant. Cohabitatio uiri et mulieris extra matrimonium non potest carere incommodis, amimae et famae et rei familiaris etc Vt interim omnium scandalorum sandalosissimum sandalum  \pend
\section*{TIMOTHEI EPISTOLAM. }
\marginpar{[ p.67 ]}\pstart cncubinatum illum epi coporum, id est plebanorum, fon tem omnis Christianae iacturae et damnationis et infelicitatis in orbe, in omnibus statibus, quibus abunde consultum esset, si maritorum lege uiuerent sacerdo tes, inuenirent honestam, proborum et diuitum (si uel lent) puellam, nutritam in disaplina, fidelem marito, fi liis, domui: familiae praeesset utiliter et cum parsimonia reponeret pro hae redibus filiis, quae de uictu et amictu superessent. Disceret a marito uerba dei, do oeret paruulos, honoris consequendi intuitu, quem indubie sperarent prae fliis aliorum. Esset amabilis omni bus et honorabilis. Sacerdos eam pudice diligeret et prolem, uacaret quiete literis, doceret sine timore cum fidutia omnem ueritatem, subditos diligeret, diligenter corrigeret, haberet fauorabiles, beneficos, promptos et obedientes doctrinae etc bona infinita sequeren tur. Nunc autem proh dolor et pudor saoerdos parro chus quotidie uel saepe celebrans et sacramenta mini strans, et publicum scortum nutriens, de decimis, bla tionibus, eleemosynis, nomne se damnationis mancipi um nouit certissime crimine et publico scandolosus? Non uult nec potest dimttere nisi alia surroganda infeliciori miseria. Orationes suas si dicit, illic omnia contra conscientiam et interminantia imperium offendit, cum nausea et horrore omnia. Odit  \pend
\section*{COMMEN. IN PRIOREM }
\marginpar{[ p.68 ]}\pstart sacrarum litteraaum studia, otio solo sese et uitiis accommodat, uritur importunius, grauatur libidine insatiabili, temporibus religoni consecratis. Sequitur conscieniiae uermis, uel desperatio infernalis, fit ex hoc procliuior ad omnia uitia. Confunditur coram deo, conscientia, populo et familia. Tractat sacramem ta ut porcus. Qui d interim agit infelicissima meretrix? Filios forsan generat. si non aborsus procurat: deserere filios et scortatorem non potest, non fert redire ad parentes uel consanguineos, non patrem inueniet, non matrem assequetur. Nouit se despectam omnibus, etiam a filiis facile se nouit eiiciendam, si li bidinosior superuenerit. Miserrima post obitum saoerdotis, ideo sua maxime curat, furatur, exigit, et praeparat, uesteis et multas et meretricias, ut post illum uel cum illo aliis placeat, quos facile admittit, scens se non ob id fore adulteram. Filios non amat, negligit, non docet, de animae suae salute desperat, ad turpitudines et uitia procliuis nimium. Vicinis uiris et foeminis odibilis, rixosa. onerosa, dum se contemeni non patitur, dum spetiosior et ornatior incedit, quam honestae matronae, dum irritat sacerdotem inuicinos, se quoquomodo dehonestantes uel etiam negligentes. Filios indulgentius nutrit, et fouet, non cohibet insolentes. Ingerit se primam coetui foe\pend
\section*{TIMOTHEI EPISTOLAM. }
\marginpar{[ p.69 ]}\pstart minarum, iactat conditionem et uictum ac libidinem Filii eorum nutriuntur inter scandala, ut non sintam plius scandala. A parentibus ut illegittimi negligun tur, sine spe consequendi honoris, et extorres hone statis prorsus. fiunt ex hoc rebelles, indisciplinati, blasphemi, rixosi, lusores, etiam parentum contempto res. Non ferunt disciplinam, nec uerbo, nec facto. Ab omnibus despiciuntur ut spurii. Nihil facile boni eis alicubi confiditur, peiores omnibus coaetaneis, audaoes ad omne flagtium, diuinorumomnium auersatores Expelluntur ante tempus in exteros, libertatem captant ad nephas omne, cuius semina susceperunt ex educatione, domoque ; paterna. Quid populo acadit? Intuetur miser dei populus et uulgus, plebis sacerdotem et pastorem, et episcopum suum, speculum omnimodae abhominationis. Verbum dei si audit, non credit, contemnit, magis detrahit, blasphemat, secta tur pastorem per abrupta gradientem in praecipiti um. Correctiones non suscipit, ridet potius cum culpa redarguit doctorem, incalescit magis et dolum suspicatur et subuersionis imposturam. Sancta dogmatd non internoscit, sed abiicit omnia. Euitat ecclesiam, non obtemperat iussis, fidem ignorat, uitia sectatur, conscientid cauteriatur, Christianismus deperditur, confissio obmittitur, decimae negintur, ae  \pend
\section*{COMMEN. IN PRIOREM }
\marginpar{[ p.70 ]}\pstart remoniae ut sordidae propter ministrum irridentur, in terim male audit sacerdotium et omnis religio Christi ana, uerbum deicontemnitur, non praedicatur, non discitur adulterinae doctrinae surrogantur, uitia palam defen duntur, excusatur error, scandalum iustificatur, obturam tur ora piorum, exquiruntur defensionum causae, decano rum fauores ambiuntur redimuntur muneribus, exco gtantur grauamina pauperum, exiguntur oblationes, ag grauatur stola, peractiones mortuorum commendantur, et purgatorium incalescit que modicis et acadentalibus haec saoerdotis famula uel familia non ualeat enutriri. Turpe lucrum cupidus episcopus imperat, diooesanus uel eius uicarius et censum auget quotamnis. Cogitur ta lis sacerdos maiorum promerert fauores per fas et ne fas, quo subsistere possit, uel inuitis subditis. Offitio suo diuino in publica ecclesia assistunt et arca alta re spurii filii, breuiter exemplum sunt flagitiorum et uela nem malit.ae magistratibus et populi rectoribus. Haec omnia et his plura ac maiora introduxit nobis miserrima necessitas caelibatus et stulta religio contra stimulum calcitrandi, et angelicam puritatem imponendi cis, qui carne et sanguine, uclint nolint opprimantur et ardent. Et paucissimorum priuilegium, tot milibus con mu ie faciunt, reclamantibus conscientia, experientia, deique obedientia. His omnibus abunde consulcretur  \pend
\section*{TIMOTHEI EPISTOLAM. }
\marginpar{[ p.71 ]}\pstart si deiuerbis et praeceptis obaudiretur. Honestissimum matrimonium fuit sacerdotibus legs Mosaicae. Condi tor legis dominus plem et matrimoniumubique et plurimum commen dauit. Patriarcharum nota coniugia et pro phetarum adeo et apostolorum. Christus toto suo euange lio inter tot saoerdotum uitia, nec semel uxorum uel unius uel plurium mentionem habuit, praesciuit omnia, matrimonium commune uoluit omnibus. Et primitiua apostolica ecclesia hunc celibatum nesauit, orientalis ecclesia ignorauit hactenus et ignorat. Papistica nuper orta tyrannis haec animarum naufraga cepit, spiritussancti et uerbi dei flatu ppediem exufflanda tela aranearum istorum, cum pufillis clericis interdicitur ob quaestumquod pinguibus pontificibus est perpetuum ad luxum. In sum ma potuit uideri priscis aliquibus episcopis commodum caelibatus introductio in clerum ut purius uiuerent, di gnius sacrosancta tractarent misteria, expeditius scri£ pturae operam nauarent, terrenis actibus minus inquie tarentur, et ut populus minus grauaretur uel ecclesia sustentatione uxorum et filiorum. Sed cum omnia in diuersum ierint, eoque deuentum fit, ut impurissime uiuant, lud brio et nausea habeant sancta homines irreligosissimi, ludo potius et poculis as idue incum bant quam scripturis, otio uero inertissimo prolabam tur in publica crimina spuriis et meretricibus non  \pend
\section*{COMMEN. IN PRIOREM }
\marginpar{[ p.72 ]}\pstart solum grauetur, sed et publice infametur ecclesia uerbum dei ueniat in contemptum, et clerus in odi um et scandalum toti Christianae religiont maxime pateat. Obligantur sub aeternae damnationis iusussima, certissimaque poena, et ira dei maxima urgere de bet, ut otius et sine dilatione hoc scandalum inexpluibilis perniciei tollature medio, si non concilio ge nerali (quod sustinere nolunt qui lucem odiunt) sal tem prouinciali. Sin minus et hoc cesserit, quilibet episcopus, ut nunc sunt, in sua diocoesi, huic pestilen tissimo morbo suoeurrat, communi consensu et synodo. Quod si neglexerint episcopi, debent, et obligantur sub dei certa maledictione incurrenda alioqui per seipsos, si continere non possunt, experientia et complexione naturali admoniti, propria autoritate sibi uxorem legittimam ducere, publice praemissata mensufficiente instructione plebis, ut sciant dei prae ceptum (et magistratuum habito fauore, si consentiant, alioqui uel resignent officio, uel inuitis omnibus deo acquiescant) et id non solum licere, sed et necessarium esse et ineuitabile dei mandatum. Et non id tam propria autoritate, quam omnipotentis dei iussu necessarium, et sine quo matrimonio saluart neutiquam possint, nec praeesse populo, nec satiffacere uerbo det, nec apostolicae traditioni obedi  \pend
\section*{TIMOTHEI EPISTOLAM. }
\marginpar{[ p.73 ]}\pstart re. Quae hoc loco et alias habet episcopum, pastorem populi in spiritualibus oportere habere unam uxorem, nisi possit capere, quae a solo deo paucis ex rara gracia donantur, et magno periculo seruantur. Cuius denique continentie difficultatem imo impossibilitatem publice testatur abhominabilis spurci tia totius ferme cleri:non solum plebanorum et canonicorum et episcoporum aliquorum, sed et mona chorum et fratrum et monialium et nonnarum pa tens impudicitia, in irrisionem etiam totius status cleri calis iustissimam. Nihil contra hanc castitatem ut nunc seruatur satis dure et seueriter doceri uel scribi po test ab aliquo earum rerum experto: sunt enim omnia quae dicuntur longe atrotiora, quam pro rei dignitate proponuntur. Excessum hunc aliquis tractabit idoneus declamator fusius et ordinatius. Clarissima ue ritas est, negari nequit, potestque de scriptis abun de probari,  \pend\pstart Sobrium ) Non ideo uxor mandatur ut uacet libidini sacerdos, sed ut medetur ei. Minus enim a carne et satana infestantur infra cancellos benidicti a deo matrimonii quam promiscua illa, liberaque libidine palantes ut canes et equi. Quin et hic statim Apostolus ad continentiae remedium adiicit coniugi episcopo, ut sobrius sit. Cibo et potu utatur ad cor\pend
\section*{COMMEN. IN PRIOREM }
\marginpar{[ p.74 ]}\pstart poris non uoluptatem, sed necessitatem parsimoniam seruet ne uel ecclesia expensis grauetur, uel corpus et ingentum, quo semper aptum sit ad spiritualia exertitia studia scilicet sanctae scripturae, praedicatones, consul tationes, con solationes, orationes, lectiones, scriptitatio nes, et c. huiusmodi. Exemplum sobrietatis praebeat nec preciosa, nec multa, nec curiosa fercula, sed sem per utatur cibo et potu pro corporis ratione et offi ciorum executione. Non uno die se exinaniat ieiunio, ut sequenti die denuo oneretur: ne utriusque opus intermitti contingat, sicut nunc faciunt satanice ieiunam tes. Vigiliis fistorum abstinere uolunt, et festis diebus repleri, quando toti spiritualibus se dedere deberent cum sint carnalissimi et ideo omnibus criminibus exer cendis instructi festis diebus, et id religionem putant.  \pend\pstart Prudentem ) Vt prospiciat plebi, cui specu lator praeponitur. Circunspectus, ut praeuenire malum et induaere sancta, quo pacto possit intelligat. Σώφρωνx pudicum dicunt, et quadrat praedictis, ut sit maritus eatenus sobrius ut et pudicus esse possit uisu uerbo, signo opere, et in omni conuersatione sua praeferat pudicitiae exemplar.  \pend\pstart Ornatum ) Non sericis uestibus argento et auro: Haec enlm omnino dedecori sunt sacerdoti, sed pudicis et modesissimis moribus, id quod pu\pend
\section*{TIMOTHEI EPISTOLAM. }
\marginpar{[ p.75 ]}\pstart dicum diximus, zichtieg und sittig. In moribus et factis uenerabilem, et maturitate legittima gratiosum.  \pend\pstart Hospitalem ) Non publicus tabernarius, sed peregrinorum et indigentium benignus susoeptor, et tractator beneficus, ne honesti uiri communi diuersorio sub turba et clamoribus praepediantur a studiis et quiete. Non solum sacerdotum et mona chorum sed piorum et qui indigent. Ad quod certe con uenientissimum est non habere pellicem, sed uxorem fide lem et matrisfamilias et filios disciplinares, ne toti ecclesiae cui praefertur inuratur passim diffamia dum peregrini praeter Christianam charitatem tractantur Ea ratione decimae eis debentur non tam pro se, fami lia, hospitalitate, quam pro orphanis et egenis infir mis et senibus pauperibus etc.  \pend\pstart Doctorem ) Appositum ad docendum, qui possit dooere, non alterius doctrinae quam diuinae sacrae scri pturae sapientiae utilis ad dei cultum. Non permotum in gymnasio, non empta facultate, sed spiritu sancto, de sacris edoctis instruat plebem cui superintendere debet ex officio. Non uenator, eques, lusor, sed uerbi dei doctor. Hoc prae oïbus aliis suis officiis peculiare est episcopo pascere oues Christi quas amat, non philo sophiam, non grammaticam, non poësim, non humanas consti tutiones, sed ea quae pertinent ad dei cultum, non damnet non blasphemet, sed doceat.  \pend
\section*{COMMEN. IN PRIOREM }
\marginpar{[ p.76 ]}\pstart Non uinolentum ) Non uinosum, qui uini potui indulgeat nimis, et impediatur a modesta in facto et uerbo doctrinae et corripiat cum importu nitate, aut conuiciis more ebriosorum. Non combibitorem, qui in aliis haec debeat corrigere. Obtun dit ingenium, inepte laetificat consilii inopem. Foedum et iracundum pariter facit uinum immodice sumptum Etiam pudicitiae et castitati aduersatur uinositas.  \pend\pstart Non percussorem ) Non promptum ad uindicandum, impremeditatum ad ulciscendum, sed potius animo tranquillum, longanimem, sustinentiam habentem in iniuria. Non paratum ad arma inferen da, sed in bono malum uincat, percutiat neminem, uel no oeat. Euangelicam patientiam monstret euangelii prae o et doctor Christianae mansuetudinis.  \pend
\phantomsection
\addcontentsline{toc}{subsection}{\textit{Non turpiter lucri cupidum. }}
\subsection*{\textit{Non turpiter lucri cupidum. }}\pstart Hoc audiant episcopi nostri, qui lucrum capiunt ex casibus reseruatis ex concubinatu sacerdotum, ex sacramentorum administratione, ex spuriorum progignitione, si non per se talia capiunt, sed tamen subornant officiales, quibus id rei foedae committunt, uel consentiunt ob lucrum. Communius tamen est turpi ter lucri cupidum, quam turpis lucri cupidum. Omne lucrum cupere episcopum dedecet, qui expendere debet potius quam lucrari et thesaurizare. Bona  \pend
\section*{TIMOTHEI EPISTOLAM. }
\marginpar{[ p.77 ]}\pstart enim episcoporum ecclesiae sunt, indigentium de ec clesia, reponi in thesauros non debent: multo minus turpe lucrum uenari.  \pend\pstart Sed aequum, modestum. 7 Mansuetum potens aequanimiter ferre infirmorum et detrectan tium et dehonestantium, non acutum ira, uel ad uin cendum propensum gerens animum.  \pend\pstart Non litigiosumq Vel contentiosus bonicon sulat uerba et fasta etiam contentiosorum, sed paci ficum Non bella et lites suscitet, sed omni diligentia componat. Quomodo autem litibus carebunt, qui diuitiis inhiant et possessionibus? Dirimant po tius contentiones. Eras: Alienum a pugnis. Vide Ann. Non cupidum 7 Vel alienum ab auaritia, ne quaestum quoquo modo sectetur, ut simoniaci. Liberalem in egenos, non lanae potius abrasor, quam uerbi dei administrator. Ab omni specie cupiditatis abstineat.  \pend
\phantomsection
\addcontentsline{toc}{subsection}{\textit{Sed suae domui bene praesit. }}
\subsection*{\textit{Sed suae domui bene praesit. }}\pstart Quomodo sine uxore bene praeerit, et cum spu riis uel externo famulicio. Vult Apostolus ut familiam suam, uxores, prolem, famulos bene moderentur, ne scandolosa domus toti ecclesiae praeiudicent. Si non habet familiam, saltem canonicis suis praesit antistes quos hodie nouimus iniquissimo iure exempros ab  \pend
\section*{COMMEN. IN PRIOREM }
\marginpar{[ p.78 ]}\pstart episcopo, et omnium scanda lorum antesignani et sa aerdotum totius episcopatus exemplaria insolentiarum.  \pend
\phantomsection
\addcontentsline{toc}{subsection}{\textit{Filios habentem subditos cum omni castitate. }}
\subsection*{\textit{Filios habentem subditos cum omni castitate. }}
\phantomsection
\addcontentsline{toc}{subsection}{\textit{Erasmus. Liberos habeat in subiectione cum omni reuerentia. }}
\subsection*{\textit{Erasmus. Liberos habeat in subiectione cum omni reuerentia. }}\pstart Disciplinatos habeat filios, qui moribus religiosis et castis, et reliquorum filios instituant, et exem plum honestatis exhibeant. Obedientes patri, reuerentiam exhibentes. Quid hodie de filiis presbyterorum iudicandum? Debet in propria sua domo discere exer atio filiorum quod docet in ecclesia subditorum, un de et sequitur:  \pend
\phantomsection
\addcontentsline{toc}{subsection}{\textit{Quod siquis propriae domui praeesse non nouit, quomodo ecclesiam dei curabit? }}
\subsection*{\textit{Quod siquis propriae domui praeesse non nouit, quomodo ecclesiam dei curabit? }}\pstart Quomodo curabit multos, et quasi alienos, qui propriis prodesse non potest? Et tamen ad id positus est, scilicet ut prosit, qui in paruo infidelis, quomodo in plurima constituetur? Qui familiar bus non prospi cit et necessariis, quid alienos curabit? Qui sensu tam tenui ut pauca cogitare nequeat, quomo do ece lesiam totam reget? I deo prius exercitari debuit in doma propria, in liberorum moderatione, ut regimen populi ualeat susinere.  \pend
\section*{TIMOTHEI EPISTOLAM. }
\marginpar{[ p.79 ]}\pstart Non nouitium'I Nouus infide, nuper con uersus, nondum exercitium spiritus inse expertus. Cum enim alios docere et moderari debeat in his quae ad deumeexpedit exercitatum habere sensum in spiri tualibus Nemo autem repente fit summus, et studium sacrarum litterarum non subito proficit: ueruntamen id non a baptismatis tempore computandum, quod mul ti olim tardius baptizabantur, sed a conuersionis tempore ad fidem. Ambrosius enim cathecuminus longo tempore Christianus erat antequam baptizare tur et episcop. nondum lotus eligebatur spiritus pmotus.  \pend
\phantomsection
\addcontentsline{toc}{subsection}{\textit{Ne in superbiam elatus in condemnationem incidat calumniatoris. Oportet autem illum bonum habere testimonium ab his qui foris sunt, ne in opprobrium et in laque. }}
\subsection*{\textit{Ne in superbiam elatus in condemnationem incidat calumniatoris. Oportet autem illum bonum habere testimonium ab his qui foris sunt, ne in opprobrium et in laque. }}\pstart Non expedit nouitium religionis subito principem et maiorem deligere religionis, ne uel religio uitupe retur, quasi ut destituta idoneis pro functione huiusmodi personis, quam pbi et sapientes omnes declinent, et cogantur aliunde pellicere pmouendos, uel ipse no uitius honoris sui sortem ferre humiliter non potens, efferat se insolentius et statim sibi licere putet quae libuerint, alienique ; ansam captent calumniandi quasi ea ratione acesserit ecclesiam ut honoraretur.  \pend\pstart Incondemnationem incidat calumniatoris Facile  \pend
\section*{COMMEN. IN PRIOREM }
\marginpar{[ p.80 ]}\pstart bonos inuidiam patitur, inuidia uero temerarium iu dicium parit et calumniam, cum quicquid a neophy tis agitur, quantumlibet bene, insimulatores perpetuo omnia eorum bona inuertentes, cum nihil aliud possunt, nouiciatum obtendunt. Sed et facile praesu mant, nimium et audenter nout et insueta, eleuanturque existimatione propria, sicque incidunt in iuditium calumniatoris.  \pend\pstart Oportet autem illum ) Hoc mihi ad episcopum quemlibet referendum uidetur, ut scilicet etiam ab his qui foris sunt, et ab omnibus bonum testimonium habeat, ne religio ob ministri improbitatem, calumniam sustineat. Specialiter tamen si a gentilibus uel Iu daeis conuersus sit ad fidem, omnino necesse est,ut il li eidem nihil inhonestum et turpe impingere possent. Vt non dicant: Ece quales preficiunt sibi Chri stiani, quales nos nouimus improbos, hos ipsi honoribus dignantur. Flagitiosus erat nobiscum, nunc bonoratus apud Christianos. Ecr quales religiones quales ministri. Sicque in laqueum et in opprobrum incidit calumniatoris. Sequitur:  \pend
\phantomsection
\addcontentsline{toc}{subsection}{\textit{Diaconos similiter pudicos, non bilingues, non multo uino deditos, non turpe lucrum sectantes. }}
\subsection*{\textit{Diaconos similiter pudicos, non bilingues, non multo uino deditos, non turpe lucrum sectantes. }}\pstart Ministri episcoporum diaconi uocati sunt, ut in  \pend
\section*{TIMOTHEI EPISTOLAM. }
\marginpar{[ p.81 ]}\pstart Actis. Non quod praedicarent uel euangelium proferrent, sed quia oblationes fidelium distribuebant. Seruiebant ecclesiae et episopo, et de oblationibus uiuebant sicut episcopi: sicut nunc sunt coadiutores plebanorum, quatinus liberius episcopi uacarent miristerio uerbi in docendo. Talis fuit Laurentius, et in primis Stephanus, qui tamen acepta oportunita te non taoebat a uerbo dei praedicando. Quia autem ministri erant ecclesiae, et ab ea sustentabantur sicut et uiduae, quae uere erant uiduae, id est, destitutae so latio filiorum et orphani etc, ideo prudentes uult esse, et omni honestate conspicuos, uenerandae conuersationis, in uerbo, facto, et exemplo: ne opinio mala dehonestaret eorum offitium, ne diuidentes collectas eleemosinas incontinentiae suspitionem incurrerent, si mulieribus pauperibus dona ecclesia largirentur.  \pend\pstart Non bilingues I Ne scilicet quaestus gracia uarie loquantur de Christianis dogmatis, ne aliud affirmarent coram episcopo, aliud coram simplicibus, neque uerbum dei adulterent, et ad graciam loqua tur, quatenus tanto magis collectarum comporte ut, et ne obloquantur episcopo, uel ambiendo episcopatum graciam uerbis fucatis dupliabusque ; mercentura uulgo.  \pend
\phantomsection
\addcontentsline{toc}{subsection}{\textit{Non multo uino deditos. }}
\subsection*{\textit{Non multo uino deditos. }}
\section*{COMMEN. IN PRIOREM }
\marginpar{[ p.82 ]}\pstart Ne luxu soluantur, a debita grauitate, ne dedeco ri sint episcopo, ecclesiae, et officio suo, ne obeum ma£ le audiat episcopus et ab operibus minime fideles retrahantur ob dispensatoris improbitatem.  \pend
\phantomsection
\addcontentsline{toc}{subsection}{\textit{Non turpe lucrum sectantes. }}
\subsection*{\textit{Non turpe lucrum sectantes. }}\pstart Non enun turpe lucrum expendi debet in eleemosynam, sed honeste parta. Non excogitent modos ad eliciendum a fidelibus supra neoessitatem ecclesiae. Praeuidebat enim in spiritu bonorum temporalium affluentiam irrepturam in ministros ecclefiae in perniciem, quam interim uidemus creuisse in immen sas abusiones.  \pend
\phantomsection
\addcontentsline{toc}{subsection}{\textit{Habentes mysterium fidei in conscientia pura. }}
\subsection*{\textit{Habentes mysterium fidei in conscientia pura. }}\pstart Cum passim et quasiin augulis seorsum instrue re deberent eos, quos eleemosynis reficiebant, omnino fidei Christianae mysterium probe sare debebant, et reddere rationem eorum quae ab episcopis docbantur respondere de fidei articulis, sicut Stephanus assiste re episcopo, defensare, et testimonium ueritat is et uitae praestare ibi uidebant id utile et honorificum Christo et euangelio et ecclesiae parati inquam semper de fide respondere. Si uero sine omni fructuuiderent, fiuerem tur saltem pura et sinoera consaentia sua certi coram deog dei gloriam quaererent et profectum fidelium, et  \pend
\section*{TIMOTHEI EPISTOLAM. }
\marginpar{[ p.83 ]}\pstart quod in se essent parati implere: si quo minus promo ueretur negotium fidei ipsi candidae egissent functionem suam et fideles ministri dei et ecclesiae inuenirentur,  \pend
\phantomsection
\addcontentsline{toc}{subsection}{\textit{Atque hi probent prius et sic ministrent. }}
\subsection*{\textit{Atque hi probent prius et sic ministrent. }}\pstart Post episcopatum primi honoris habebatur diaconus, quasi episcopi et ecclesiae procurator fidelis, di oerem hodie decanorum camerarios agere, sed longe miserabiliore comparatione. Oportebat ergo probari fidem eorum et probitatem ac sanctam industriam, si con mitti eidem possent tuto distribuendae collectae, et si idoneus minister etia uerbi, episcopo concors fide et moribus, et dignus ecclesiasticorum bonorum admini strator, quod non debebant passim et sine delectu commendari cuiuis. Probandi igitur erant ministrantes.  \pend
\phantomsection
\addcontentsline{toc}{subsection}{\textit{Nullum crimen habentes. }}
\subsection*{\textit{Nullum crimen habentes. }}\pstart Derogasset criminis diffamia epi scopo, et ecclesiae fideles absterruisset a fide et minime operibus iacturam aoeepissent pauperes ecclesiae, fadlcque ; fraudem suspicari contingeret in dispensatore et procuratore crimi nos fa mil amque et statum episcopi dehonestusset.  \pend
\phantomsection
\addcontentsline{toc}{subsection}{\textit{Mulieres similiter pudicas, non detrahentes, sobrias, fideles in omnibus. }}
\subsection*{\textit{Mulieres similiter pudicas, non detrahentes, sobrias, fideles in omnibus. }}\pstart Videtur sine dubio de uxoribus diaconorum hoc dia, quod eadem ratio sit hic quae prius de episcporum familia: quia enim de ecclesiae bonis et oblatiombus  \pend
\section*{COMMEN. IN PRIOREM }
\marginpar{[ p.84 ]}\pstart uiuebant, dect bat maxime eas tanquam diaconorum collegas, pudicas esse, honestae conuersationis, quietas, pacificas, modestas, sine arrogantia et dissoluti one aliqua in moribus, ne ecclesia turbetur et scandalizetur.  \pend
\phantomsection
\addcontentsline{toc}{subsection}{\textit{Non detrahentes, non calumniosas. }}
\subsection*{\textit{Non detrahentes, non calumniosas. }}\pstart Rixolas, discordiam seminantes, delatrices et mendosas. Hoc enim uergeret in turbationem ecclesiae cum episcopi et diaconi a talibus omni diligentia dehortari publice et priuatim debebant.  \pend\pstart Sobrias. vt prius de diaconis et episcopis dictum. Omnino enim foedum in muliere temulentia et uini abusus, mater turpitudinis et libidinis. Non decebat talium uitiorum exemplar relinqui subditorum uxoribus a foeminis ecclefiasticis.  \pend
\phantomsection
\addcontentsline{toc}{subsection}{\textit{Fideles in omnibus. }}
\subsection*{\textit{Fideles in omnibus. }}\pstart Vt cum ministrorum ecclesiae adiutrices eas esse oporteat, similiter ut diaconi omnia fidelissime dispensent, nihil pro se per fraudem detrahant, non po nant fid elibus offendiculum cum de ecclesia uiuant ei£ dem quam fidelissime ministrent.  \pend
\phantomsection
\addcontentsline{toc}{subsection}{\textit{Diaconi sint unius uxoris uiri, qui filiis bene praesint et suis domibus. }}
\subsection*{\textit{Diaconi sint unius uxoris uiri, qui filiis bene praesint et suis domibus. }}\pstart Non patientur se notari de incontinentia, ut uxores plus quam unam habcant, ut liccbat eo tem\pend
\section*{TIMOTHEI EPISTOLAM. }
\marginpar{[ p.85 ]}\pstart pore reliquis Christianis, uel ad minus tollerabatur tanquam reliquias Iudaeorum, licet Romanis legibus non probabatur: unde et cum tempore abrogata est, huiuscemodi uxorum pluralitas. id tamen nec Apostolus poterat pati in episcopo et diacono: maxime quod insuper etiam grauamina et displicem tia fidelibus contigisset ex tot uxorum nutritione. Nihil hic de digamia uel trigamia Apostolus insinuasse mihi uidetur. Adolescens enim amissauxore pri£ ma, quare perpetuo ureretur? Sicut nec uiduas iuuem culas matrimonio uetabat, sed praecipiebat nubere. Et hic eadem ratio. Ne satanas tentet et uincat, filiis suis bene praesint. Scientia, moribus, et discipli na erudiat filios, ut sint speculum honestatis caeteris adolescentibus.  \pend\pstart Et suis domibus J Omnem familiam honestissime gubernet, eadem ratione qua et episcopus. Videmus quales esse debeant moribus et fide et exercitio, quotquot ecclesiae censibus nutriuntur, qui scili cet in ministerium ecclesiae neoessarii sunt. Ncmo enim nisi ministrantes ecclesiae ab eadem sustentari debent et omnino nemo otiosus qui laborare potest, ex quo ecclesiae utilitas nullaprouenit, et qui non sunt usui in opere a Christo et Apostolis instituto, uel ad hoc insti uantur aliquando idonei ministri ecclesiae. Sunt autem  \pend
\section*{COMMEN. IN PRIOREM }
\marginpar{[ p.86 ]}\pstart neoessarii ecclesiae episcopi, diaconi, scolaftici:uel iuuem tuti instituendae praefecti. Nomino episcopos ut sunt uere, qui uerbum dei tracti nt et docent: diaconos qui discunt et cooperantur episcopis, non uideo alios. Cantus horarum cnonicarum, ut nunc fit, est prorsus inutilis, posset tamen institui pro eo melius et neceFarium aliquid. Monachi omnes laborant manibus, et siuerbo det horis aliquibus perdiem insistant, laudabi le. Temporales passiones deseruisse conforme quidam euam gelio, sic tamen, ut non resardant posthabita mendia tate, sed laborent, et eleemosynas potius dare et fa cere studeant, quam requirere et mendicare. Siquis alius ecclesiae neoessarius minister deoernendus sit, il Ie de ecclesia uiuat, nemo inficiabitur. Sed odosorum prorsus nemo firatur: Qui non laborat, non mandu oet, qui salict laborare possunt. Egens nemo sinatur in ecclesia nec otiosus..  \pend
\phantomsection
\addcontentsline{toc}{subsection}{\textit{Qui enim bene ministrauerint, gradum si bi bonum acquirunt, et multam fidutiam in fide quae est in Christo lesu. }}
\subsection*{\textit{Qui enim bene ministrauerint, gradum si bi bonum acquirunt, et multam fidutiam in fide quae est in Christo lesu. }}\pstart Diacont alicuius ecclesiae si ministerii sui functionem laudabiliter absoluunt, si officium suum cum laude et pro bitate peregerint, si uidebunturmunus suum fideliter obiuisse, non solum laudem, sed et gradum sibi meliorem acquirent. Erunt procildubioido nei ad  \pend
\section*{TIMOTHEI EPISTOLAM. }
\marginpar{[ p.87 ]}\pstart su cdendum episcopo, si non illi cuius diaconi sunt, transferentur eligendi in aliam ecclesiam, uel ex co operatoribus fient fortassis, coëpisopi uel in ciuitate una pro multitudine populi. Episcopt, id est, parrochi constituentur plures, duo, assequeturque apud uni uersam Christi fidelium congregationem, quae salioet est in Christo ecclesiam, autoritatem magnam. Et mul£ tam fiduciam in eundem collocabunt fideles, ut qui tam strennue et fideliter diaconatum administrauerrit dignus euadat maiori gradu.  \pend
\phantomsection
\addcontentsline{toc}{subsection}{\textit{Erasmus: multam libertatem quae etc. }}
\subsection*{\textit{Erasmus: multam libertatem quae etc. }}\pstart Videtur uelle ut diaconus qui a prinapio ministe rii sui omnia ad praescriptum et ordinationem epi scopi agere debebat, nec pro suo iuditio dispensare omnia, sed quasi ad rationem acapiendam ab episoo po uel ecclesia. Si inquam talis diaconus insigni fide dexteritate, diligentia, et uiglantia omnia semper agere cum sua familia uisus fuerit et inuentus, acquiret ex hoc sibi libertatem eidem permittendam ab episcopout omnia uel multa saltem pro arbitratu suo fruatur curare, fidelis toties repertus Christi ob sequio, ut nemo sibi fidem agendorum abrogare pos sit, sed fiducialiter et libere agat quae in honcrem dei et ecclesiae fuerint peragenda. Iste sensus mihi hactenus apparet.  \pend
\section*{COMMEN. IN PRIOREM }
\marginpar{[ p.88 ]}
\phantomsection
\addcontentsline{toc}{subsection}{\textit{Haec tibi scribo fili Timothee, sperans me uenire ad te cito. }}
\subsection*{\textit{Haec tibi scribo fili Timothee, sperans me uenire ad te cito. }}
\phantomsection
\addcontentsline{toc}{subsection}{\textit{Erasmus: Haec tibi scribo, sperans fore, ut ueniam ad te cito. }}
\subsection*{\textit{Erasmus: Haec tibi scribo, sperans fore, ut ueniam ad te cito. }}\pstart Hoc ex I'aodicia scribebat Ephesum, nuper ibi institutum episcopum instruendo, quae sciret, ad id tempus necessaria magis, suo tempore personaliter de omnibus coram edocturus, quae circa episcopi, dia conorumque officia essent ordinanda.  \pend
\phantomsection
\addcontentsline{toc}{subsection}{\textit{Si autem tardauero ut scias quomodo oporteat te conuersari in domo dei. }}
\subsection*{\textit{Si autem tardauero ut scias quomodo oporteat te conuersari in domo dei. }}\pstart Domum dei non templum marmoreum notat, sed quamlibet congregationem fidelium, Christo sinceri ter credentium, et uerbum dei toto corde amplectentium et per opera fidei confitentium semper. Inter quos, non super quos conuersatur episcopus apostolicus, cos ducendo, et insttuendo infide Chri sti, et doctrina sancta.  \pend
\phantomsection
\addcontentsline{toc}{subsection}{\textit{Quae est ecclesia dei uiui, columna et firmamentum ueritaris. }}
\subsection*{\textit{Quae est ecclesia dei uiui, columna et firmamentum ueritaris. }}\pstart Domus illa dei in qua te oeu ministrum et curatorem et superintendentem consituit dominus, congrega tio scilicet omnis fidelium, non tua Ephesina solum, sed ubique eodem spiritu congregata multitudo sancto rum, est reuera ecclesia quam deus condidit, gubernat  \pend
\section*{TIMOTHEI EPISTOLAM. }
\marginpar{[ p.89 ]}\pstart constituat, et dooet spiritu oris sui, uerbo dei. Non enim per se, sed in Christo subsiftit et durat, basisque ́ est et sedes ueritatis. Colitur enim in ea et agnosa tur unus uerus et unius deus, cum omnes alie eccle siae et congregationes hominum, deos falsos et ido la colant. Iudaei denique a ueritate dei, quae est per Christum Iesum aberrauerunt et exoecati sunt et ceciderunt. Nostrae autem ecclesiae doctrina et fides fundamentum est et lapis et columna immobi lis et semper permansurae ueritatis. Ipsi enim ced derunt, nos autem surreximus et erecti sumus credentes in Christum dominum et deumnostrum: uiuum et uerum, qui est ueritas et uita, illamque ueritatem firmissime credimus et praedicamus et usque ad mortem defendimus.  \pend
\phantomsection
\addcontentsline{toc}{subsection}{\textit{Et manifeste (citra controuersiam) magnum est pietatis misterium, quod manifestatum est in carne, iustificatum est in spiritu, apparuit angelis, praedicatum est gentibus, creditum est in mundo, assumptum est in gloria. }}
\subsection*{\textit{Et manifeste (citra controuersiam) magnum est pietatis misterium, quod manifestatum est in carne, iustificatum est in spiritu, apparuit angelis, praedicatum est gentibus, creditum est in mundo, assumptum est in gloria. }}\pstart Si de me mibi quicquam tribuere possem citra hunc locum tam obscurum, et si cogita tione mea non aberrem, distinguerem aliter, euaderemque duas difficultates: primam quae ecclesiam dicit columnam ue  \pend
\section*{COMMEN. IN PRIOREM }
\marginpar{[ p.90 ]}\pstart ritatis. Secunda quae sequitur, et sic ordinarem textum: Vt saas quomodo oporteat te conuersari in domo dei, quae est ecclesia dei uiuentis. Hic distin guerem, et subsequeretur hoc modo: Columna et et firmamentum ueritatis, et manifeste magnum est pietatis mysterium. Deus manifestatus est in carne iustificatus est in spiritu, uisus est angelis, praedicatus est gentibus, creditus est in mundo, assumptus est in gloria. Vel si mauis sequi aliam lectionem Chrysosto.et Theophi gnere neutro, et obmisso qui deus. Lict enim priora intelligi possint ut diximus de ecclesia et columna simplicor tamen et indubius sensus est, ut ueritatem non fundari in ecclesiam nec fulariuel firmari, sed potius dei ecclesiam inniti, fir mari et su per ueritatem ipsam dei uerbum Christum et petram ec clesiae et non econtra. Faalis igitur et oertus sensus et insignis plane apparet, si sic lectionem ordinemus Vt scas quomo do te oporteat conuersari in ecclesia illa tibi commissa, hominum salict congregatione, non infidelium gentilium, dolorum, qui mortuos deos et nihilum colunt: sed inter homines qui credunt in deum uiuum et uerum, qui sunt ecclesia dei uiui tibi ut ministro eius commendata. Et sequitur Columna et firmamentum ueritatis. Et manifeste magnum est pi£ etatis sacramentum, quod nunc tibi prodo, quasi di\pend
\section*{TIMOTHEI EPISTOLAM. }
\marginpar{[ p.91 ]}\pstart oeret omnis Christianae doctrinae et atholicae ucri tatis columna, cui innitur, et firma mentum cui subsistit, et mysterium pietatis et fidei. Hoc unicum est magnum ualde, et citra controuersiam magnum quod paucis uerbis tibi describo, ut in ecclesia praedioes summamfidei et basim: deus manifestatus est in car ne, uel mysterium pietatis diuinae in Christo nobis ex hibitum, cum olim multiphariam aparueriti deus, Guerbum deiangelis specie hu ma na, igne, spiritu, fa cibus, nube et oetaera. Nunc maximo pietatis et cle mentiae suae misterio decreuit manifestari, nobis uel mundo in carne, incrnationem sui faciens in Christo ut uisibilis deus hominibus in terra, aliquando aparuet uisibilis, carne circundatus.  \pend\pstart Iustficatus in spiritu ) Per fpiritum sanctum datum hominibus, iustus, deus aparuit, ut iustificans per gradam fidem per uerbum suum, qui in spiri tu sancto uoluit omneis peaatores iustif cari quos omneis damnauerat, nunc per spiritum suum iustificat et ipse iustus innotescit, et pius qui durus prius uideba tur, ut scilict iustifioetur in sermonibus guincat cum nostro sensu carnali iudictur seuerus et iniustus. Apparuit angelis ) Dei gracia misterium pictatis euangelium a saecilis ignotum angelis, per Christum hoc tempore innotuit prinapatibus et potestanbus in caele  \pend
\section*{COMMEN. IN PRIOREM }
\marginpar{[ p.92 ]}\pstart fibus, qui hactenus ignorauerunt tantam dei ad bo minis pictatem et in Christi misterio didicerunt, quae prius abscondita eis erant et ignorata, quae per euam gelii praedicationem et per ecclesiam eis clarior in notuit, Ephesiorum 3.aap.  \pend\pstart Praedicatum est gentibus IOmnibus populis hoc pie tatis misterium et euangelium regni praedicatum est. Dela ta est haec tanta dei gracia omnibus hominibus nemi ne excluso. Christus saluator datus est omnibus genti bus pignus et arra salutis: non solum Iudaeis, sed et genti bus coepit praedicari, et praedicatum est hoc misterium.  \pend\pstart Creditum est mundo I Non solum Iudaeis, sed orbi uniuerso, contra carnem et sanguinem, cui hoc misterium non reuelatur nisi a patre caelefti, qui trahit ad se per Christum quos eligit et donat fidem, quos iustificare decreuit per fidem in Christum.  \pend\pstart Assumptum est in gloria) Glorificatum est, et laudatum ubique et diffusum inomnem terram. et post agnitum misterium etiam glorificatum est, Christo ascendente in caelum, et ad dexteram sedem te, et euangelium illucescente cum gloria, qua omnis prudentia, potentia et iusticia oppressa iacet, ut illud solum et unicum gloriosum sit in caelo et terra Hoc fidei misterium et pietatis sacramentum omnis ueritatis sedes, basis, columna, fundamentum, cui in\pend
\section*{TIMOTHEI EPISTOLAM. }
\marginpar{[ p.93 ]}\pstart nititur omne quod scimus, credimus et operamur Hoc siquis credat, ueneretur et amplectatur, iustus est et perfectus Christi seruus et Christianus.  \pend
\endnumbering\beginnumbering\section{CAPVT IIII.}\pstart \huge\textbf{H}\normalsize Anc summam dostrinae teneas, doceas, inculcare studeas, Christianae fidei breuissimos articulos, deum scilicet incarnatum, uel in carne manifestatum, per spiritum omnia iustificantia, non lege carnalis literae. Angelis etiam admirabilem nobis dei graciam collatam, post factum innotuit. Euangelium Christi et recondliatio dei prae dicata est gentibus sine delectu, mundo creditum est, domino miserante et hac credulitate et fide omnibus iusticiam dante. Et assumptum est hoc pietatis euangelium et beneficium in gloriam, in toto orbe, in terra et caelo, ut eo nemo fraudetur fidelis. Tale aliquid et prius aeque breuiter descripsit, cum diceret:unus deus, unus etiam conciliator dei et hominum homo Christus Iesus, qui dedit semetipsum precium redemptionis pro omnibus. Haec columna fidei et euangelium regni, praeter haec alia non quaerenda non docenda, non inculcanda fidelibus.  \pend
\phantomsection
\addcontentsline{toc}{subsection}{\textit{\huge\textbf{S}\normalsize piritus autem manifeste dicit, quod in posterioribus temporibus desciscent quidam a fide. }}
\subsection*{\textit{\huge\textbf{S}\normalsize piritus autem manifeste dicit, quod in posterioribus temporibus desciscent quidam a fide. }}
\section*{COMMEN. IN PRIOREM }
\marginpar{[ p.94 ]}\pstart Spiritus Christi qui in disaipulis adhuc manet, per quem futura nouerunt, prophetant quoque de antichri stis, qui a Christo abstrahunt, et a fiducia in redmpto re nostro nobis donata. Ille spiritus prophetiam reuelauit futuram, post haec nostra tempora, ut parum curaturi sunt homines etiam nomine Christiani, de fi de in Christum. Non agnoscent misterium illius pietatis diuinae, se putabunt non Christi spiritu per fidem iustificatos. Non erit gloriosum apud eos quod deus euangelium graciae eis transmisit, sed in operi bus legis iustificari uolent, fi non solum legem Mosi, sed et plusquam Mosaicae legis praecepta seruanda proponant et contraria magis, cum Christus fidem in deum et perfectam fiduciam docuerit mediumue rae iustificationis esse, non ex operibus quae nihil uale ant adueram pietatem obseruare, uidelioet docentes delectum ciborum, temporis, uesfium, locorum caelibatus, his rebus iusticiam collocantes. Posteriora tempora melius dicitur quam nouissima, ne quis putet in fine mundi haec olim tantum futura, cum etiam tempore apostolorum, et interim semper speudoapostoli et sophistae et superstitiosi, haec credi derunt, seruauerint, ct docuerint Desasoere a fide non unius dei, sed a Christianismo, a uerbo dei, ab agnitione salutis a Christo nobis paratae, et libertate spiri\pend
\section*{TIMOTHEI EPISTOLAM. }
\marginpar{[ p.95 ]}\pstart tus filiorum dei a cultu dei, quo uere uult coli in spiritu, fide, spe, et charitate et operibus caritatis.  \pend
\phantomsection
\addcontentsline{toc}{subsection}{\textit{Attendentes spiritibus erroris, et doctrina daemoniorum in hypocrisim ( simulationem) loquentium mendacium, et cauteriatam habentes con scientiam. }}
\subsection*{\textit{Attendentes spiritibus erroris, et doctrina daemoniorum in hypocrisim ( simulationem) loquentium mendacium, et cauteriatam habentes con scientiam. }}\pstart Cum ipsissimo uerbo dei, e doctrinae Christi, qui ueritas est adhaerere deberent ipsi erroris spiritui et humane praesumptionis superstitiosis traditionibus attendunt: quicquid enim secus uiam est, uel a uia temdens erroneum est. Spiritus autem Christus uia est, cuius uitam, mores et uerba aemulantes non errabimus unqua  \pend\pstart Doctrinis daemonioru l Daemones legimus docu isse primum quidem ut uerbo dei Eua non crederet, suo proprio abundaret, praeter dei scriptum saret et age ret, docuit idololatrarum cultum et mores, sacrificia, cae libatum. Vestales et eunuchos suos parauit. Christo quoque suadere tentabat de cibo, arrogantia, supstitione et idololatria. Tales doctrinae daemoniacae sunt, diui nae autem sunt disoere mititatem, humilitatem, fidem in deum spem et charitatem in deum et proximum Quicquid his obuiat uel impedimento est, daemoniacum sit neoesse est.  \pend\pstart Loquenttum mendatium per simulationem IQui religiosiores uideri uolunt, adinuenientes de corde suo doctrinas atque regulas uiuendi proprias  \pend
\section*{COMMEN. IN PRIOREM }
\marginpar{[ p.96 ]}\pstart quasi Christus insufficienter docuerit, et cum Aposto lis non satis religose uixerit, et sibi praeter traditio nem euangelicam quaedam alia ordinandaet instituë da commiserit, cum id de suo corde confinxerint:uolen tes uideri aforis sancti cum intra sint spurcitia reple ti, ut hypocritae omnes solent, qui plura uolunt osten tare quam Christus iusserit, obmittentes necessaria, Et cum obmittant quae Christus docuit et praecepit supererogatione tamen consilia sibi praefigunt, quae nec ipsa complent. Vnicum studium habent conftitu tiones perpetuo ordinare, nihilque ; seruare.  \pend
\phantomsection
\addcontentsline{toc}{subsection}{\textit{Et cauterio notatam habentium consciem. }}
\subsection*{\textit{Et cauterio notatam habentium consciem. }}\pstart Conscii sibi sunt hi doctores, quod quae alios docent facere, ipsi non possunt, uel nolunt, uel necesa ria non sunt, et quae Christus non docuit. Ea ratione introducentes supra uel praeter, uel contra Christi doctrinam, ut emungatur aliquid, ut territi, et augu stiati, promptius ad uota obediunt, uel ut magis hono rantur, uel ut sanctuli uideantur, fingentes se facre, quod alios docent et praecipiunt. Omnibus his mo dis consaentia talium cauteriata deprehenditur. De us autemt non irridetur, sed potenter tormenta intentat talibus. Et iuditio percutienttis serui et mactantis in absentia domini iudicabunt, portioque ; eorum erit cum infidelibus.  \pend
\section*{TIMOTHEI EPISTOLAM. }
\marginpar{[ p.97 ]}
\phantomsection
\addcontentsline{toc}{subsection}{\textit{Prohibentium nubere, abstinere a cibis quos deus creauit ad percipiendum cum graciarum actione fidelibus. }}
\subsection*{\textit{Prohibentium nubere, abstinere a cibis quos deus creauit ad percipiendum cum graciarum actione fidelibus. }}\pstart Quid illi daemoniad praedicatores, et cauteriats consaentia dooeant per ordinem designat. Primo omnium docuerunt, non quidemut nemo matrimonium ineat, hoc enim fuisset impossibile et contra scripturam totam. Non oportet quod pecatum sit omnibus, sed religiosioribus Christianis quibus sic celibatum suaserunt ut perpetuo seruandum uouerent, quasi rem turpem matrimonium habentes, et indignum tractantibus sacramenta, et sine pecato magno fieri non posset. Huius doctrinae fundamentumcum nullum sit in sacra scriptura, cum multorum malorum causa sit, cum contra Christi et Apostolorum doctrinas, merito dicitur diabolica. Cum nubere, deus, natura, scriptura, usus, neoessitas approbent omnia, tumque sanctum et deo placitum, nullis prohibitum, neque aliquibus non necessarium, nisi special ssima et rara gracia praeuentis ut non urantur. Prohibite autem nuptiae, ineffabilium malorum causam praestant, quorum nemo nisi diabolus pater et patrator et causa esse potuit, id quod Paulus manifeste saiuit, docuit et scripsit: ideo piaculum grande est prohibere praecepto nuptias.  \pend
\section*{COMMEN. IN PRIOREM }
\marginpar{[ p.98 ]}\pstart Abstinere a cibis quos deus creauit ad percipien dum cum graciarum actione ) Abstinere a crapula et ebrietate, non onerare corpus et animam, non reddere imbecillum ad functiones necessarias deus praeoepit. Iudaeis denique genera ciborum quaedam pro hibuit, pro tempore legis et cohabitatione pacific, ad abominationes cauendas, ad praeueniendum uale tudinem aduersam corporis, qui alioqui non bene con plexioni hominum conueniunt, et idololatrarum su perstitionem testabantur, hoc deus prohibuit. Nunc autem in nouo testamento spiritus sanctus reuelauit, et Christus docuit, omnia munda mundis, nihil quod intrat in os coinquinat hominem. Cibi non insani, de se immundi non sant, si cum gracarum actione acpiun tur, deoque ; donante in escas hominibus, omnes come di possunt, ut nemo eos prohibere possit. Si tamen ob castigationem corporis et parcimoniam tenendam ali quis libere abstineat, non quidem malefadet, tantum ne iin nundum reputet et illicitum et magna oerona dignam huiusmo di abstinentiam. Possumus enim abstinentiam seruare non solum in pisabus et legumi nibus sed etiam in carnibus. Non enim refert quid, sed quantum comederis. Potest in leguminum esca peaari potest omnis generis caro licite manducari. Quicquid comedimus cum gradarum actione manduce\pend
\section*{TIMOTHEI EPISTOLAM. }
\marginpar{[ p.99 ]}\pstart inus, agnoscentes largitorem bonorum omnium, qui omnia nobis in usum bonum contulit si non ad libidinem et superfluitatem utamur, et contra dei mam datum. A principio enim herbas conoessit, post dilu uium carnes addidit. Sanguinem negauit, postea animalia immunda prohibuit, per Christum autem omnia mundauit. Et regnum dei dixit, non esam et potum, sed fidem et charitatem etc.  \pend
\phantomsection
\addcontentsline{toc}{subsection}{\textit{Fidelibus et his qui cognouerunt ueritatem: quia omnis creatura dei bona est. Et nihil reiiciendum, quod cum gracia rum actione percipitur. Sanctificatur enim per uerbum dei et orationem. }}
\subsection*{\textit{Fidelibus et his qui cognouerunt ueritatem: quia omnis creatura dei bona est. Et nihil reiiciendum, quod cum gracia rum actione percipitur. Sanctificatur enim per uerbum dei et orationem. }}\pstart Contra Iudaeos et Iudaizantes non solum loquitur, quorum iam haeresis ut noti damnabatur, sed contra futuros praedicatores qui docturi essent ali quos cibos esse malos. Et quod religi sius esset hoc cibo uesci quam illo, quod deus placaretur hoc, non illo cibo sumpto offenderetur mortaliter, hoc tempo re comedente hunc cibum, alio tempore alium, introducentes uota contra huiusmodi escas tanquam deo grata, quasi placeret deo magis hic cibus quam ali us: licet autem unus cibus magis impediat ingenium, sa cras meditationes quam alius, nullus tamen hoc facit ni si immodice sumptus, quod iure dei uetatur. Sunt  \pend
\section*{COMMEN. IN PRIOREM }
\marginpar{[ p.100 ]}\pstart igitur omnes tales doctrinae falsae contra deum adin uentiones humanae et daemoniorum. Omnis autem cibus semper et omnibus bonus est, si cum gratiarum actione ad naturae sustentationem sumatur.Omnis autem aibus semper malus, si nimius et ad libidinem sumatur, eatenus ut uotum illud uel aliud si cut mendadum humanum, in his rebus assumi non debet. Si cibus superfluus offeratur, semper illicitus est. Non est ut abstinentia uoueatur, si ad necessitatem utilis et moderatus, semper assumatur, nec deo placere ut necessariis naturae detrahatur. Obligamur omnes et semper ad moderatum usum cibi et omnium bono rum temporalium. Non oportet uouere uel abnegare quod dei uoluntati resistit. Cibus nunquam sumendus sine graciarum actione. Christus hymnum et laudes dixit, post cibum et coenam exactam. Siquid immundum uel insanum in cibo sobrie sumpto, oratione et uerbo dei audito et credito, sanctificatur, et fit remediabile et innoxium. Scienter tamen ue sanum cibum et sine necessitate acpere, peauto non caret. Non est deus temptandus, non omnia omnibus conueniunt.  \pend\pstart Nota non omnibus est creatura omnis bona, non omnium cibus sanctificatur per uerbum dei et orationem, sed fidelibus tantum, id est his qui co  \pend
\section*{TIMOTHEI EPISTOLAM. }
\marginpar{[ p.101 ]}\pstart gnouerunt ueritatem. Ignorans enmi omnem crea turam esse bonam, et in conscientia sua iudicans malam et non abstinens pecut, quia contra fidem agit. Nemo igitur contra conscientiam instigandus ad abum, sed prius instruendus est, ut fidelis ueritatem sa at. Imo et talium scandalo cauendum, si ex infirmitite, non contempta et temeritate credere non pos sint. Si enim esca scandalizat fratrem meum non manducabo etc.  \pend
\phantomsection
\addcontentsline{toc}{subsection}{\textit{Haec praeponens fratribus, bonus eris minister Iesu Christi, enutritus in uerbis fidei et bonae doctrinae quam assecutus es. Erasm. De his, si commonefecris fratres, bonus eris minister Iesu Christi, enutritus in sermonibus fidei, bonxque doctrinae, quam usque secutus es. }}
\subsection*{\textit{Haec praeponens fratribus, bonus eris minister Iesu Christi, enutritus in uerbis fidei et bonae doctrinae quam assecutus es. Erasm. De his, si commonefecris fratres, bonus eris minister Iesu Christi, enutritus in sermonibus fidei, bonxque doctrinae, quam usque secutus es. }}\pstart Bonus minister est, qui bene et bona praedicat, qui nutritus est in uerbis fidei, et inuerbis bonae doctrinae. Verba autem fidei sunt Euangelium, id est quicquid in tota sacra scriptura et canonica continetur, quae nos doct fidere et sperare et amare deum, quae Christum praedicat, uel futurum exhibet, in quem credere iusticia fidei sit.  \pend\pstart Bona doctrina est, quicquid, ex sacris scripturis operandum dooetur, uel in eius intellgentiam dirige  \pend
\section*{COMMEN. IN PRIOREM }
\marginpar{[ p.102 ]}\pstart re potest, et non illae litterae gentilium, qui abusiuo no mine bonae litterae nominantur, quae bonum non red dunt, nec ad bonitatem instituunt. Haec autem de cibis non diiudicandis et generaliter bonis doctrina diatur fratribus proponenda, fidelibus salioet, ut pri mo fratres sunt, fideles sint, credant in Christum legis perfectionem, et poterunt erudiri de aessatione legalium. Non enim Iudeis oeessant legalia seruanda, nisi credant in Christum, alioqui stante infidelitate magis pec cbunt, si legem suam non seruent, contra conscientiam agentes. Sunt autem primo docndi de Christo, postra de gracia et fide: sic tandem de libertate Christi na, sic et superstitiosi docendi. Bonus minister, qui bona bene praedicat ad utilitatem et fructum fidei. Malus igi tur minister erit, qui non est assecutus bonam doctrin am qui non est enutritus in uerbis fidei, quantumlibet alio qui multa didioerit, qui nonChristi Iesu, sed sua gloriam et utilitatem sectatur: qui quaerit quae sua sunt, non quae Christi, qui dooet non nubendum, qui suadet uouendam perpetuam castitatem, quod in uerbis fidei nusqua continetur, qui delectum doct in abo, tempore, et caeteris ceremoniis non a Christo institutis, uel qui uerba fidei non doct qui proiicit margaritas ante porcos scienter etc hi mali ministri non sunt audiendi, sed instruen deuel eiiciendi.  \pend
\section*{TIMOTHEI EPISTOLAM. }
\marginpar{[ p.103 ]}\pstart Quam assequutus es y Quam a me didicist, non nouum aliquid comminsoens, non neuas quaest ones inquirens, non noua mendata proponens, sed quae a me didicisti, uerba fidei, haec diligenter fratribus in culca Sic nec ipse Paulus aulus erat aliquid dogre, quod non in se fecit, uel docuisset uel reuelasset Chri stus Nostri autem non sic agunt, breuiter euangelium recitant, deinde sua ad longum ingerunt, nisi quod nunc dei gracia talcs non audiuntur a fratribus, quia semel gustauerunt doctrinam uerborum fidei.  \pend
\phantomsection
\addcontentsline{toc}{subsection}{\textit{Ineptas autem et anileis fabulas deuita: exerce autem te ipsum ad pietatem. Nam corporalis exercitatio ad modicum utilis est: pietas autem ad omnia utilis est, promissionem habens uitae quae nunc est et futurae. }}
\subsection*{\textit{Ineptas autem et anileis fabulas deuita: exerce autem te ipsum ad pietatem. Nam corporalis exercitatio ad modicum utilis est: pietas autem ad omnia utilis est, promissionem habens uitae quae nunc est et futurae. }}\pstart Didicerat a Iudaea matre sua Timotheus multas fabulas, quas ualde Iudaei admirantur unice, et nunc Cabula dici uolunt aliqua, et pueris ac puellis legenda permittunt, sicut et nos legendas sanctorum et apocrypha. Vult Paulus nihil horum praedicari, sed tantum uerba fidei et sacras scriptu ras, et imbui oertissimis ucritatil us non fabulis Non sunt pueris fabulae poëtarum proponendae, sed fidei  \pend
\section*{COMMEN. IN PRIOREM }
\marginpar{[ p.104 ]}\pstart fomenta, et morum, multo minus in ecclesia aliquid praedicandum, nisi canonicum ac probatum. Pietas qui dem est cultus dei. Is uero deum uere colit, qui fi de et charitate fertur totus in deum, et ex eadem fi de et charitate propensissimus fertur in omnia cha ritatis officia proximo impendenda. Is deum colit, qui sua mandata facit, qui obedit suis iussis. Non qui saepe et multum uersatur in humanis quibuscunque  traditionibus, sed qui deum, et proximuum propter deum diligit. Haec est uera pietas, hic uerus et legittimus dei cultus, haec religio Christiana. In hisse exeroere debet episcopus, non in cantu, oeremoniis, uestibus, ornatibus, et caeter. plaoere, sed in pieta te. Hic salioet plenus fide et charitate litterarum studio et peritia instruat populum, sit sollicitus ne oues errant in fide in moribus placidis. Et ita inhortetur ut nemo inuidia fouens laedat charitatem in ecclesia, nemo alteri noceat, uer bo aut facto, nemo egeat necessario uictu et amictu. Haec sunt pietatis officia, istae oblationes deo placetes, his sacrificiis promeretur deus, hic cultus dei et uita Christiana, et digna, imo necessaria episcopo.  \pend\pstart Corporalis exerdtatio ) Sicut cor poralis castitas et ui rginit as, tetunium, uigiliae, peregrinationes, paululum et tin casu prodesse possunt, sed et obesse saepe sic et de multis aliis Pictas autem de qua diximus gra  \pend
\section*{TIMOTHEI EPISTOLAM. }
\marginpar{[ p.105 ]}\pstart ta est deo et hominibus: prodest nobis et proximo non est inauis et infructuosa. Imo quicquid non in gloriam dei et proximi utilitatem referri potest, in utile est, et corporalis exercitatio non spiritualis dicitur.  \pend\pstart Promissionem ) Promissionem huius uitae habet, quia eadem mensura, qua mensi fuerimus etc. De operi bus minime in extremo iuditio praemiabimur. Fides in Christo hic nos beat et in futuro. Cha ritas proximi, est impletio omnis legis, quae omnino non praeiudicare debet charitati, nec exercitatio cor poralis quaecunque , nec uotum, nec promissum. Cultui enim diuino quid praeferri potest? Obedientia autem dei, nónne cultus dei?  \pend
\phantomsection
\addcontentsline{toc}{subsection}{\textit{Fidelis sermo, et omni acceptione dignus. In hoc enim laboramus et male dicimur, quia speramus in deo uiuo, qui est saluator omnium hominum, maxime fidelium. }}
\subsection*{\textit{Fidelis sermo, et omni acceptione dignus. In hoc enim laboramus et male dicimur, quia speramus in deo uiuo, qui est saluator omnium hominum, maxime fidelium. }}
\phantomsection
\addcontentsline{toc}{subsection}{\textit{Erasmus: Indubitatus sermo dignusque ́ qui omnibus modis approbetur. Nam in hoc et laboramus, et probris afficimur, quod spem fixam habemus in deo uiuente, qui est seruator omnium homi num maxime fidelium. }}
\subsection*{\textit{Erasmus: Indubitatus sermo dignusque ́ qui omnibus modis approbetur. Nam in hoc et laboramus, et probris afficimur, quod spem fixam habemus in deo uiuente, qui est seruator omnium homi num maxime fidelium. }}
\section*{COMMEN. IN PRIOREM }
\marginpar{[ p.106 ]}\pstart Attentum uult discipulum et episopum ut sibi cer to credat et doctrinam, ut uerissimam amplectaiur, non diffi dat siquos uideat alia docentes et agentes, uerosque apostolos, et uerum euangelium persequentes fa£ ctis et uerbis, licet enim ad horam periclitari uidean tur propter ueritatem, quia tamen sperant, et omnem fiduciam collocauerunt in deum, omnia moderantem et omneis homines inse credentes sa uantem, qui deus uiuus est, non idolum, non ignarus, non inuidus, non tyramnus. Ideo et nos qui irfideles sumus sine dubio suo nondestituemur auxili, nec ueritas per nos praedi cata aliquando succumbet, sed triumphabit, et nos cum illo in gloriam dei qui omnium saluator est, maxime fidelium.  \pend\pstart Fidelis sermo I Peculiare uerbum apostolo aliquid asseuerare uolenti ne ab aliquo uel hesitetur uel uel negligatur ut inutile uel casu dictum. Vult autem intelligi euangelium quod praedicuerat Ephesiis et edocuerat Timotheum propter quod sine intermissi one perpetiebatur multa.  \pend\pstart In hoc enim et laboramus, et maledicimur, Propter hoc euangelium multa passi sumus, et ad huc patimur, tanqua falsi doctores, subuersores ueri tatis, speudoapostoli et apostate a lege paternarum traditionum. Et mala multa de nobis loquuntur et im precantur mala.in tantumut etiam nobis id uitio detur  \pend
\section*{TIMOTHEI EPISTOLAM. }
\marginpar{[ p.107 ]}\pstart et crimini, ꝗ ausi simus tales deceptores et falsi do ctores et subuersores uerborum dei adhuc sperare in deo uiuo, cuius ipsi nos blasphematores uideamur. Ipse dominus deus noster, cui seruimus, cui obedimus in praedicatione euangelii, in quem confidimus, cui nos totos credimus, ipse seruator est oïm hoïm etc.  \pend\pstart Saluator est om. I Pro omnibus sollicitus est, omnes suas creaturas diligit, suum filium ex amore maxi mo pro omnibus tradidit, solem suumsuper omnes diffundit, salutis consilia et remedia nemini negat, ci£ tra condignum punit, omnes uult saluari, uiam sa lutis aperuit omnibus.  \pend\pstart Maxime fidelium I Toto corde sibi credentium fiduciam omnem collocantium in eum, qui uerbis suis credunt, qui praeoeptis acquiescunt, uoluntatem suam faciunt. Non facientes poenitent, non poenitentes dolent se non poenitere. Orant pro peautis in no mine Christi, iusticia Christi et meritis confidunt, de se sperant. Hi tales fideles sunt et eorum certissimus sal uator deus est, qui per Christum saluat fideles tanquam filios, amicos et intimos quos dimittere non potest.  \pend
\phantomsection
\addcontentsline{toc}{subsection}{\textit{Praecipe haec et doce. Nemo adolescentiam tuam contemnat: sed esto exemplum fidelium in omni uerbo et conuersatione, in charitate, in fide, in castitate. }}
\subsection*{\textit{Praecipe haec et doce. Nemo adolescentiam tuam contemnat: sed esto exemplum fidelium in omni uerbo et conuersatione, in charitate, in fide, in castitate. }}
\section*{COMMEN. IN PRIOREM }
\marginpar{[ p.108 ]}
\phantomsection
\addcontentsline{toc}{subsection}{\textit{Erasmus: Praecipe haec et doce. Nemo tuam iuuentutem despiciat, sed esto forma fidelium, in sermone, in conuersatiotione, in dilectione, in spiritu, in fide, in puritate. }}
\subsection*{\textit{Erasmus: Praecipe haec et doce. Nemo tuam iuuentutem despiciat, sed esto forma fidelium, in sermone, in conuersatiotione, in dilectione, in spiritu, in fide, in puritate. }}\pstart Haec praecipe et doc Confidentiam in deum in omni labore et maledictione, pietatem ueram seruandam prae omnibus. Nihil praeter scripturam docendum, fabulis non attendendum, matrimonium non interdi cendum, cibos non respuendos tanquam animae damno sos etc. Similiter quae de Christo docenda praescripsi breuibus, haec inquam praecipe et doce, non humanas doctrinas uel traditiones.  \pend\pstart Nemo adolescentiam tuam contemnat ) Actas senectutis, uita immaculata, deus mentes magis com probat, quam aetates, non tanquam inexpertus spiri tualium loquaris. Non adolescentium more agas et doceas, licet adolescens sis, factis ostende et doctrina et moribus ut adolescentia tua promereatur ho norem senectutis: quod adolescentia detrabit autori tati, hoc morum et doctrinae maturitate aoedet. Ideo per autoritatem praecipias et doccas, ut non huma na, sed ex spiritu proueniens doctrina credatur, ut non ob adolescentiam doctrina spernatur, sed ob maturitatem et soliditatem doctrinae ueneretur ado  \pend
\section*{TIMOTHEI EPISTOLAM. }
\marginpar{[ p.109 ]}\pstart lescentia, et a nemine sperni possit.  \pend\pstart Sed esto exemplum fidelium ) Hoc est tandem uiiliter praeapere et docere, si prius feceris, si exem plar honestatis, in uerbo et conuersatione praestiteris, si comprobaueris uerba factis, si uita non doctri nam condemnet, sed perjuadeat, non uerba sufficiunt coepit facere et docere princeps pastorum Christus. Exemplar autem sit primo in sermone, ut nihil nisi scripturas sacras doceat, et solidam doctrinam, sem per de honestis loquatur, obloquatur nemini. Omnia interpretetur mitius, increpet detractatores. Sit apud eum est, est: non, non. Secundo in conuersatione, ut humiliter, mansuete, pudice, sobrie, amanter, ad omnes se exhibeat, paratus omnium commodis promouendis. Sollicitus de necessitate proximi magis quam propria corporali et spirituali. Tertio in dilectione dei et proximi. Propter deum in corde, non solum et opere prosit ex animo, non hypocritice, non ob commodum priuatum, etiam inimicos, externos, immeritos, ingratos. Quarto in fide, ut totus ex deo pendeat, se deo totum committat, omnia quae eueniunt deo acepta firat, non dubitet de gracia dei, Chrisui beneficia credat, et iugiter meditetur, dei gloriam quaerat in omnibus, uerbis det assistat, credat discat, doceat ueritatem fdei  \pend
\section*{COMMEN. IN PRIOREM }
\marginpar{[ p.110 ]}\pstart et scripture, non feratinfideles et caetera. Quinto in puritate, fides hoc cor mundum creat. Verba spurca non patitur, opera turpia exhorrescit, nihil mixtum, uarium, impurum, duplex, admittit, sed omnia sunt simplicia dextra, certa, pura, et sincera quae agit et loquitur. Castitatem eatenus seruet, ut libido regatur et non regat.  \pend
\phantomsection
\addcontentsline{toc}{subsection}{\textit{Donec uenero, attende lectioni, exhortationi, et doctrinae. }}
\subsection*{\textit{Donec uenero, attende lectioni, exhortationi, et doctrinae. }}\pstart Non praestoleris ocosus aduentum meum, noli mihi reseruare, quae tu poteris praestare. Lectioni interim instesscripturarum sacrarum, unde possis tibi et sub ditis proficere. Non erunt in manibus tuis nisi sani ꝙ dioes, non philosophia, non quicquid mundani docent, sed lex et prophetae et psalmi, idest canon scripture ucteris. Euangelia apostolorum traditionibus disceban tur donec describerentur et euulgarentur per discipu los, haec legebantur non fabulae gentilium uel Iudaeorum.  \pend\pstart Exhortationi) A d fidem Christi et mores aposto licos, non cogendo, sed hortando instituas imperfe ctos, uicia infrmitatis dulabus uerbis perfuade, non manda, non mineris, non increpas, sed exhorteris suauiter et instanter.  \pend\pstart Et doctrinae) Ex scripturis proba, ostende, cuince, doceas fic se habere omnia iuxta dei legem, et  \pend
\section*{TIMOTHEI EPISTOLAM. }
\marginpar{[ p.111 ]}\pstart sacras literas, ut tu praedicas, ut non sint inoerti, et opi nantes, sed fide scripturarum solidentur in doctrinis sanis apostolicis, quales a me suscepisti.  \pend
\phantomsection
\addcontentsline{toc}{subsection}{\textit{Ne neglexeris quod in te est donum, quod datum est tibi per prophetiam cum impositione manuum, autoritate sacerdotii. }}
\subsection*{\textit{Ne neglexeris quod in te est donum, quod datum est tibi per prophetiam cum impositione manuum, autoritate sacerdotii. }}\pstart Donum prophetiae hoc est in terpretandi sacras literas affluentius ci datum erat, quam caeteris, ideo episcopus ordinatus est, cuius est docere legem dei, nihilo ninus tamen lectioni instare iubetur nedesit lo corum collatio, ut certior interpraetatio habeatur. Non enim quilibet passim concinata enarratio sufficit sed quae desumitur ex scripturis, nisi ubi manifesta sit le ctio Hanc propheciae gratiam uide ne negligas, ne talentum creditum su odias, sed usu et exercitio quotidiano usuram comportes domino tuo debitam et red dendam et exigendam.  \pend\pstart Cum tibi presbyteri et seniores imponerent ma nus, ut esses epi scopus per prophetiam habebant de te, quod idoneus esses huic officio. Et quod gratiam ac donum spiritus suscepi Jes docendi ecclesiam et donum prophetiae. Bo utere solliate et ne negligas, sed do ce exhortare, et fac iuxta acceptam a deo gratiam per impositionem manuum, quoniam certum est,  \pend
\section*{COMMEN. IN PRIOREM }
\marginpar{[ p.112 ]}\pstart adesse tibi haud dubium quin in profectum ecclesiae tuae ne sis ut seruus ne quam domino, et ecclesiae infidelis.  \pend
\phantomsection
\addcontentsline{toc}{subsection}{\textit{Haec meditare, in his esto, ut profectus tuus manifestus sit omnibus. Attende tibi et doctrinae. Insta in illis. Hoc enim faciens, et te ipsum saluum facies et eos qui te audiunt. }}
\subsection*{\textit{Haec meditare, in his esto, ut profectus tuus manifestus sit omnibus. Attende tibi et doctrinae. Insta in illis. Hoc enim faciens, et te ipsum saluum facies et eos qui te audiunt. }}\pstart Admonui de lege domini praedicanda, non fabulis et Christianae fidei articulos, instare lectioni, exhortationi, ad fidem, spiritum, dilectionem, puritatem conuersationem et sermonem. Haec inquam medita re, ut apposite, utiliter, et sufficienter de eis loqui possis, non temere, sine praemeditatione Nondum commen daui sacrificia, missas, baptismos et ceremonias, de quibus non est tanta necessitas. In his quae tibi proposui esto, studium omne confer, jet sollicite cura ut fi£ ant, et ita fiant, ut pro ea proficias patenter in omni bus tibi subditis, ut prosit ouibus pa toralis cura tua ut sciant ex commodo suo spirituali se habere episoopum, doctorem et curatorem. Ne semper in eodem maneat populus, sed de bono proficiat in melius, in fi£ de et conuer satione Christiana crescat indies.  \pend\pstart Attende deniq interim et tibi, et doctrinae, ut cul pa non redarguat doctorem, non contrauenias exem plo, eis quae docueris, et ut doctrina tua sana sit, et ut  \pend
\section*{TIMOTHEI EPISTOLAM. }
\marginpar{[ p.113 ]}\pstart praeter talem nemo subintret, alius aduersa docens sed et ceu firma et certa teneatur et credatur ab omnibus, ut proficiant non deficiant.  \pend\pstart Insta in illis, ut et docendo et agendo, tu et po pulus proficiat: illa inqua non alia propone, inculca, defende, quae si pro uocatione et gratia dei feceris, et officio tuo incubueris etc.  \pend\pstart Teipsum saluum facies, euades iuditium nequa et infidelis serui, et pastorum negligentium, et sua quaerentium, gregem deserentium et perdentium, non audics grauem comminationem et dannationem, te omnium eximes periculo, etiam si non acquieuerint, etiam si contemnent et reiiciant, non credant et damnant. Tu tamen liberabis ani mamet conscientiam tuam, principi pastorum obediendo.  \pend\pstart Si uero audierint, et obaudierint, si monitis et doctrinae tuae dociles animum apposuerint et pro fecerint in ea, etiam tu eis causa salutis eris, ut una tecum saluentur, et salutis suae fidem ministrum grex grata agnoscet coram deo.  \pend
\endnumbering\beginnumbering\section{CAPVT V.}\pstart SEQVITVR licet exemplaria quaedam adiecerint quarto capiti quae sequuntur, usque ad Viduas honorandas.  \pend
\section*{COMMEN. IN PRIOREM }
\marginpar{[ p.114 ]}
\phantomsection
\addcontentsline{toc}{subsection}{\textit{\huge\textbf{S}\normalsize Eniorem ne increpaueris, sed obsecra ut patrem, iuuenes ut fratres nus ut matres, iuuenculas ut sorores in omni castitate. }}
\subsection*{\textit{\huge\textbf{S}\normalsize Eniorem ne increpaueris, sed obsecra ut patrem, iuuenes ut fratres nus ut matres, iuuenculas ut sorores in omni castitate. }}
\phantomsection
\addcontentsline{toc}{subsection}{\textit{Erasmus: Seniorem ne saeuius obiurges sed adhortare ut patrem, iuniores ut fratres, mulieres natu grandiores ut matres, iuniores ut sorores cum omni castitate. }}
\subsection*{\textit{Erasmus: Seniorem ne saeuius obiurges sed adhortare ut patrem, iuniores ut fratres, mulieres natu grandiores ut matres, iuniores ut sorores cum omni castitate. }}\pstart Modestia et lenitas iuncta autoritate, necessaria sunt in episcopo, et ut neminem non obseruet, omnes curet, omnibus prospiciat et medeatur. Non negligat etiam si maiores sint aetate, si quid contra ea quae decuerit ab eis fieri et committi uideat, hoc est contra legem dei, id autem faciat comitate et pia admonitione, non seuera correptione.  \pend\pstart Seniores non facile ferunt a iuniore mo neri, multo minus increpari seuerius, et tamen eorum insolentia neutiquam dissimulari debet, ne ob senium in autoritatem et exemplum trahatur a iunioribus procliuioribus ad mala sectanda. Corrigi ergo debent, sed non saeuitia, non duricia uerborum, sed blandiciis, lenitate et quasi reuerenti insinuatione admonendo potius contumeliis irrogandis, quae ferri non possint a seni bus, sed potius irritantur mags exasperati. Seruan  \pend
\section*{TIMOTHEI EPISTOLAM. }
\marginpar{[ p.115 ]}\pstart da igtur reuerentia aetatis, sed tamen non intermitten da admonitio transgressionis, sed leniter et obsecra tione, non imperio uel increpatione. Si uero senes omnino stulti sint, et inueterati dierum, publico scan dalo patrocinantes autoritate senectutis suae, Danielis exemplo utendum est, et pro rei indignitate tractandi. Semper tamen decet, iuuenem praelatum mo destia in uerbis, de quibus hic agituy ad senes, quatinus magis condolentia et dolor malitiae et rigoris neoessitas urgere uideatur, quam contemptus senectutis et praesumptio praelationis, quam iuuenes declinent.  \pend\pstart Iuuenes ut fratres) Solemus uariis modis allo qui et corrigere fratres insolentes, iuxta cuiusque indolem ut unum alloquemur durius, alteri familiarius, et affectu dulcori conuenimus, aliquando et ir ritamus uolentes turbari, ut citius resipisct. Nunquam tamen infamia publica eos inuri patimur, praeue nimus incommodum famae et corporis, oontra ex ternos acusatores pro uirili excusamus, defendimus et uindicamus familiariter duri, si expediat: publice fideles et propugnatores. Sic concliamus fratrum animos, sic emendamus, sic proficimus, omnia ex sinoe ra, naturalique ́; fraternitate consulentes et adaptantes. Sic ad Christianos omnes et coëtaneos affid debe\pend
\section*{COMMEN. IN PRIOREM }
\marginpar{[ p.116 ]}\pstart mus et agere ut consulamus ut facile corrigantur, si amorum tantum causa admo neantur uel etiam corripiantur  \pend\pstart Anus ut matres ) Honorifice alloquendo, commouendo, dedecus suum ostendendo. Erubescentiam uel ruborem iniciendo, ex rei indignitate proposita non contemptu, odio, maleuolentia uerbis mode stis, ne irritentur potius quam emendentur, et ad desperationem adigantur et obruantur tristitia, bis contemptae.  \pend\pstart Iuuemculas ut sorores, in omni castitate ) Amanter alloqui consueuimus sorores, et sincere, sincque formae et libidinis habita ratione, sed tantum ut consulamus honori et rebus agendis. Sic dum Episcopus parrochus necessitate adigitur ad alloquendum iuuenculas, id breuiter faciat sincere grauiter et cau te, omnia leuitatis signa cauens. Verecunde et com posite ad disciplinam se habens circa eam, ne aliquo peoato suspitionis sinistre iudicium proferat, et ca stitatis lubrica fama periclitetur, etiam si nihil mali intercedat. Ab omni enim specie mala abstinendum maxime episcopo. Lubrica autem est fama pudoris inter iuuenculas praecipue et adolescentemepisco pum, quantumlibet grauis sit moribus. Nemo enim satis, et omnibus horis sapit hac in causa alioqui lapiens.  \pend
\section*{TIMOTHEI EPISTOLAM. }
\marginpar{[ p.117 ]}
\endnumbering\beginnumbering\section{CAPVT V.}
\phantomsection
\addcontentsline{toc}{subsection}{\textit{\huge\textbf{V}\normalsize Iduas honora, quae uere uiduae sunt. Si quae autem uidua filios habet aut nepotes habet, discat primum domum suam regere et mutuam uicem reddere parentibus, hoc enim acceptum est coram deo. }}
\subsection*{\textit{\huge\textbf{V}\normalsize Iduas honora, quae uere uiduae sunt. Si quae autem uidua filios habet aut nepotes habet, discat primum domum suam regere et mutuam uicem reddere parentibus, hoc enim acceptum est coram deo. }}
\phantomsection
\addcontentsline{toc}{subsection}{\textit{Erasmus: Viduas honora, quae uere ui duae sunt. Quod si qua uidua liberos aut nepotes habet, discant primum propriam domum pie tractare, et uicem rependere maioribus. Hoc enim est honestum et acceptum coram deo. }}
\subsection*{\textit{Erasmus: Viduas honora, quae uere ui duae sunt. Quod si qua uidua liberos aut nepotes habet, discant primum propriam domum pie tractare, et uicem rependere maioribus. Hoc enim est honestum et acceptum coram deo. }}\pstart Ex consequentibus contextuque ; sermonis intelligitur, is honor deferendus uiduis, ut praecaucat ne uere uiduae, quae sunt marito, filiis, nepotibus et di uitiis destitutae pariter sunt, id est neminem habent sui curam habentem, istis uiduis honorem habeat episcopus eatenus, ut non paciatur talem contemni, male auditam indigere, odiri, turbari, molestari. Ipse episcopus siuere sit uidua ecclesiae nomine talem foueat, defendat, iuuet, consoletur, et de sumptu eccle siae et oblationibus collectis in Synaxi fidellum prouideat talibus uiduis in communi de necessariis. Et id intelligendum de omnibus eadem miseria labo\pend
\section*{COMMEN. IN PRIOREM }
\marginpar{[ p.118 ]}\pstart rantibus pupillis, orphanis, infirmis, peregrinis, elerias, et qui sua pauperibus erogarunt et postra egent necessariis. In summa caueat episcopus, ut nemo inter Christianos sibi subditos indigeat et mendioet Honora id est non patiaris miscre, despectique eam uiuere  \pend\pstart Si qua autemn ) Si aliqua sit non uere uidua, solum quidem marito uiduata, supersunt autem ei filii et nepotes, istam non curabit necessario episcopus, sed filii et nepotes, si indigeat, uictum minisirabunt et amictum. Ipja uicissim eos regat, instruat, moneat, et si potest, nutriat et foueat, et eduoet. Similiter eisdem obsequium humanitatis praestans, quale ipsa a suis susaepit parentibus. Debent enim patres thesaurizare, filii non coutra Nec susaipiatur ad collegium uere uiduarum, sed domui et familiae suae potius prae sit et incumbat Hoc enim placet deo, ut prouectiores aetate et prudentia iunioribus praesint, et qui aliis benefacere possunt, non patiantur ocosi sibi benefi eri et offitium capere.  \pend\pstart Disat primum domum suam regere ) Non intromittat se in ministerio ecclesiae deseruiendo dimissa familia filiorum, nepotumque suorum: istis potius praesit, et benefaciat, et regat domum propriam Quam ubi bene rexerit, et regendo exercitatior fuerit, postea ecclesiae descruire legttime poterit si  \pend
\section*{TIMOTHEI EPISTOLAM. }
\marginpar{[ p.119 ]}\pstart haec ad hoc per episcopum sexagenaria delegatur.  \pend
\phantomsection
\addcontentsline{toc}{subsection}{\textit{Quae autem uerem uidua est, et desolata, speret in deum, et instet obsecrationibus et orationibus, noctu dieque ́. Nam quae in delitiis est, uiuens mortua est. }}
\subsection*{\textit{Quae autem uerem uidua est, et desolata, speret in deum, et instet obsecrationibus et orationibus, noctu dieque ́. Nam quae in delitiis est, uiuens mortua est. }}
\phantomsection
\addcontentsline{toc}{subsection}{\textit{Erasmus: Porron quae uere uidua est ac desolata, sperat in deo et perseuerat in obsecrationibus ac precationibus noctu dieque ́. Porron quae in delitiis uersatur, ea uiuens mortua est. }}
\subsection*{\textit{Erasmus: Porron quae uere uidua est ac desolata, sperat in deo et perseuerat in obsecrationibus ac precationibus noctu dieque ́. Porron quae in delitiis uersatur, ea uiuens mortua est. }}\pstart Describitur hic uere uidua, quae salicet desolata estomni solatio huius mundi, caret marito, filiis, nepotibus, possessione: sperat autem in domino, et omnem spei suae fiduciam collocat in deum, qui pupillum et uiduam susapiet. Non utitur hoc mundo, non sectatur gaudia huius mundi, libidines carnis, sed continuo uacat orationibus, et audientiae uerbi dei, seruitio indigentium, consolatione afflictorum, sine intermissione curat quae domini sunt, exo nerata molestiis huius uitae et curae domesticae, et quo modo placeat marito, sed deo uacat soli, Talis uere uidua est, id est, desolata apud mundum, deo autem coniunctissima, die, noctuque sicut de Anna legimus, quae Christum suscepit in templo oblatum. Vidua uero lla quae in delitiis uersa,  \pend
\section*{COMMEN. IN PRIOREM }
\marginpar{[ p.120 ]}\pstart tur, non seruiens deo, nec uacans orationibus, sed abutens libertate, delitias sectatur ocio, multiloquio, uinolentiae, libertate in occasionem carnis tradita uacat. Talis ut maxime suauiter iuxta tarnem ui uat, sed mortua est spiritu et gratia dei priuata. Christo non uiuit, quo d pietati non uiuit. Tales ab ccelesia nutriri mini me debent, cum sint opprobrium et diffamia ecclesiae. Sed corrigi potius et puniri, si quo modo reuiuiscat in Christo per poenitentiam.  \pend
\phantomsection
\addcontentsline{toc}{subsection}{\textit{Et haec praecipe, ut irreprehensibiles sint. }}
\subsection*{\textit{Et haec praecipe, ut irreprehensibiles sint. }}\pstart Etiam hic hoc pertinet ad tuum offitium, ut doceas eas esse tantae maturitatis et circumspectionis, ut nihil de eis sinistri suspicari et dici possit, ne laxius uiuant ad libertatem suam in uiduitate, et scan dalo sint ecclesiae, etuideatur quod non ob religionem et castitatis amore abstincant ab iterando coniugio, sed ad satisfacien dum libidini sine timore mari ti, quo carent.  \pend
\phantomsection
\addcontentsline{toc}{subsection}{\textit{Si quis autem suorum et maxime domesticorum curam non habet, fidem negauit, et est infideli deterior. }}
\subsection*{\textit{Si quis autem suorum et maxime domesticorum curam non habet, fidem negauit, et est infideli deterior. }}
\phantomsection
\addcontentsline{toc}{subsection}{\textit{Erasmus: Quod si qua suis et maxime familiaribus non prouidet, fidem abne}}
\subsection*{\textit{Erasmus: Quod si qua suis et maxime familiaribus non prouidet, fidem abne}}
\section*{TIMOTHEI EPISTOLAM. }
\marginpar{[ p.121 ]}
\phantomsection
\addcontentsline{toc}{subsection}{\textit{gauit, et est infideli deterior. }}
\subsection*{\textit{gauit, et est infideli deterior. }}\pstart Licet haec sententia interiecta, uidedlur dici de uidua, quae non est uere uidua, quae filios et nepotes habet, et curare iubetur familta, potius quam ut uel ecclesiae iungatur alenda uel eidem ministratura, sed potius suae domui praesit et inslituat ac re gat, sicut reuera hoc de uidua Paulus retulit. Verumtamen generaliter eatenus hoc pronunciauit, ut uideatur uoluisse id uniuersaliter, et de uiris et mulieribus uiduis et maritis esse suscipiendum, et de omnibus communiter, ut neminem sic uelit curare aliena, ut propria familiae suae necessaria, et domesticorum quibus dei prouidentia, rector et prouisor constituitur curam reiiciat, et ea tamen ra tione orationi uacare eligat, ut non cogatur proximi necessitati seruire.  \pend\pstart Hic iubetur paterfamilias et materfamilias dome sticorum curam diligenter gerere, non secus acsi a do mino spiritualiter eorum curam demandatum aoepisset Quod si abiiciat huius munus a deo impositum curam dae familiae, instituendae prolis, regendae domus, mo e stias sustinendi pro proximo, et s ecundum carnem de quibus deus mandauit, et uelit obmissis istis uacare ppriae deuotioni ecclesiae ministrare, religosius ad speciem fieri et deo sccundum proprium sensum seruire  \pend
\section*{COMMEN. IN PRIOREM }
\marginpar{[ p.122 ]}\pstart hic tantum abest ut deo placeat, ut magis infideli sit deterior, qui curat familiam quod iussit, natura, et deus: hic autem quod sibi placitum est eligit, et quae mandauit deus, obmittit, et deum mendacem et mi rabilem facit, aliud uelle, aliudque iubere credit quod est contra fidem. Infideles enim naturae ductu prospiciunt proli et proximis, et id deus praecepit. Id hii abundantes suo sensu despiciunt, et propter tra ditiones proprias humanas, transgrediuntur mandatum dei.  \pend
\phantomsection
\addcontentsline{toc}{subsection}{\textit{Vidua eligatur non minus sexaginta annorum, quae fuerit unius uiri uxor in bonis operibus testimonium habens, si filios educauit, si hospitio rece pit. Si sanctorum pedes lauit, si tribulationem patientibus subministrauit, si omne opus bonum subsecuta est. }}
\subsection*{\textit{Vidua eligatur non minus sexaginta annorum, quae fuerit unius uiri uxor in bonis operibus testimonium habens, si filios educauit, si hospitio rece pit. Si sanctorum pedes lauit, si tribulationem patientibus subministrauit, si omne opus bonum subsecuta est. }}
\phantomsection
\addcontentsline{toc}{subsection}{\textit{Erasmus: Vidua allegatur non minor annis sexaginta, quae fuerit unius uiri uxor, in operibus bonis hominum testimonio comprobata, et si filios educauit, si fuit hospitalis, si sanctorum pedes lauit, si afflictis subministrauit, si in omni opere bona fuit assidua. }}
\subsection*{\textit{Erasmus: Vidua allegatur non minor annis sexaginta, quae fuerit unius uiri uxor, in operibus bonis hominum testimonio comprobata, et si filios educauit, si fuit hospitalis, si sanctorum pedes lauit, si afflictis subministrauit, si in omni opere bona fuit assidua. }}\pstart Ecclesiastici primitiuae Ecelesiae erant omnes pau\pend
\section*{TIMOTHEI EPISTOLAM. }
\marginpar{[ p.123 ]}\pstart peres, orphani, pupilli, et imprimis uiduae, orbatae, uel desolatae uiris et nepotibus uel filiis. Et bona ec clesiae erant, quae sabbatis post auditam fidei doctrinam colligebantur, et collecta bene dicebantur, et per di aconos ad episcopi dispositionem ecclesiafticis talibus dispergebantur, de quibus et episcopus et diaconi cum sua familia, uxore scilicet et filiis sustentabantur ad suf ficientia, non abundantiam uel luxum. Peregrini etiam fra tres reficiebantur et uirgines pauperes et infirmi. Suc cessu temporis testamentis fidelium et munifica largitionc diuitum puisum est huic usui et educationi indigentium alicubi quidem decimis, alicubi possessionibus agrorum et domuum etc. Tandem deuotio principum et Christiana religio eorum terras et prouentus in usum ecclesiastico rum contulit et commendauit religiosioribus uiris, et pbatae sanctimoniae, quos et monachos appellabant, ut illi dispensarent quotamnis et semper egentibus et sacris literis incumbentibus et in spiritualibus populo seruiem tibus, orphanis, uirginibus, uiduis, hospitibus. Sunt enim usque hodie in nobilibus religiosorum monasteriis ordinationes de lectionibus sacrae scripturae, per totum annum in choro, hospitalarii, infirmarii, uirginum puiso res, scolastia qui praesunt pueris instituendis, plebani. Episcorum uero munus eratpponere ad populum, semp tamem mansit festis dieb, collecta et oblationes pro cccle  \pend
\section*{COMMEN. IN PRIOREM }
\marginpar{[ p.124 ]}\pstart siasticis, sed pro temporis longitudine et morum imita tione omnia inuersa sunt, ut bona ecclesiastica uene rint in maximum abusum quem uidemus. Episopi prina pes facti sunt, monachi satrapae, religosi caeteri plus quam mundani euaserunt. Et qui spirituales dicuntur, toti funt facti animalia uentris. Hodie nonuirgines, non ui duaesexagenariae eliguntur, sed meretrioes et iuuemculae perditissimi nebulones et lenones et adulteri nutriuntur, imo impinguantur, luxuriantur de bonis illis ecclesiae. Pauperes omnino negliguntur a clericis quicquid corrodere possunt, hoc reponunt in omnesus, unde augeatur in immensum multitudo pditissimorum scortatorum spuriorum et moretricum, omnia sub praetextu nominis et cultus diuini. Qui uero saeculares et prophani dicuntur et non ecclesiast a, hi pro ho spitibus et infirmis uiris et mulieribus, pro horren dis hominum plagis et miserandis infirmitatibus, pro pupullis et orphanis sollicitudinem gerunt, uere ecclesiastici, re et officio et cura insigni et chri stianissima etc. Viduae non omnes assumebantur ad tutelam et annonam ecclesiae, sed eligebantur ad praeoeptum apostolorum ut hic uidemus, uiduae illae tantum quales hic describuntur, non minorii aetatis quam sexagenaria: ne semel susoepta in curam ecclesiae, cum probro magno ad coniugum redeat, et ne im  \pend
\section*{TIMOTHEI EPISTOLAM. }
\marginpar{[ p.125 ]}\pstart pudiatiae suspitio inferatur, et quod ad aetatem ista ui£ ribus destitutae manuum labore se nutrire non possent Eligebantur deniq uere uiduae, non illae quae filios et nepotes regere debebant et diuites, sed destitutae.  \pend\pstart Quae fuerit unius uiri uxor ) De qua non possit esse suspitio intemperantiae in tali aetate, qualis notaretur, si multoties nupsisset, et de maritorum possessionibus ditior euadere potuisset. Laus erat non solum in ecclesia prima, sed et apud gentiles, si nuptiis secundis abstinuisset mulier. Et non parua no ta habebatur, et dedecus si repetisset maritum, hoc et apostolo placuit, ut continentur uiuant et non ite rent matrimonium si sint prouectae aetatis. Mandat autem et uult uiduas iuniores nubere quod tutius sit et ne aessarium remedium incontinentiae et periculum ma£ gnum cauetur et scandalum multiplex. Quod autem hic quasi peaatum punire uidetur secundas nuptias, cum secundo nuptam a stipe ecclesiastica reiicit, hoc agit consulere uolens honestati et fidelium famae apud infideles nolens id honoris deferre eis foeminis, quae ob impudicitiae notam a gentilibus reproba bantur, cum hic tandem ageretur de corporali subsi dio, quod talibus ab aliis poterat et modis aliis conpensari. Interimque ; consulebatur paci et tranquillita u reliquarum uiduarum ecclesiasticarum, quae non si  \pend
\section*{COMMEN. IN PRIOREM }
\marginpar{[ p.126 ]}\pstart ne murmure huiusmodi admississent suo contubernio uel ae quari eis sustinuissent. Et conscientiae interim nihil adimebatur, quam ut liberam seruaret, uel inuitis uel damnantibus gentilibus, qui sccundum mundumsapiebant, improbantes secundas nuptias in mu lieribus, non in uiris. Ipse fidelis existens Apostolus, non solum permittit, sed et iubet secundas nuptias iunioribus uiduis, ne deterius aliquid contra deum peoeetur, quod hoc indulto possit praeueniri. Aroet autem huiusmodi ab annona ccclesiastica, consulens famae, paci, et ca stimoniae non ut apud deum ream, sed infirmorum infamationi obnoxiam apud gentes, sicut et sterilitas foeminae et uxoris apud Iudaeos probro non carens. Similiter infra eadem ratione ar oetur, quod constat carere noxa apud deum et homines, sed non apud Iudaeos: quibus sicut benedictio erat ge£ nerasse filios et promissionis donum, sic caruisse maledictio reputabatur, et ut infausta et misera plange batur. Et sterilis diu mulier si peperisset tarde gaudi um omni et toti uiciniae proferebat maximum.  \pend\pstart In bonis operibus testimonium habens) Exigeba tur exercitium exactum in sanctis operibus, maximeautem charitatis et miseriae, et ut oelebraretur eius dignitas toti ecclesiae a qua sustendanda erat, ne forte fauo re et falsa laude promotam dicerent, sed conspicua  \pend
\section*{TIMOTHEI EPISTOLAM. }
\marginpar{[ p.127 ]}\pstart essent multis testibus sancta opera eius.  \pend\pstart Si filios educauit ) Non dicitur si peperit, licet id uideri possit significasse, sed si educauit, supplendum laudabiliter, religiose, honeste. Quod si utuidetur filias educasse, est simul quoque genuisse, ut sterilis aroe atur a stipe, hoc in gracia Iudaearum contigisse respondebitur. Mihi hoc probabilius uidetur, maxime cum sterilis tanta aetate molestiis puerorum carens, sibi laborando prouidere poterat de senectutis sustentamento. Similiter si filios generans, amisisset nondum eductos. Malo haec dioere quam addere uerbis Apostoli, uel laudabiliter uel fideliter, lioet impro borum filiorum uel filiarum mater etiam possit detesta bilisuisa fuisse contubernio illo sanctarum uiduarum Negari enim non potest infelicissimum esse apud Iudaeos uxoris sterilitatem.  \pend\pstart Si hospitio recipit ) Hospitalitas in ueteri et no ua lege impense laudetur, et Apostolorum tempore susaepisse praedicatores cuangelii, magnum pietatis offitium habebatur, nec fidelis aparebat, qui huiusmo di aliquando exclusisset. Nec digna ecclesiastico nutrimento, quae plantatores ecclesiae repulisset.  \pend\pstart Si sanctorum lauit pedes ) Sine calciamentis ambulantium primum hospitalitatis et facile beneficium est lauari pedes. Hoc Abraham et Loth  \pend
\section*{COMMEN. IN PRIOREM }
\marginpar{[ p.128 ]}\pstart docuerunt exemplo suo, ut fatigati quiescant suauius Exemplo quoq suo memorabili Christus commenda uit hic officium charitatis et praeoepit mutuo a disa pulis exhiberi ob humilitatis exerantium, quae charitatis fomentum esse solet. Sanctorum dicit pro praedicatione euangelii discurrentium et fidelium in cau sa fidei laborantium, uim peregrinationes et exilia sustinentium pro ucritate.  \pend\pstart Si tribulationem patientibus subministrauit ) Si enim ex euangelica doctrina omnibus petentibus tribuendum est, si praedicatoribus nihil possidendum si pro euangelio omnia dimittenda, si multa patienda, si de auitate in ciuitatem fugiendum, si non claudenda uisoera ab egeno inspecto, si in opera misericordiae profundenda sunt omnia, si nihil proprium habendum, quod non sit paratus impartiri egentiChri stianus, certe ecclesiastica disciplina spiritum sanctum ab apostolis edocta est, ut huiusmodi de paupertate non fraudentur necessariis, sed Christi praeaeptis obedientes: Christi prouidentia non egeant ne aessarii. Non enim potest id fieri, si iuxta euangelium uiuatur, ut desit alicui uitae neoessarius prouentus, et suppetentia, quibus oentuplum et in hac uita promit titur. Dicitur autem tribulationem patientibus, non oblectamenti, curiositatis, instabilitatis, auaritiae, su\pend
\section*{TIMOTHEI EPISTOLAM. }
\marginpar{[ p.129 ]}\pstart perbiae, consolationis gracia discurrentibus subministranda huiusmodi. Sed pro ueritate et euangelio patientibus, hospitalitas laudatur et numeratur. Fuerunt enim et tuncut nunc pseudoapostoli, qui domus uiduarum deglubebant uariis doctrinis, seducebant, ueros Apostolos persequebantur, commo dum et solatia propterea sectabantur, non his, sedpa tientibus tribulatione m ministranda necessaria.  \pend\pstart Si omne opus bonum subsecuta est ) Si nullum opus bonum obmisit, quod commode faoere potuit, et omnibus pietatis officiis dedita fuit. Si in omni ope re bono assidua. Non sufficit unum uel pauca egisse bona, nullum obmisisse laudabile si modo non defuis set facultas. Huiusmodi uiduas ecclesia nutrire deb t Et non solum uiduae, sed omnes etiam uiri his recita tis laudibus digni, de bonis ecclesiae debent diuino iure sustentari. Ecclesiastica bona non sunt templorum et aedificiorum et pinguium ac otiosorum clericorum sed pauperum ho iestorum et necessitati habentium. Maxine autem illorum, qui profuderunt in usus pau parum sua propter euangelium et in pietatis operi bus assidui, habuerunt temporalia sua neglectui ac cotemptui.  \pend
\phantomsection
\addcontentsline{toc}{subsection}{\textit{Adolescentiores (luniores) autem uiduas deuita. }}
\subsection*{\textit{Adolescentiores (luniores) autem uiduas deuita. }}
\section*{COMMEN. IN PRIOREM }
\marginpar{[ p.130 ]}\pstart Noli ad stipem ecclesiae susapere, et ad contuber nium sanctarum uiduarum iuniores sexaginta annis, periculum enim est, ne carnis illecebra ante id etatis non tota oessauerit, ne uicta libidine uel se matri monio iungat, uel inurat diffamiam ecclesiae intem perantia uitae, ne suspitio subeat cohabitationum uel collegium caeterarum, ocio dissolata, uel titillatione li£ bidinis. Non est tutum imo suspitione non uacat, frequem tatio episcopi ad uiduas iuniores, a qua omnimodo antistiti et doctori abstinendum, Non exoepit Aposto lus negotium confessionis et secre ti conalii conscientiae et scrupulandorum explicandorum, quod hactenus multum ualuit, et ut neoessarium a monachis et ple banis praedicatum, inomnibus Apostolorumepistolis et ordinationibus neglectum, ut nec semel hoc iniumxerit uel admonuerit quod hodie solum inculcant reli giosi uiri etc. Viduas intelligi puto omnes iuniores foeminas extra matrimoniumuiuentes, ut sunt maxi me sanctimoniales, et ubicunque sandalum, uiolentaque ́ suspitio nasci possit.  \pend
\phantomsection
\addcontentsline{toc}{subsection}{\textit{Cum enim luxuriatae fuerint in Christo nubere uolunt, habentes damnationem, quia primam fidem irritam fecerunt simul autem et ociosae discunt circuire domos, non solum otiosae, sed et uerbo}}
\subsection*{\textit{Cum enim luxuriatae fuerint in Christo nubere uolunt, habentes damnationem, quia primam fidem irritam fecerunt simul autem et ociosae discunt circuire domos, non solum otiosae, sed et uerbo}}
\section*{TIMOTHEI EPISTOLAM. }
\marginpar{[ p.131 ]}
\phantomsection
\addcontentsline{toc}{subsection}{\textit{sae et curiosae loquentes quae non oportet. Erasmus: lasciuire coeperint aduersus Christum. et: Primam fidem reiecerunt }}
\subsection*{\textit{sae et curiosae loquentes quae non oportet. Erasmus: lasciuire coeperint aduersus Christum. et: Primam fidem reiecerunt }}\pstart Iuuenculae uiduae adepta per mariti mortem libertate, et discusso iug timoris et uerecundiae, exper taque uoluptate, dum licentius uiuere inaptunt, et non uitare ooeasiones impudicitiae, incipiunt periclitari, et uina libidine, cohabitatione et incauta familiarita te. Et sic sensim prolabuntur in luxuriam et lapsum crnis sandalizantque et infamant Christi ecclesiam, et Chrissianum nomen quod propter cas blasphe matur. Et quod secretius egerunt, publica diffamia de tegitur, et luxuria earum propalatur. Quam ignominia dum euadere et contegere uel excusare aupiunt, nubere uolunt, et matrimonium praetexunt uel de nouo ineunt.  \pend\pstart Habentes damnationem ) Primum quidem in consaentia, et penes se consaae sibi turpitudinis. Deinde etiam publice reiecisse fidem comprobantur, dum fi ctam caslitatem fuisse testantur, et non ueram, nunquanque ; ex animi sententia et ex fide continuisse uidentur, sed ex necessitate dudum aliud facturae si per timorem et ze lum mariti licuisset, si tamennon simile temptasse ma ritata credatur:facile enim peiora persuadentur.  \pend\pstart Simul autem et otiosae discunt enim circu. IDu plo plaxt his uiduis licentia post coactam continen  \pend
\section*{COMMEN. IN PRIOREM }
\marginpar{[ p.132 ]}\pstart tiam gaudent se a nemine impediri, domi solitudo displicet, inuitatio passim plaoet. Dum non est eis curam da familia domus, ociosae sunt, parta dilapidant donec extat aliquid, ociosae circuncursant per alienas aedes, rumores circunferunt et comportant, et au gent suo more mendaciis, uerbosae, insigniter loqua ces et nugaces, uanae et leues, neminem uerentes, li bertatem suam in carnis delitias ordinantes, curiosae rerum nouarum, et quicquid passim dicitur uel agitur perquirentes, totam diem ociose, uerbose, et curiose exigunt, et uanae leuesque ́, dissolutae, domum redeunt in confusionem reliquarum uiduarum et de mortui mariti, propriamque ́. Garriunt perpetuo uana et otiosa, et quae eam prorsus dedeoeant. Tandemque ́ abnegata fide, et infidelium foeminarum more uiuem tes, aduersus Christum instituunt blasphemiam, con funduntque et diffamant ecelesiam, prodentes insolentiam, hactenus timore poenae tantum dissimulatam.  \pend
\phantomsection
\addcontentsline{toc}{subsection}{\textit{Volo autem iuniores nubere, filios gignere, domum administrare, nullam occasionem dare aduersario, ut habeat maledicendi causam. lam enim quaedam conuersae sunt retrorsum satanam secutae. }}
\subsection*{\textit{Volo autem iuniores nubere, filios gignere, domum administrare, nullam occasionem dare aduersario, ut habeat maledicendi causam. lam enim quaedam conuersae sunt retrorsum satanam secutae. }}\pstart Ad precauendum tot malis, uolo magis ut iunio res uiduc non tan libere uiuant, sed denuo nubant  \pend
\section*{TIMOTHEI EPISTOLAM. }
\marginpar{[ p.133 ]}\pstart marito, ut quae se regere nesaunt in offtio et hone state, per uiros seruentur, amputent lasciuiae causas, procosque et lenones euadant. Habeant quod domi curent. Mariti copia, ardori libidinis medentur, et probrum ac diffamiam effugiant. Filios et filias gignendo humilientur carne et sollicitudine. Habeant ocupationes dignas iuuenculis, filios educando, cu rando, instituendo: prae illis uoluptatibus, aliae omnes oessent, pruritus et libidines. Curiosa sit circa domus necessitates, familiam sollicite regat, matrisfami lians curam gerat et maturitatem. Exemplo praesit ancillis et filiis. Vni uiro decenter obseruiat, casttque ; placere meditetur, et inuigilare, ut quae maritus do mui prouiderit, ipsa utiliter administret. Haec Apostoli uoluntas nemo dubitet praeceptum esse, si alioqui periculum sit incontinentiae, ocii, diffamiae et indi sciplinatae conuersationis cum concordet ratio, natura, aequitas, et uerbum dei quod non omnes apiunt.  \pend\pstart Nullam oausionem dare aduersario maledicti gracia ) Sic enim nubentes euadunt non solum peri culum incontinentiae et luxuriae, feruoremque urentis libidinis, sed et diffamiam et obloquutionem declinabunt et ocasionem non habebunt aduersa rii et detractores, malesuspicandi uel maleloquendi de eis, et ecclesia nihil confundetur. Amputanda  \pend
\section*{COMMEN. IN PRIOREM }
\marginpar{[ p.134 ]}\pstart enim quantum fieri possit causa omnis iustae suspitio nis, maxime mulieribus, quibus ob sexus infirmitatem facile imponitur: sunt enim leues et propensis a fectibus, diffiale reluctantes. Vix illibaia fama defenditur, dum omnia eurant, multominus dum libertate indulgetur.  \pend\pstart Iam enim quaedam deflexerunt post satanam ) Admo net nos nota omnibus neetssariis, quod uideamus ex huiusmodi uiduis tunioribus quasdam iurpiter uiuere, et sui diffamiam incurrisse, et uidu arum ecclesiasticarum oellegum sedasse, quod licentio sa conuersatione sua propositum castitatis abieoerint, et ecclesiam confuderint, et post manifi statam libidinem et fidem repulsam, nuptiis sese excusare uel in totum post satanae tentamenta abiisse, experientia docuit, et a Christo cui sese uicturas in castitate instituerant, ora tionibus et bonis operibus incumbendo seruire proposuerant, turpiter et cum probro ecclesiae recesserunt quod ob uotum emissum et recissum desauerunt a Christo non dicitur, quod coactae fuerint ad redeundum et per petuo seruandum quodinstituerant non habetur. Non ex communtata, non incarcerata, non reducta aliqua, non con pulla scribitur. Suo periculo domino suo oecidit uel ste tit. Non denique cancellis obseratae erant, ut nunc, alioqui sicdisaurrereet nugari non potuissent, uel abire post  \pend
\section*{TIMOTHEI EPISTOLAM. }
\marginpar{[ p.135 ]}\pstart satanam. Omnia libere seruabantur, et sine uotis usque  id nouissima tempora, post millenarium, quando pro tabiliendis examinibus monac horum uotorum uinculo constringi oportere, definierunt, et obtinuerunt per tyrannidem in laqueum conscientiarum sine dei uerbo Iii tmipt uina nueuangener conuersationis et uitae.  \pend
\phantomsection
\addcontentsline{toc}{subsection}{\textit{Quod si quis fidelis, aut si qua fidelis habeat uiduas, subministret illis ut non grauetur ecclesia, ut his quae uerae uiduae sunt sufficiat. }}
\subsection*{\textit{Quod si quis fidelis, aut si qua fidelis habeat uiduas, subministret illis ut non grauetur ecclesia, ut his quae uerae uiduae sunt sufficiat. }}
\phantomsection
\addcontentsline{toc}{subsection}{\textit{Erasmus: suppeditet illis, et non oneretur ecclesia, ut iis, quae uerae uiduae sunt, suppetat. }}
\subsection*{\textit{Erasmus: suppeditet illis, et non oneretur ecclesia, ut iis, quae uerae uiduae sunt, suppetat. }}\pstart Siquis uir fidelis uel mulier, cognatam habeat ui duam, quam de suis faciltatibus commode nutrire possit, et uitae neoessaria ministrare, non eam patiatur ob trudi ecclesiae, quae solet uere uiduas, id est, deftitutas omni subsidio alere, de communibus ecclesiasticis prouentibus, ne fortassis dum omnes passim tam uiduae quam aliae ecclesiasticae personae alendae fuerint, non suppetant uel sufficiant donationes collectae uel prouentus. Iterumuidemus ad quid sint bona ecclesi affica, decimae et oblationes ordinanda, et caetera et quae olim consuetudo, quid nostris temporibus emendandum, qui ccclesiaficisint uere.  \pend
\section*{COMMEN. IN PRIOREM }
\marginpar{[ p.136 ]}
\phantomsection
\addcontentsline{toc}{subsection}{\textit{Qui bene praesunt presbyteri duplici honore digni habeantur, maximem qui laborant in uerbo et doctrina Dicit enim scriptura: Non alligabis os boui trituranti. Et. Dignus est operarius mercede sua. }}
\subsection*{\textit{Qui bene praesunt presbyteri duplici honore digni habeantur, maximem qui laborant in uerbo et doctrina Dicit enim scriptura: Non alligabis os boui trituranti. Et. Dignus est operarius mercede sua. }}
\phantomsection
\addcontentsline{toc}{subsection}{\textit{Erasmus: Boui trituranti non obligabis os. }}
\subsection*{\textit{Erasmus: Boui trituranti non obligabis os. }}\pstart Sunt qui male praesunt presbyteri, crisoepi salicet et diacont, qui nullam habent curam de subditis, de pauperibus uiduis, de uerbo det praedicando, conso lendo afflictos, fragiles in conscientia instruendo. Sibi solis incumbunt, ociantur, libidinantur, fastum gerunt, ambiunt honores et quaestum, ab umunt ec clesiastica bona. Quid illis debetur, quamut eum sirt sal infatuatum, et ad nihilum ualent ultra, quamut eiiciantur de sedibus, excommunicentur, confundartur, accusentur publice, si non nimis publice sese acru sant. Et pedibus concilcentur, destituantur honore, diuitiis, officio, et publica abiectione conculcentur ab omnibus: ne si permittantur in sua scandalosa uita, et negligatur populus, et obtendatur peaant um autoritas quasi patrocinium iniquitatis, et fat sicut sacerdos, sic et populus. Et lupus oueis dispergt, cum pastor lupus fuerit, et mercennarii foedus inerint cum lupis in perniciem gregis Qui ueron bene prae  \pend
\section*{TIMOTHEI EPISTOLAM. }
\marginpar{[ p.137 ]}\pstart Junt episoepi, qui hic dicuntur et ubique in Paulo presbyteri, nisi ubi manifiste aetatem indicat solam nomen presbyter, tales duplici honore dignt habeantur. Fpiicopi uero ut bene praesint, sint prestyteri, id est, senieres, non tam aetate quam morum grauitate et sapientia. Iuuienis enim erat Timotheus et Titus: sed cani sunt sensus bomns, tales episcpi et senes honore presbyterii, duplici honore digni habeantur. Honor dupliciter sumitur, pro reuerentia uul gari, ut est decedere de uia, nudare caput, flecti re ge nu, quae gratitudinem indicant, et appreciationem et agnuionem donorum dei in illo. Et est magis honorantis quam honorati. Alias ut hoc loco honor est munus, elecmosyna, et sustentatio, prouisioque ; necesa riorum honesta et sufficiens, quae debetur dectori, non ut profundat in usus mundanos, sed uiuat modeste, et sit omnis parsimoniae et maturitatis exemplar et speculum uitae Christanae, hoc est cuangelicae, imo apost licae, quorum sucessores esse de bent et exhibere. Hoc duplici potius honore quam duplici prouisione honorentur, cum nhil superflut cos decrat, sed uictum quil us tegeniur habentes et c. Si qui ucro fuerint prestyteri et antisttes uel ecclesiae praeficti, qui non solum praesunt integritate uitae, consi lio et prudentia, ac in adminislratione bonorum ec  \pend
\section*{COMMEN. IN PRIOREM }
\marginpar{[ p.138 ]}\pstart clesiasticorum et Christianae conuersationis, quiett fouendae incumbant, et rectores gerunt, sed ctiam episcopos et doctores exhibent. Hoc est, laborent legendo, docndo, exhortando, incilcando et dispem sando doctrina Christi, euangelicam, apostolicam uerbo et doctrina prosint sicut praesunt ecclesiae. istis duplex sic honor omnium maxime habeatur, et referatur, tantoque diligentius, quanto sibi suisque sacris il lis addictus exercitiis minus prospicere potest et de uitae necessariis prouidere. Duplici quoque honore dignus dicitur, ut et eo magis quae sufficiunt et simplicem et honestam compctentiam assequantur. Non enim duplum preciosorum uictum sibi et uestitum con parare debet, qui frugalitatis specimen sit. Hospitali tatis enim et eleemosynarum largtio aliunde sumenda de bonis scilicet ecclesiasticis, cuius dispensator sit  \pend\pstart Dicit enim scriptura: Non alligabis ) Non infre. nabis os boui trituranti, Deut.25. Mos est a multis bobus circumactis uel equis trituram facere, si ergo laborantibus crudele fuerit, negare cibum, omnino enim qui ecclesiae opus exequitur, qui docet, instruit hortatur, sacrameeta ministrat, priuatim et publice is dignus estut ab ecclesia nutriatur et foueatur. Quid autem de illis qui nihil tale laborant, nunquam sacra legunt ut docoeant, uon prae dicant, non monent, non con  \pend
\section*{TIMOTHEI EPISTOLAM. }
\marginpar{[ p.139 ]}\pstart solantur conscentias, sed omne opus corum est in ligan dis et excarnificandisque et domandis conscientiis, tot praeceptis et excommunicationibus. Nihil loquuntur ad populum nisi leges suas, toti sunt in augendis, oelligem dis, recensendis et absumendis sumptibus et pecuniis numerandis. Nec prae dicant unquam, nec euangelium prae dicari sinunt, nisi coacti timore amittendi luxus, quae stus et fastus. Paulus de dispensatoribus uerbi dei haec atat, non de cantoribus, organistis, monachis et non nis. Quod si tot ecclesiarum presbyteri sacras scriptu ras ad fructum aliquem legerent uel audirent, ut idonci surrogari possent ad pastorale officium et digni iudicri forsan uictu. Nunc autem boues triturantes id est agricolae toto anno laborantes, frau lantur durissimo labore suo Vinitores aquam bibere coguntur, cum uxoribus, filiis et famulis dura inedia macrantur ut uixuiuere possint an non hoc est boum ora constring et infrenari uix tole rabiliter laborantes, et clericos ociosos et scandalissi mos fouere in perniciosissimum soelus mundi etc. multa?  \pend\pstart Dignus est enlm operarius mer. ) Lucae io dicit Christus et grande piaculum fraudare uel usque ad noctem laborantem meraede. Operarium dicit non ocosum, non stertentem, non aliorum laborum partis auare inhiantem et solis uoluptatibus operam dantem. Non sunt operarii, qui tota die in foro stant ocosi, sed qui pondus ferunt dici et astus, ꝗ ludunt alea et colludunt mu  \pend
\section*{COMMEN. IN PRIOREM }
\marginpar{[ p.140 ]}\pstart lierculis et illudunt uino et lauticiis, qui otiosi, et ad nihil utiles laborantium sudores exigunt et necessariis fraudant. Operarios dicit qui laborant in uerbo et doctrina, qui praesunt, et bene, non obsunt, qui aduigilant gregi, qui in ecclesiae utilitatem proficiunt. Hodiernus cantus et uigiliae nec deo placet, nec hominibus prodest, sed impedit a melioribus, ab oratione, a uerbo dei ediscendo, au diendo et dooendo, nec populo, nec clero, uel in minimo utilis. In eum usum a patribus nunquam insitutum fuisset hoc officium.  \pend\pstart Meroede sua) Merces respondet obsequio et operi, fraus erit et dolus et rapina, meroedis nomine exigere, cum nihil egeris utile, uel gratum. A ple be sustentari uelle cui non seruias, nihil prosis, sed ob sis maxime et impedias profuturos, et aliorum mer cedes operum tibi rapias, deque ; sanguine et sudoribus aliorum sagineris, ut hodic oernere est in praela tis ecclesiarum, et beneficiosis qui sceleratissimi ele rici et populi oppressores et opprobrium Christi, uiolenter pauperculis pastoribus ecclesiarum mercedes deglutiunt, et ad ructum uorant. Et qui infe liiores et indigni magis, residere et uicarium agere sustinent, hi deliguntur et populo prae ciuntur a ne bulonibus istis. Et id non diutius, quam euangelio  \pend
\section*{TIMOTHEI EPISTOLAM. }
\marginpar{[ p.141 ]}\pstart praedicando pure et libere destiterint, statim abiiaem di, dum Christum et mores Christianos, simul ac fidem pietatemque ; coeperint asserere, ne populus sapiat ali quando et meroem neget oüosis. Dignus est itaque operarius, sed indignus certe, fallax, nequam, infidelis, ociosusque saudator, benefacta inficiens, et opera tores studiosos absterrens, et sibi soli intendens, suaeque libidini fastui et auaritiae, dignior certe ut mitta tur mola asinaria etc. et ut conclcetur ab hominibus etc. et ligatis manibus et pedibus mittatur etcae. et partem cum infideli acipiet etc.  \pend
\phantomsection
\addcontentsline{toc}{subsection}{\textit{Aduersus presbyterum accusationem noli accipere, nisi sub duobus et tribus testibus. Eos qui peccant coram omnibus argue, ut et caeteri timorem habeant. }}
\subsection*{\textit{Aduersus presbyterum accusationem noli accipere, nisi sub duobus et tribus testibus. Eos qui peccant coram omnibus argue, ut et caeteri timorem habeant. }}\pstart Inoertum iuxta Erasmum, an de seniore loquatur, an de episcopo, mihi de utrisque uidetur acipiendum. Episcopi enim sunt ecclesiae seniores, qui eidem com nuni condlio et sapientia pro autoritat sua prospicere habent de omnibus. Non enim magnae ciuitatis tot populi, unius prudentiae et spiritut debent niti, sed de seniorum concilio, senatum quasi ha beant, cum quibus iuxta deiuerbum et regulam tractet omnia. Sintque ; unius ciuitatis hoc modo plures epsoopi, imon presbyteri, non tam aetate quam ser su  \pend
\section*{COMMEN. IN PRIOREM }
\marginpar{[ p.142 ]}\pstart et moribus, qui omnes apud congregationem fidelium in autoritate habeantur, etiam si non omnes dooeant et praesint uerbo et doctrina, assistant tamen episcopo et presbyteri sint plebis honorabiles et seniores eccle siae. Son enim legimus discrimem fuisse tempore aposto lorum et status intra ecclesiam ut pars una religiosa, cle ricalis, spiritualis et ecclesiastica fuerit. Alia autem pro phana, saecularis, mundana, carnalis. Fuerunt enim omnes Christiant religiosi, spirituales utentes hoc mun do tanquam non utentes, uacantes omnes scripturae di uinae prot empore congruo, et de labore manuum ui ctitantes: id enim docuit apostolus Paulus uas electio nis et apostolus in signis et exemplo firmauit, qui non laborat in ecclesia, non manduct inquit. Quod si qui saeculariter, mundane, carnaliter et sine dei timore et disciplina ecclesiastica uiuebant, Christiani non erant, sed post satanam ire iudicabantur et euitabantur tanquam haeretici, et ad sacrosancta mysteria non admittebantur. Seniores itaque Christianorum bonorum presbyteri dioebantur, inter quos qui uerbo dei et doctrina ualebat, et donis spirituss. nobilior, is episcopus dioebatur, et erat magis offitio docendi prae cellens, quam potestate et imperio, De istis presbyte ris ecclesiae honoratioribus, scilicet et autoritate et sapientia probitateque prestantioribus hic dicitur:  \pend
\section*{TIMOTHEI EPISTOLAM. }
\marginpar{[ p.143 ]}\pstart Aduersus presbyterum et grandaeuum et maturum uirum acusationem non recipias, nisi sub duobus aut tribus testibus. Poterant insolentes a talibus presby teris correpti et admoniti facile malignari, et comminisci mendacia, redundaret hoc in probrum ecclesiae, eleuaretur autoritas uirorum. Ideo non farile admittendi calumniatores, ne ansa praebeatur personis leuibus temere lacescendi seniores ecclesiae.  \pend\pstart Omnes christiani, pie, religioseque uiuentes, Eccle sia nominabatur, licet adessent hypocrytae et speu dochristiani, sed secretis cordium dei est, non ecclesiae iuditium. Non uno aut duobus testibus, sed pluribus tam spectabilis uiri a facie acusari debebant, ne facile concinnaretur mendacium contra honestissimos ab indignis et malis uindicare se uolentibus de cassgatione, sicut de Quadrato et Ale xaudro papa.  \pend\pstart Pecuntes autem coram omnibus argue ) Ad pre sbyteros referendum eadem ratione uidetur, ut siquis inueteratus dierum senex, sic publice insolenter pecet, ut testibus fide dignis et multis crimen testari possit cum id cedat in multorum scandalum et totius ecclesiae et Christianit atis, licet hi sic conuicti publice et coram omnibus argui debent, ut caeteri timo rem habeant et non  \pend
\section*{COMMEN. IN PRIOREM }
\marginpar{[ p.144 ]}\pstart personarum aceptio locum in ecclesia habere uide atur, ne omnia hypocritice geri in ecclesia uideri possint ab in idelibus, et ut criminum punitio, scandalum tollat innocentis ecclesiae, quae non tam ex innocentia uitae, quan'ex zelo gloriae diuinae, et hone statis Christianae dinoscitur. Non sunt omnes in ecclesia innocentes, sed criminum osores, et uindices et zelatores omnes esse debent uere ecclesiastici. Et cum uiderint etiam maioribus et presbyteris non para et conniueri, sibi quoque non parcendum noue rint, si criminosi fuerint. Nihil hic dicit Paulus de correctione episcopi in specie: tum quia Timotheus probatae probitatis erat, tum quia in primitiua eccle sia episcopi assumebantur ex pluribus sanctis sancti ores :uerum presbyteri nomine satis indicat coram pluribus uel ommibus arguendos etiam episcopos si sandalosi fuerint. Et quanto lues contagiosior, si epi scopus criminibns obnoxius sit, tanto adhibitis suffici entibus attestationibus, durius et seuerius arguendus. Hic iterum et semper episcopi et presbyteri nomine ecclesiae praefectus et parrochus quilibet et plebanus intelligi debet. Ex horum enim doctrina et moribus pendet profectus, detrimentumque ecclefiae. Longeque maturius deligendi essent, testimonio et sufragio plebis quam magstratus ciuitatum, ubi non  \pend
\section*{TIMOTHEI EPISTOLAM. }
\marginpar{[ p.145 ]}\pstart parua cura adhiberi solet, ut quam integer fama et moribus et rerum experientia et spectata probitate probatus et conspicuus deligatur, misi maioris aesti mare uelint rectorem auium quam doctorem anima rum. Verum sicut superius Paulus de uiduis, quae ue re uiduae sunt tractans eas specialiter, et docuit tamen generalem regulam, scripsit de domesticis maxime omnium fouendis, sic hoc loco uidetur loquens de philosophorum arguitione et contestatione subin ferre generalem regulam de omnibus publicis et scandalosis in ecclesia pecatoribus, ut tales coram omnibus arguantur ab episcopo, ut caeteri timorem habeant similia committendi. Similiter crimen esse non facile, imo nunquam recipiendum, nisi appareaut et conuincant illud testes idonei. Sic regulae genera les notantur cum de negotiis et personis specialibus sermo incidit nonnunquam, plures tamen testes con tra senes et spectuales uiros quam contra leues personas semper exigendi sunt.  \pend
\phantomsection
\addcontentsline{toc}{subsection}{\textit{Iestificor (Obtestor) coram deo et domino lesu Christo et electis angelis cius, ut haec custodias sine praecipitatione (et sine praeiuditio) iuditii, nihil faciens iuxta propensionem animi) in aliam partem declinando. }}
\subsection*{\textit{Iestificor (Obtestor) coram deo et domino lesu Christo et electis angelis cius, ut haec custodias sine praecipitatione (et sine praeiuditio) iuditii, nihil faciens iuxta propensionem animi) in aliam partem declinando. }}
\section*{COMMEN. IN PRIOREM }
\marginpar{[ p.146 ]}\pstart Su b testimonia et praesentia dei omnia cernentis, et domini Iesu Christi, qui nobiscum est semper et ubique , usque ad consummationem saeculi, et quasi in pro spectu angelorum, et ministrorum dei, qui fidelium san ctis negotiis iugiter intersunt, ego temi Timothee, adiuro et obtestor, ut praeceptor discipulum, ut non patiaris praeiudicium aliquod fieri uerbo dei, quod me referente audis et suscipis, quin praecise insistas illi tanquam regulae ueritatis et Christianae doctrinae ne in aliquo detrimentum capiat pura ueritas, et labe factetur, adultereturque ; sicut a speudo disapulis. Nihil facias in aliam partem declinando, hoc est, praeter regulam fidei et scripturae, et iuxta humanum sensum et iuditium nihil dooeas aliquando.  \pend\pstart Erasmi translatio uult, ut nihil iudicet, etiam sub testibus, siue paucis, siue multis cum praecipitatione iudicii, inconsiderate et inconsulte, et friuole uel festinato. Cum iunior esset Timotheus et inge nio feruentiore, et praeuolatili, et ut non indnlgeret iudicando, affectionibus, et naturae inclinationi et propensiori animo in ea quae libent et plaoent, sed iudicet circunspecte et mature, et deliberatione di ligenti praehabita, et non iuxta sensum naturae et ad personarum aceptionem, sed iuste et aequaliter, fideli terque pronunciaret, sicut in oculis domini dei, et iesu  \pend
\section*{TIMOTHEI EPISTOLAM. }
\marginpar{[ p.147 ]}\pstart Christi et angelorum, ut similiter nosque a dominort Christo Iesu iudicandos cogitenus: qua enim mensu ra etc. Discimus hoc loco primum quam inuitus ueni at Apo tolus ad mandandum et praecipiendum, obte ffatur potius et armatis precibus contestatur. Deum semper prae oculis habeat iudicaturus, nihil praecipi tauter, nihil praeter regulam uerbi dei, nihil ad gra dam uel displioentiam, sed iuste iudicet episcopus Angelos dei missos in ministerium et testimonium actuum no strorum, quos et reuereri semper debemus, et cauere a peaatis, non solum deo et Christo domino aspi cientibus, sed etiam exeratibus angelicis qui assistunt semper: et tamen in caelis semper aspiciunt patrem etc. Id uero faciendum episcopo, non solum in iudiciis sed in ommbus aliis administrationibus, ut prouide, fideli ter, sinoere et candidu ipu omna qui praetsi ecclesiae  \pend
\phantomsection
\addcontentsline{toc}{subsection}{\textit{Manus cito nemini imposueris, neque  communicaueris peccatis alienis. }}
\subsection*{\textit{Manus cito nemini imposueris, neque  communicaueris peccatis alienis. }}\pstart Apostoli cum mitterent aliquos ad praedicandam fidem et euangelium uel officium uerbi imponebant, ute bantur impositione manuum, et dabatur talibus spiritus sancti donum, ad hoc idoneum et sufficiens in prinuti ua ecclesia. Et quia subintrabant multi spendoapst li et mali Christiani sese admisoebant, et ob quaestum et glo ria temptabant opus apostolicum, ideo apostolus praemo  \pend
\section*{COMMEN. IN PRIOREM }
\marginpar{[ p.148 ]}\pstart net Timotheum, ut pspiciat sibi ne his officium p raedica tionis deleget uel committat, nisi prius fidem et mores et studia pbe cognouerit, ne malae et erroneae doctrinae illius et seductionis ipse causa uideri possit, qui talem ad id ordinauerit, neque error humanus et falsitas spi£ ritus sancti doctrina aestimari possit, qui manuum im positioni uel dabatur uel dari significabatur: erat enim attestatio quaedam sanctitatis in doctrina et uita, quae non cito exhiberi decebat. Qui enim doctorem malum instituerit, errorum illius et seductio nis in populo causa erit, taliumque peaatis communiabit semper. Ideo sumopere discutienda essent praedicatorum doctrina et studia, mores et uita, hoc est praelatoribus ecclesiae, quorum officium primari um et prindpale est, docre uerbum ueritatis.  \pend
\phantomsection
\addcontentsline{toc}{subsection}{\textit{Te ipsum castum custodi. }}
\subsection*{\textit{Te ipsum castum custodi. }}\pstart Caste et pure hactenus uixisti, et iuuentutem tu am disciplinis seuerioribus instituisti ad castitatem fa dlius seruandam. Eam custodiam serues, non ut ideo castus sis, cum dei donum fit castitas uera, sed in seruitutem redigere, et abstinentiis cruafigere et con primere fomitem ut hactenus, non obmittas quod con ducant ualde castitati seruandae in integra adhuc aetate, dum feruet sanguis et caro. Verumita serua absainentiam ut non mmis, ne per corporis imbeall  \pend
\section*{TIMOTHEI EPISTOLAM. }
\marginpar{[ p.149 ]}\pstart tatem et nimiam attenuationem periclitetur ualitudo cor poris, et impediantur opera necessariae charitatis.  \pend
\phantomsection
\addcontentsline{toc}{subsection}{\textit{Noli adhuc aquam bibere, sed modicouino utere propter stomachum tuum, et frequentes tuas infirmitates. }}
\subsection*{\textit{Noli adhuc aquam bibere, sed modicouino utere propter stomachum tuum, et frequentes tuas infirmitates. }}\pstart Ne posthac aquam solam bibas, quae obest corporis tui ualetudini, ne dum castigas nimium, sub one re portando sucumbas. Non propter se bona macratio corporis, sed propter exercitia uirtutum. Aqua multis rationibus obest, et tui corporis conditionem non fert nimiam. Vino tibi opus est ob stomachum de bilem. Eo utere ad necessitatem, sed modice, nisi id corporis debilitas exigeret, ego a te non exigerem ui ni potum. Nos puelli bibimus uinum, et adolescentes potamus, uiri inebriantur. Monachi et clerici nephas et miseriam putant aquae potum, et poena maximam. Episcopi exquisita uina potantuel electa Sic praestamus apostolicos. Nimius etiam aquae potus in uitio est. Non quid bibas uel comedas, sed qua tum curandum. Modus imponitur potationi, ut quam tum seruiat incolumitati bibant. Non ad medicos, purgationes, confortationes mittit, sed diaetam indicit.  \pend
\phantomsection
\addcontentsline{toc}{subsection}{\textit{Quorundam hominum peccata manifesta sunt praecedentia ad iuditium: quorundam autem sequuntur. Similiter et }}
\subsection*{\textit{Quorundam hominum peccata manifesta sunt praecedentia ad iuditium: quorundam autem sequuntur. Similiter et }}
\section*{COMMEN. IN PRIOREM }
\marginpar{[ p.150 ]}
\phantomsection
\addcontentsline{toc}{subsection}{\textit{fact a bona manifesta sunt, et quae aliter se habent abscondi non possunt. Frasmus: Quorundam homiuum peccata ante manifesta sunt, precedentium ad iudi cisi, quosdamuero et subsequuntur. Con similiter et bona opera ante manifesta sunt, et ea quae secus habent, occuliari non possunt. }}
\subsection*{\textit{fact a bona manifesta sunt, et quae aliter se habent abscondi non possunt. Frasmus: Quorundam homiuum peccata ante manifesta sunt, precedentium ad iudi cisi, quosdamuero et subsequuntur. Con similiter et bona opera ante manifesta sunt, et ea quae secus habent, occuliari non possunt. }}\pstart Praeueniunt iuditium, et non indigent attestati onibus ea hominum multorum pecuta quae manifiste, in uerecunde, audacter committuntur in prospectu hominum et totius ecclesiae, huiusmodi peanta ad publi cum etiam iuditium delata puniantur, nec test moniis multis ibi incumbendum, et inquisitione scrupuloia illa tu quoque solus punire potes, et ecclesia debet. Quae uero peaata sequuntur iuditium id est, nesciuntur nisi idoneis et pauais testbus comprobentur, haec pos sunt abscondi et corrigi secreta admonitiont uel corre ctione ab episcopo, et aliis quae alioqui deo iudici et extremo iudicio seruabuntur. Non omnia hic puni ri et emendari possunt, non sunt de foro episopi uel ecclesiae nisi manifesta.  \pend\pstart Similiter et bona opera quaedam manifiste fiunt non egentes discssione, nec sinistre ob timorem oculte nuper electionis criminari debent, qusi licet bona  \pend
\section*{TIMOTHEI EPISTOLAM. }
\marginpar{[ p.151 ]}\pstart sint, tam em quia mala in conscientia esse possunt, ideo nibil sint curanda et abseondenda. Lioet enimomnia nostra opera defectuosa sint, et dei iuditio expendem dadeprehendantur mala, quia tamen speciem mali non habent, non sunt obmittenda uel ocultanda. Quae uero aliter se habent, bona sunt tantum in ap£ parentia, et reuera mala peritius in conscientia, illa suo tempori et iudici obscondi non poterunt. Nihil enim ocultum quod non sactur.  \pend\pstart Quidam homines nunc fic uiuunt, ut ex corum factis manife ste malis, antequam moriantur, uel a deo iudicentur damnandi peaatores agnoscantur a nobis: qui seilicet sine omni terguersatione omnibus criminibus et publicis fese infamant. Quorundam pecata mala quaedam sunt, sed hoc ut uel dubia, uel oculta non iudicantur, sed deo tantum post hanc ui tam diiudicanda reseruantur, de quibus nihil ad nos suo domino et suo tempori stant uel cadunt etiam. Confimiliter aliquorum hominum opera manifeste bona uidentur, ut nemo ea nisi malus criminari debeat uel possit, illa sine dubio laudanda a nobis sunt et de eis nemo debet uerecundari, neced propter ma los omittere qui omnia etiam bona opera infamare gaudent et student, sunt tamen bona opera in speciem et hypocritica, quae apparent bona, sed non sunt. Illa  \pend
\section*{COMMEN. IN PRIOREM }
\marginpar{[ p.152 ]}\pstart deo manifesta quidem sunt, sed hominibus ignota, suo tamen tempore apparebunt qualia fuerint. Nihil enim opertum quod non reueletur. Abstinentia autem a uino et similibus, nec manifeste bona nec mala est, sed ob ocultas causas, necessaria uel indigna, deo grata uel ingrata apparebunt suo tempore et discuaentur.  \pend
\endnumbering\beginnumbering\section{CAPVT VI.}
\phantomsection
\addcontentsline{toc}{subsection}{\textit{\huge\textbf{Q}\normalsize Vicunque sunt sub iugo serui, dominos suos omni honore dignos arbitrentur, ne nomen domini et doctrina blasphemetur. }}
\subsection*{\textit{\huge\textbf{Q}\normalsize Vicunque sunt sub iugo serui, dominos suos omni honore dignos arbitrentur, ne nomen domini et doctrina blasphemetur. }}\pstart Serui empticii infidelium dominorum conuersi ad fidem unius dei ab idolalatria, pacienter ferent iugum obedientiae et laborum et status sui sortem. Nec quia conuersi ad Christum fraudare uelint dominos seruitio suo, et transfugae uel profugi fiant ob fidem dei, quasi contemnentes dominos ob infidelitaten, qui bus se seruire indignum putent; sed magis humilibus obsequiis seruituti satis faciant, cauentes ab omni of fensa eorum, honore debito colentes et dignos arbitrantes, a deo sibi praepositos cogitantes, non casu, hic esse dei uoluntatem, ut boni exeroeantur a malis ut mali cum deo placuerit, conuertantur a bonis, ut  \pend
\section*{TIMOTHEI EPISTOLAM. }
\marginpar{[ p.153 ]}\pstart inclarescatuirtus seruorum dei, et Christi gloria a suis se ftr toribus illustretur, ut seminarium seditionis non detur ethnicis et oaasio infidelibus malignandi, ne nomen det optimi maximi et solius per idololo tras blasphemetur, quasi non aequus sit qui ubeat do minos, suis seruis destitui et passim gentiles irritari a scruis, parem infringi, obedientiam cassari, maiores sperni, pronissa negari, commertia rumpi, deumque  uerum insolentia seruitorum blasphemari. Similiter et doctrina euangelica infamabitur si eius professo res, mundi ordinem et magistratuum et dominorum obedientiam, detrectauerint, et omnis religio Christiana, propter eos maleaudiet. Et scandalizabuntur in doctrina Christi idololatrae, qui potius ad meliorem mentem exemplo fidelium debuerant aedi ficari et allici et conuerti, unaque nobiscum aliquando dominum deum nostrum glorificare inseruis suis Voluit et ludaeos dominus obedire regibus infidelibus, pro eis rogare, tributa pendere. Et Christus Caesari iussit reddi debita. Et Petrus in Epistola mul ta talia, cum baptismus quidem pecati dominium au ferat, sed ab heri dominatu non eximit. Siue illud ser uitium licitum per se sit uel illicitum, disputandum non erit, sed potius humiliter obsequendum. Si seruus commode et modo honesto liber fieri poterit, non ob  \pend
\section*{COMMEN. IN PRIOREM }
\marginpar{[ p.154 ]}\pstart mittat. Si dominus seruum manumiserit, ut et ipse a seruitute pecati liberetur a deo, optime froerit, et ad id induoendus, non contemptu, uiolentia, temeritate et diffamia iugum excutiendum.  \pend
\phantomsection
\addcontentsline{toc}{subsection}{\textit{Qui autem fideles habent dominos, non contemnant quia ftatres sunt, sed magis seruiant quia fidelessunt (inseruiant quod fideles sint) et dilecti, quia beneficii participes sunt. }}
\subsection*{\textit{Qui autem fideles habent dominos, non contemnant quia ftatres sunt, sed magis seruiant quia fidelessunt (inseruiant quod fideles sint) et dilecti, quia beneficii participes sunt. }}\pstart Serui conditione mundana, si habuerint dominos Christianos, et in Christo iam confratres sibi factos non ideo minus reuerenter seruiant et obediant qua si indigne ferendo par paribus obsequi. Pares quidem sunt in Christo et fratres in fide, sed non in uitae hu ius fortunis et dei prouidentiis, fic enim unicuique  homini deus prospexit et praeordinauit uitae ordinem et statum et exercitium in quo creatori seruiat, conmensurauitque grauamina ut incertum habeatur, ser uor um ne an dominorum conditio tollerabilior fuerit, pro dotibus naturae tam differenter passim impartitis, quin potius tanto promptius et alacrius eisdem seruiant, quanto ob Christi religionem mitius et mam suetius tractantur, et ad prophana seruitia non com pellant, sed padifice, quae Christianae religionis sunt exercitia sancta permittunt, a quibus et diliguntur in  \pend
\section*{TIMOTHEI EPISTOLAM. }
\marginpar{[ p.155 ]}\pstart Christout fratres instruuntur, ut discipuli et tanquam filii bencficiis afficiuntur, et omnium bonorum partiapes habentur. Di plici enim nemine obe dire debent fidelibus dominis promptiores. Tum quia domini eorum sunt, ium quia socii religonis esse non contempnunt. Si enim uere fideles domini sunt, seruos suos trocteburt, ut fratres potius, quam ut seruos, diligent mogis quam oentennant. Beneficiis aff cient, et inmo deratis grausminibus non oppriment. Bonorum suorum ees participcis aestimabunt, et facient. Quod si crudeliter, et sine charitate eos tracteuerint domini, non tam Chri stiant erunt, quam Domitiani, et perfidi. Verum nihilominus tollerandi sunt, etiam discholi quam diu contra Christum uniuersorum dominum nihil imperauerint. Alioqui oportebit deo obedire magis quam hominibus.  \pend
\phantomsection
\addcontentsline{toc}{subsection}{\textit{Haec doce et exhortare. }}
\subsection*{\textit{Haec doce et exhortare. }}\pstart Sic episcopum docere dect, non solum maiores et dominos, sed etiam infimae conditionis homines seruos non negligere: docere autem ea quae sunt ad pietate m, tranquillitatem et pacificam uitam, ne susatetur disadium in ecclesia, ne inobedientia iunorum ordo confundatur. Exhortari debel etiam eos quos ui derit cunctabundos et aessantes, ut crebro admom ti  \pend
\section*{COMMEN. IN PRIOREM }
\marginpar{[ p.156 ]}\pstart et consolati leuius portent onus a deo impositum, et pro quo dei fauorem et beatitudinem consequi possint. Non doce eos libertatem carnis et mund quae sae pe nooet, sed quae est in Christo et spiritus lon gammitatem, pacientiam, obedientiam.  \pend
\phantomsection
\addcontentsline{toc}{subsection}{\textit{Si quis aliter docet, et non acquiescit sanis sermonibus domini nostri lesu Christi, et ei quae secundum pietatem est doctrinae, is inflatus est, nihil sciens, sed inseruiens circa quaestiones et uerborum pugnas, ex quibus oritur, inuidia, contentiones, maledicentiae, suspitiones malae, superuacaneae conflictationes hominum mente corruptorum, et qui a ueritate priuati sunt. }}
\subsection*{\textit{Si quis aliter docet, et non acquiescit sanis sermonibus domini nostri lesu Christi, et ei quae secundum pietatem est doctrinae, is inflatus est, nihil sciens, sed inseruiens circa quaestiones et uerborum pugnas, ex quibus oritur, inuidia, contentiones, maledicentiae, suspitiones malae, superuacaneae conflictationes hominum mente corruptorum, et qui a ueritate priuati sunt. }}\pstart Refertur hoc aliter ad doctrinam apostolicam quae tam in superioribus quam insequentibus descri bitur, si quis iuxta sensum proprium et legis nationem et rationem erroneam hominum et philosophorum do cu crit, et non credit uel appreciatur ac uelut doctrinae salutis et ueritatis acquiescit, sanis, immaculatis, aequissimisque domini nostri Iesu Christi sermonibus, sed introduxerit alios etc. Istis enun tantum sermonibus adhaerendum, haec tantum docenda, haec est sola doctrina pietatis, secundum quam domino serui\pend
\section*{TIMOTHEI EPISTOLAM. }
\marginpar{[ p.157 ]}\pstart re et plaoere poterimus, de qua certi sumus, quia Christi est, qui ueritas est, qui in me loquitur, a quo acepi. Eam scire et sapere pietas est, eam negare impietas.  \pend\pstart Is inflatus est ) Superbit. Superbia est sensum proprium praeferre uerbis dei, scientia animali et diabolica inflatus contra deum. Omnis enim eruditio alia a fide, inflat, et tumorem mentis excitat, in se confidere, de proprio iuditio et sensu praesumere docet, contemptui habet simplicitatem spiritus et scripturae. Cunque crasso carnis intuitu eam attingere nequiuerit, despicit et stultam iudicat. Non enim quae spiritus sunt sapiunt carnales, sed dum se credunt sapientes stulti facti sunt.  \pend\pstart Nihil sciens) Hoc unum sciet, si quicquam scit quod ni hil sciat, qui dominum ignorat, qui in Christum non credit, qui nititur humanae rationi erroneae, cum sit omnis homo mendax, et comparatus iumentis insipientibus quid scire poterit? Scientia non est nisi ueritatis agnitio. Christus autem ueritas et lux mundi, praeter quem omnia tenebrae et ignorantiae et stultitia.  \pend\pstart Sed insaniens circa quaestiones et pugnas uerbo rum) Certum est longo tempore post apostolorum tempora quaestiones et pugnas uerborum pau\pend
\section*{COMMEN. IN PRIOREM }
\marginpar{[ p.158 ]}\pstart cas fuisse, lioet nunquam defuerint, sed per nostros Parrhisienses haec stulticia disseminata est in orbem et insanire fecit omnia ingenia, lauguereque circa quaestiones futiles et insanas uerborum pugnas, de quibus tot miseri quaeflionarii et sophstae nihil inte grum uel oertum reliquerunt, sed dubiis inuoluerunt omnia. Iusto autem dei iudicio palpabiles istas tene bras inciderunt, quia sanos sermones neglexerunt, sacras litteras philosophicis studiis postposuerunt. Veruntamen iram dei solam, in hac causa uideo de qua alias.  \pend\pstart Ex quibus oriuntur inuidiae, contentio. Jvt tac am de tot haeresibus ante tempora Augustini, de in uidia episoporum tempore Arrianorum, de contentionibus et blasphemiis atque conuiciis. Quid nostris temporibus uidimus, inuidias inter uniuersita tes. Iu uniuersitatibus, inter uias antiquorum, realium, modernorum, praedicatorum, minorum, propter quaesti ones illas et pugnis uerborum, quot et quanta scandala publica, quot umposturae et insaniae. Tot autores in sententias quorum nullus non alios omnes mipetit et uel licat ac irridet, id unum hodie mirabile uidemus, cum concordant omnes aduersus sacras litteras et euange lium, et pro uiribus omnibus ratiocinantur, sed frustra, et sine spe aliqua obtinendi. Infatuauit enim  \pend
\section*{TIMOTHEI EPISTOLAM. }
\marginpar{[ p.159 ]}\pstart dei sapientia omnes et transtulit regnum suum ad paruulos et plebeios et mulierculas, de his multa di ci possunt.  \pend\pstart Suspiciones malae ) Dum alter alterum haereti cum credit et proclamat excommunicatum, suspectum de haeresi, et in ira dei, et potestate diaboli, Tot sunt enim hodie suspitiones malae, tot criminationes haereseos ob solas eorum tradituones et quaestiones sophisticas.  \pend\pstart Superuacneae con. ) Dict aliquis conflictationes augeri hodie, dum non audiuntur uniuersitatum determi nationes. Sed quis pbabit superuacaneas cum sint de le ge dei, de euangelio, de fide, peaato, sacramentis, libero arbitrio. Item de solis scripturis diuinis, non de philo sopho et auerroë, astrologia, metaphysica, logica. Et omnia solum contra philosophastros et quaestionarios Quis dicet superuacane as conflictationes, si etiam conten damus circa linguas, arca literam interpretandam, uete ris et noui testamenti ad inquirendum sensum germa num scripturae? Et quoniam homini dicentur quae sola uersantur circa litteras diuinas et organa spiritus san cti in qu'bus est uita et ueritas.  \pend\pstart Mente corruptos quis dicet eos qui ex uerbo dei, cui toti incumbunt digladiantur, tanto zelo contra luxum, quaestum ecclesiae satraparum, contra tot pu blicos, notissimos, pestilentissimosque abusus corum  \pend
\section*{COMMEN. IN PRIOREM }
\marginpar{[ p.160 ]}\pstart qui sibi arrogant regimen ecclesiae, quam conculcant uerbo et exemplo et opere, quomodo correpti men te qui fidem inculcant, spiritum erigunt, charitatem persuadent, opera uere bona docent, commendantur et urgent fortissime.  \pend\pstart Et qui a ueritate priuati sunt T Quibus adempta est ueritas. Quid enim iustius eis contingere potuit, quam quod ablata est ab eis ueritas tam indigne tractata, tot dubiis inuoluta, tot difficultatibus conspurcata, et omnia sua specie et gracia destituta, ut ludibrio haberi coeperit, Christus, et ueritas, et scriptura sacra? Quis enim eam aliquot annorum cen turiis digne tractauit? Quis linguas inuestigauitsine quibus nihil agitur in sacris? Sed defuerunt libri, prae ceptores, imo studiosi sacrarum rerum deerant, dum humana omnes et uana mirantur, dum honores et diuitias ambiunt, et fides non inuenitur in mundo, dum nemo audet contradioere magistris nostris, dum ueritatis Apostoli concremanttir, dum homo peauti sedem Christi ocupat, et fastus supprimit Christi humilitatem, traditiones humanae legem diuinam et caetera multa.  \pend
\phantomsection
\addcontentsline{toc}{subsection}{\textit{Existimantium quaestum esse pietatem. }}
\subsection*{\textit{Existimantium quaestum esse pietatem. }}\pstart Dei cultus est pietas, fide, spe, charitate deum o lere, nihil antiquius ducere quam deo plaare, uer\pend
\section*{TIMOTHEI EPISTOLAM. }
\marginpar{[ p.161 ]}\pstart bo dei insistere, meditari die ac nocte in lege domini. Scrutari testimonia dei, promissionesque ; diuinas et toto corde confidere, proximo bene uelle et be nefacere, et ut seipsum diligere, de se ipso diffidere peccata propria plangere, et dei misericordiam implorare et sperare. Et in Christi meritis praesumere In hanc pietatem ueram alii existimauerunt implo rare sanctos, ecclesiarum patronos, et reliquias ornare, missas canere et legere in honore eorum, canto res instituere, horas canonicas exoluere, anniuersaria instituere. Census pro augendo sacerdotum numero emere, lumina incendere, campanas pulsare. Imagines preciosas procurare, offerre argentum, aurum, blada, uinum, fructus. Exstruere templa, dedicare altaria, quis negabit, haec nunc a multis existimare pietatem, quis nesciat hoc esse quaestum? Nonne his modis deus, imo sancti coluntur? Quis ferat hunc cultum dei subuerti? Dixi de quaestu confessionum, et confessionalium et confessorum, indulgentiarum, sacramentorum omnium, ietuniorum, dispensationum, uotorum. Nonne religiosissimus ac dei cultus est mendicantibus fratribus sacerdotibus elee mosynas largiri? Cum eis jepeliri, testamenta confiscere confraternitates emere. Certe haec existimauerunt dei cultum acceptissimum, peccatorum conciliano\pend
\section*{COMMEN. IN PRIOREM }
\marginpar{[ p.162 ]}\pstart nem, amplissima merita, arram beatitudinis certi ssime uiam lalutis et deo gratissimus cultus, quaestus hospi taliorum, pro uariis infirmis. Et quid non in statu clericorum quaestum resipit, nisi si euangeltum praedioeetur et uerbum dei? Nihil de hoc quaestu infernali et satanae aul tu pro nostri temporis abusu satis arbitror dia potest sed uerbo dei interficietur idolum quaestus et mendatii.  \pend
\phantomsection
\addcontentsline{toc}{subsection}{\textit{Est autem questus magnus pietas cum suf. }}
\subsection*{\textit{Est autem questus magnus pietas cum suf. }}\pstart Copiosissime diues est, Croesoque felicior, siquis sit sua sorte contentus: praesentis uitae neoesariis tantum rebus fruitur, et deum colit in spiritu et ueritate, diues in fide et charitate et uerbo dei. Mundum cum sua concupiscentia contemnere, longerendis diuitiis non instare. Multa conquesiuit, multum thesaurizauit qui pietatem obtinuit.  \pend
\phantomsection
\addcontentsline{toc}{subsection}{\textit{Nihil enim intulimus in hunc mudum, haud dubium quia nec auferre quid possumus. }}
\subsection*{\textit{Nihil enim intulimus in hunc mudum, haud dubium quia nec auferre quid possumus. }}\pstart Diuitiae huius mundi, aurum, argentum, agri, possessiones fuerunt in hoc mundo, nondum nobis natis.animam deus creauit, aptauit nobis corpus, praefuit generationi no strae fauor dei, a nobis nihil uel nobilcum intulimus, omniaque relinquemus. Corpus terra mater suscipiet, spiritus ibit ad deum qui dedit illum. Nihil ergo solliciti simus de terrenis, nec etiam in crastinum, multominus quam aues caeli, et qui hodie insatiabilius colligunt di  \pend
\section*{TIMOTHEI EPISTOLAM. }
\marginpar{[ p.163 ]}\pstart uitias huius mundi, et quaestum sectantur et turpe lucrum, quam illi ad quos loquitur hoc loco dei uerbum Christus in Apostolo, regula pietatis. Nihil inferunt nec exportant, sed comportant semper peioribus usibus profuturas. Mundum possident et possidentur a mundo, alieni a Christo. Efferent secum ex hoc mundo con scientiam criminum, dei terribile iuditium, diffamiam ad posteros, animae damnationem, sucessorisque subuer sionem et aertum periclum damnosi luxus.  \pend
\phantomsection
\addcontentsline{toc}{subsection}{\textit{Habentes autem alimenta et quibus tegam. }}
\subsection*{\textit{Habentes autem alimenta et quibus tegam. }}\pstart Cum hoc ad Christianos omnes et Chrisu sugesso res pertineat, tamen id omnium maxime episopos atti net ut usum omnium rerum naturae finibus metiantur, ut corpus alatur, non luxurietur, oneretur, pinguescat, tractetur te nerius, sed ad necessitatemalimenta susapiat Tegatur amictu necessario et minine, uel precioso, uel curioso uel superfluo. Hic locus est neoessario no stris praedicatoribus semper et ante omnia ad hoc nostrum saeculum tractandus, cum uidemus tantum luxum in uictu et amictu, ut nobis ditissimus hic mundus non sufficiat, quomo do nostro saeculo adaucta sunt omnia, augenturque indies sine discrimine, opum uel personarum, in iacturam intollerabllem patriae, euacua tionem totius terrae, fomentum omnium uitiorum, oppressi onem pauperum, tyrannidem diuitum, fidei Chrisliana  \pend
\section*{COMMEN. IN PRIOREM }
\marginpar{[ p.164 ]}\pstart ludibrium, scandalum infidelium. Christi blasphemiam inculcanda iugiter etstrennue nemo non probabit, negabit nemo uera, damnosa, pessima, saunt omnes, detestanur ut malum. Emendat nemo, praecones ecclesiae taoent, arrident, aemulantur. Canes muti, quo rum animae requiruntur, de sandalo iudicabuntur u nt det declamatores idone os?  \pend
\phantomsection
\addcontentsline{toc}{subsection}{\textit{Nam qui uolunt diuites fieri in hoc saeculo, incidunt tentationem, et in laque um diaboli, et desideria multa, inutilia et nociua, quae mergunt homines in interitum et perditionem. }}
\subsection*{\textit{Nam qui uolunt diuites fieri in hoc saeculo, incidunt tentationem, et in laque um diaboli, et desideria multa, inutilia et nociua, quae mergunt homines in interitum et perditionem. }}\pstart Qui uolunt, desiderant, student ditescre. Non loquitur de illis qui non tantopere uolunt, sed diuitiae affluant, cor non apponunt, utuntur tanquam non utantur. Grati deo de donata affluentia, magis quam parta multa cura. Sunt enim quibus parum anniten tibus diuitiae ocurrunt, sunt qui multis laboribus uane nituntur, quos fugit mammona.. Diuites fieri cum pauperes sint, qui uolunt uinaere fata sua, ui persequuntur diuitias fugientes, praestaret tales aequanimiter ferre sortem suam, et ex neoessitate uirtutem assumere: non enim omnibus prosunt.  \pend\pstart In hocsaeculo ) Si in futuro reponere conaren, tur thesauros, ubi nec erug, nec tinea corrumpit  \pend
\section*{TIMOTHEI EPISTOLAM. }
\marginpar{[ p.165 ]}\pstart sancte agerent, tuteque ; uiuerent, si duitiis spiritualibus fide, pacientia, pace et charitate ditesoerent Christi ani essent.  \pend\pstart Inddunt in tentationem diaboli ) Volentes nolentes dei promissione et diaboli machinatione con tra conscientiam, famam, uitam, salutem domini et scriptu ram aliquando moliri contigit eos affectu et deside rio lucri turpis adactos, cum priusquam ditesare oona rentur, ea ut mala maxima abhoruissent, sicut fit plu ralibus in beneficiis, et theologis, si curtisanorum praepositi inaderint. Et bonis uiris si acaidat proles multa, quos uolunt prouchi ad gloriam huius mundi uel ad flatum suo similem, hi statim incipiunt uehemem ter temptari. Sic tentationes sunt maxima pericula maris, uiarum, latronum, praedoum, infidelium et mem dosorum, principum, coaequalium. Quanto magis di tescit, tanto maior conflatur inuidia plurimorum. Figulus enim odit figulum, et mercator comparem. Item deficiunt paulatim in timore dei, laxa conscientia. Spernunt primo dei uerbum, deinde odiunt, tandem persequuntur. Nemo simul deo seruire potest et mammonae.  \pend\pstart Et in laqueum ) Tentationem sequitur trepida tio et casus. Et qui diligit periclitari, facile perit. Qui incaute ambulat, scandalizatur. Mundus ple  \pend
\section*{COMMEN. IN PRIOREM }
\marginpar{[ p.166 ]}\pstart nus laqueis a mum di amatoribus facile illaqueatur, cum iustus uix saluetur, et oculatissimus uix euadet, quid desideriis languens unum pegatum trahit ad aliud, inuol uit inexplicabiliter? Diabolus tentator adest, semper quaerens quem deuoret rugiens leo, magis tamen unusquisque a concupiscentia sua trahitur et illicitur, suaque ; uoluntate et desideriis impuris et terrenis. Ideo Christus paupertatem, id est pecuniae, diuitiarum que contemptum docuit, non tantum spiritu, quae neaessaria in primis est omnibus Christianis, sed et rebus et possessionibus immunes apostolos uoluit et apostolicos ne in earum aliquando appetitum incidentes, qd'facile solet, errarent a fide. Vnde enim tot abusus, tot luxus tot errores in moribus quam auaritia clericorum, qui ubi scriptura sacram sibi aduersari sentiebant, deprauuerunt expositionibus fucatis ad suum quaesum et distincti onibus sucurerunt: inde tot regulae, tot dispensationes tot priuilegia, de quibus alii satis clamant. Patet enim late quomo do cupiditas et amor sui ipsius, fidem exterminauerint de terra, uerisimileque ; est de tempore apostolorum subintrasse pseudoapostolos, ꝗ quaestumma gs q̃ pietatem inqsicrunt. Conqueritur haec saepe apost.  \pend
\phantomsection
\addcontentsline{toc}{subsection}{\textit{Et inseruerunt se doloribus multis. }}
\subsection*{\textit{Et inseruerunt se doloribus multis. }}\pstart Angustas, dolores, sollicitudines, simulationes, fi dtaque uitae ausuritatem, et hypocriticam iusiciam macti  \pend
\section*{TIMOTHEI EPISTOLAM. }
\marginpar{[ p.167 ]}\pstart sunt arripere, qui cupidi pecuniae erani, ditariq̃; uolebant ut hodie. Sic enim placidissime, suauis imcque ́; uiuunt qui sunt sua sorte contenti, incurii rerum mundanarum siait qui sucumbunt affectui colligendarum diuiti arum omnium miserrimi sunt plenque ; perpeiuis anxieta tibus, quibus semetipsos inserunt nemine compellen t, ut sine doloribus mundouiuere non possint, qui Christo plaoere et acquiesoere non curant.  \pend
\phantomsection
\addcontentsline{toc}{subsection}{\textit{Tu autem ô homo dei haec fuge, sectare uero iusticiam, pietatem, fidem charitatem, patientiam, mansuetudinem. }}
\subsection*{\textit{Tu autem ô homo dei haec fuge, sectare uero iusticiam, pietatem, fidem charitatem, patientiam, mansuetudinem. }}\pstart Praedicta agunt mancipia huius mundi, saeculi amato res, et sui ipsius, dite oere cupiunt, quaestum sectantur, ue ritati resistunt, falsitatem et errores praedicant, super bi unt sensu et habitu. Tu uero homo dei, a deo electus donatus gracia multiplici, spiritu dei ornatus, per prophetiam ordinatus, deo totus dicatus, et seruus ec clesiae datus, tu haec omnia fuge quae deo aduersantur, tibi et ecclesiae obsunt. Sectare uero iusticiam, tota uita tua iusta sit et uir tuosa, ut sis sine querela coram hominibus, nemini iniuriam aliquatenus irrogans soli quoq deo plaoere cupiens ut improbitatem omnem causas in opere et opinione, ut sit tibi saluus honor et au toritas et doctrina. Hic iusticiam generaliter inte lligo conplectente uirtutes oës quarum ins gnores receset consequenter.  \pend
\section*{COMMEN. IN PRIOREM }
\marginpar{[ p.168 ]}
\phantomsection
\addcontentsline{toc}{subsection}{\textit{Et desideria multa et inutilia et nociua. }}
\subsection*{\textit{Et desideria multa et inutilia et nociua. }}\pstart Quanto enim magis ditantur, tanto magis deside rant. Facilius est totum deserere, quam desideriis imperare quae subsequitur, cura et sollicitudo inutilis, dum omnia timent, semper attoniti, laeti nunquam, nihil sufficit, nihil satis fortunatum, optauerant plura, iam homines, iam fata incusantur, somniis inquietantur nibil cogitant nisi mammonam, ubi erg seruitium dei? Nociua desideria sunt, quicquid contra legem et cha ritatem cogitatur in communem iacturam, dum priua tim nobis consulimus, dum proximum grauamus etc.  \pend
\phantomsection
\addcontentsline{toc}{subsection}{\textit{Qui mergunt homines in interitum etc. }}
\subsection*{\textit{Qui mergunt homines in interitum etc. }}\pstart Desideria nociua et inordinata dum agrauantur et impellunt ad nimis propense uolendum contra dei legem ad satisfaciendum auaritiae, ad tumultus sae culares obuoluut, et sollicitudinibus inquictis diuit arum tam soopulis alliduntur, quid aliud restat, nsi ut mergantur in profundum auaritiae et intereaxt, diut tiarum mole obruti, et hydropesi desideriorum munda norum, perditi deum negligant, odiunt, uendant et tra dant ut Iudas quem auaritia inuidum, sacrilegum, impium et traditorem ac parricidam reddidit et subuertit. Et hodte uiri statu religiosi, et deum alioqui timem tes, sed auaritia sic a mammona demerguntur, ut Christi amorem et neoessitatem nibili ducant, dum  \pend
\section*{TIMOTHEI EPISTOLAM. }
\marginpar{[ p.169 ]}\pstart augescant census et possessiones, quantumlibet certi sint, quam nihil ualeat omnis uita et conuersatio illo rum, ut in con scientiis suis agnoscant fucata per cos ge ri omnia, deumque uideant offendi ex omni superstitione sua. Auaritia mergit hos homines in perditionem, neglectusque proximi neoessitatis.  \pend
\phantomsection
\addcontentsline{toc}{subsection}{\textit{Radix autem malorum omnium cupiditas quam quidam appetentes errauerunt a fide, et inseruerunt se doloribus multis. }}
\subsection*{\textit{Radix autem malorum omnium cupiditas quam quidam appetentes errauerunt a fide, et inseruerunt se doloribus multis. }}\pstart Omnia mala ex amore sui pullulant, quae est cupi ditas proprii et priuati boni: non quidem ueri, sed praefixi et electi ab errante ratione. Magis tamen hoc loco insinuatur pecuntae studium, quam qui impensius appetunt, facile in multa mala prolabuntur. Ideo sine pera et aere discipulos msit, ideo aurum et argentum Petrus et Ioanes non habent, ideo Simonem cum sua pecu nia damnationi addixit, ideo quaestum interdixit, ali menta et tegumenta sola permisit, omnem luxum et curiosi tatem improbat, quae fidem Christi eneruant et annihilant.  \pend\pstart Pietatem. virtutem qua deo plaoere possis et hominibus: habet enim rememorationem inuita quae nunc agitur et futura. Hanc puto hic dici tractabilitatem ad homines, et timorem filialem ad deumut ad patrem.  \pend\pstart Fidem. Qua tota fiducia ex deo pendeas, ut taetera omnia paruifacias, siue sint prospera, siue ad  \pend
\section*{COMMEN. IN PRIOREM }
\marginpar{[ p.170 ]}\pstart uersa, quicquid egeris in dei gloriam, certus sis deo placere, et de eius bonitate sine cunctatione omne bonum sperare.  \pend\pstart Charitatem. Ad proximo inseruiendum dum possis, affectu imprimis et operis exhbitione insignis, oibus dulci, pio sinceroque corde coniunctus paratusque facere, ut tibi fieri iuste rationabilite que uelis.  \pend\pstart Pacientiam. Que necessaria piis est, quos ma nent psecutiones in hoc mundo, qua et animam posside mus, et Christi exemplo maxime commendatam uiiam ad regnum caelorum, omnibus quoque sanctis familiaris imam.  \pend\pstart Mansuetudinem. vt peccantibus facile in dulgeamus, omnia quae benefieri possunt, in bonam par tem interpraetando. Non reclamando si contra quam placeat quippiam in nos, uel dicatur uel fiat, benefa cere, etia nobis malefacientibus, et in bono malum uin cere, oppressis pro uiribus condolere et succurrere.  \pend
\phantomsection
\addcontentsline{toc}{subsection}{\textit{Certa bonum certamem fidei, apprehende uitam aeternam, ad quam uocatus es, et confessus es bonam confessionem etc. }}
\subsection*{\textit{Certa bonum certamem fidei, apprehende uitam aeternam, ad quam uocatus es, et confessus es bonam confessionem etc. }}\pstart Docui te in primis fidem uera in Christum, ut et tu doceas omnes, et ut in ea et per eam apprehendatur uita aeterna. Hanc confitere semper, et quia illius fidei inimicos plures sustinebis, aduer sarios ueritatis qui per ouseruant:am legis Mofaicae, religionemque die\pend
\section*{TIMOTHEI EPISTOLAM. }
\marginpar{[ p.171 ]}\pstart rum et ciborum uestium et obseruantiarum carnalium deo placere uolunt, et id docent, fidemque ; in Christum ignorant, ideo fortiter et strennue certamen cum eis ineas ac sustineas. Hic enim uincendi sit, ad hanc pugnam conficiendam, uocatus es. Hanc coram me ipso et multis testibus professus es, esse apprehensionem uite aeternae consumataeque ; iusticiae, ut ea sola saluandi sint omnes qui saluabuntur. Si in certami ne isto fidem istam perdideris, similiter quoque aeter nam uitam amiseris. Certa itaque pro fide, apprehende uitam aeternam ad hanc uocatus es. Quaestum magnum asse jueris, si fidei illius pietatem obtinueris et obseruaueris Foris et intra daemonum insultus expe rieris, ut tuam fidem labefactent, terreant, anxiumque a fidu tia in deum deiiciant, sed resiste in fide hac unica munitus, apprehende promissam tibi, et fidelibus uitam aeternam. Vocatus es in uitam aeternam, quam per fidem obtinebis, si pro ea dimiaueris. Confessus es fidem illum, et me docente credidisti et docendam suscepisti, dum Episcopatum idem munus praedican di et instituendi ecclesiam obiuisti. In hac confessione perseueres atque proficias. Quomodo autem uita aeterna apprehendi posset in ea fide, qua deum esse credimus, et articulos fidei non falsos esse fatemur, cumid et diabolus possit et iniqui, sed tota mente  \pend
\section*{COMMEN. IN PRIOREM }
\marginpar{[ p.172 ]}\pstart deo fidere, et fiduciam bonitatis in deum, obedien tia praeceptorum comprobare, uel de fragilitate no stra humiliter ueniam orare et sperare: quis neget esse uitae eternae uiam et arram et ianuam, ex qua nemo excluditur, de qua alias?  \pend
\phantomsection
\addcontentsline{toc}{subsection}{\textit{Praecipio tibi coram deo qui uiuificat omnia, et Christo Iesu, qui testimonium reddidit sub Pontio Pilato bonam confessionem, ut serues praeceptum, sine macula irreprehensibile, usque in aduentum domini nostri lesu Christi. }}
\subsection*{\textit{Praecipio tibi coram deo qui uiuificat omnia, et Christo Iesu, qui testimonium reddidit sub Pontio Pilato bonam confessionem, ut serues praeceptum, sine macula irreprehensibile, usque in aduentum domini nostri lesu Christi. }}\pstart Deus omnia uiuificat, quia in eo uiuimus, mouemur, et sumus: et non solum nos homines, sed cre aturae omnes, quae a deo creatore acceperunt esse, uuere, moueri, operari, et quicquid habent. Cui sub iecta sunt omnia, cui obedire debent omnia, et in gloriam et honorem suam ordinantur. Eidem obe dies, seruies, placebis, si quae de fide in eum praedicanda committo, compleueris. Iesus Christus dominus, cum ante Pontium Pilatum iudicandus astaret, testimonium reddidit ueritati, nec mortis etiam timore retractus, contestatus ueritatemse oppone re, testari: quam et tu, si ut obligaris, praedicaueris attestatusque ; fueris, praeceptum illud meum de fide et uita aeterna praedicanda, diligenter custodies,  \pend
\section*{TIMOTHEI EPISTOLAM. }
\marginpar{[ p.173 ]}\pstart et fidelibus, ut omnium maxime necessariam incul çabis, per quam et iustificentur, et uitam aeternam sine dubio apprehendunt.  \pend\pstart Sine macula ) Prae dicationem fidei et caeterorum quae imposuit discipulo, sine macuila seruat, cum quaestum non admiscet, dum sua non quaerere auditur, dum non tam laudari, quam proficere cupit. Ma culam enim inurunt non solum praedicatori, sed eti am doctrinae: quaestus, ambitio, et insinceritas notatur. Refertur etiam ad ipsum ministrum uerbi, ut non scandalizet opere, quos instruit doctrina. Irre prehensibile uerbum fidei dicitur, cum sacrarum scrip£ turarum autoritate declaratur, ostenditur et robo ratur Ne sensu proprio docere uideatur, quod reprehensionem non euaderet, quod omnis cogitatio et sensus humanus mendatium operentur. Similiter traditiones humanae uidemus sine reprehensione doceri non posse: sunt enim mutabilitati obnoxiae. Lex autem domini irreprehensibilis est, secundum quam tu quoque praedicans irreprehensibilis eris.  \pend\pstart Vsque in aduentum domini nostri Iesu Christi ) Quousq Christi fides, gratia, et gloria, toto orbe clarescat, et agnoscat totus orbis eum dominum et redemptorem, cui deus pater dederit omnia. Qui saluator esse debeat omnium in se credentiumn, qui  \pend
\section*{COMMEN. IN PRIOREM }
\marginpar{[ p.174 ]}\pstart solus colendus et inuocandus, per quem deus oma nia restituit, et recreat orbem, per eum, et in eo, con ciliari deo omnem hominem in hunc mundum uenientem. Haec inquam agnoscere, et confiteri de Chri sto, et ut redemptorem glorificare, inquam regi regum et domino dominantium seruire, uerbo suo credere, in eum sperare, omne bonum de eo consequi, longe nobilior susoeptio Christi aduenientis est, et cum maiestate cultus felicior, quam si rex et princeps ad ueniret in mundum, denuo in carne gloriosus, sicut uenerat passurus, donec aliquando in gloriam maiestatis suae orbem uniuersum, uiuos et mortuos iuste, manifesteque iudicaturus adueniet in fine mundi et subito.  \pend
\phantomsection
\addcontentsline{toc}{subsection}{\textit{Quem suis temporibus ostendet beatus et solus princeps rex regnantium, et dominus dominantium, qui solus habet immortalitatem. }}
\subsection*{\textit{Quem suis temporibus ostendet beatus et solus princeps rex regnantium, et dominus dominantium, qui solus habet immortalitatem. }}\pstart Ostendet deus apparitionem Christi toto orbi, cum Euangeliumsuum ubique praedicabitur, aoeeptabitur, seruabitur, intelligetur. Haec erit nobilis apparitio Christi, cum non iam in carne corruptibili, sed inuirtute et gloria spiritus sui subiiciet sibi omnia. Id autem fiet gratia et dono dei patris, tempore de sinai, et prefinio: omnibus ignorato modo inesf.  \pend
\section*{TIMOTHEI EPISTOLAM. }
\marginpar{[ p.175 ]}\pstart bili, dum minus sperabitur. Subito sicut fulgur ab oriente apparet in occidentem, qua illustrations et iudicio dultus sui iudicabit orbem.  \pend\pstart Beatus ) Vnice uerbo maiestatem ineffabilem balbuciens nominat. Beatus enim, qui omne bonum possidet. Is solus deus. Sicut et sunt beati dei, quo scilicet maior cogitari non potest nec nominari.  \pend\pstart Et solus potens ) Innititur uel paululum aliquod de deo clarius et specialius eloqui et dicit, solum potentem, quod in plus est, quam omnipotens et maius incomparabiliter. Vt non solum non possit om nia, sed et quicquid usquam est, potestatis uel poten tue, id solius dei sit. Nulla enim creatura aliquid po test nisi in, et per deum. Omnia enim in omnibus operatur. Nihil hic distinguitur, non de absoluta, uel ordinata potentia, non de mala uel bona. Omnia enim ipse solus potest, ipse solus operatur nihil malum, quod facit, vidit deus omnia quae fecit, non dicitur ibidem quae creauit, sed quae fecit. Et erant ual de bona. Iple itaque solus omnia potest, nos nihil pos sumus, nisi in eo, sic angeli, sic caeli, sic omnia, sunt, uiuunt, mouentur et operantur in deo: ut nihil po tentiae nobis ascribamus. Peccareuero potentia non est, sed defectus et impotentiae inditium.  \pend
\section*{COMMEN. IN PRIOREM }
\marginpar{[ p.176 ]}\pstart Rex regum ) Ipse rex est, ipse regit, gubernat, moderatur, disponit, prouidet omnia. Qui regnant in mundo, nihil possunt, nisi per deum: nec uelle bo num, nec facere. Cor non solum, sed quicquia uiuit et habet. Rex est in manu dei, et quo uoluerit uer tet illud. Vnicus iste rex regum, deus solus, potens, et faciens omnia, nobis orandus est, ut propitius, clemens, misericors, dirigat actus nostros in beneplacito suo. Et dispositionem dei in hoc mundo hu militer et pacienter amplectamur, quicquid illud tandem fuerit. Interim tamen uoluntatem suam ex uerbo suo discentes, omni studio et uiribus annitamur, ut quae praecepit dominus, haec fiant: tam a no bis quam ab aliis. Haec enim nos debere conari sci mus, deumuelle, et quicquid in manus dederit agendum bonum, hoc instanter operandum, scientes uolun tatem dei nostri, quod si prouidentia sua secus cesse rit, bonum quod instituimus, certi tamen sumus, quod haec nos deus facere uoluerit, licet non perficere. Si peccatum est, pro culpa oremus. Si processit melius, domino acceptum feramus. Et sic rex regumet dominus dominantium, agnoscetur a piis et digne coletur. Si nos et nostros actus iuxta suam uoluntatem pro uirili instituerimus, quas uires, quod uelle et pficere, sua tantum gratia assequamur, nihil nobis.  \pend
\section*{TIMOTHEI EPISTOLAM. }
\marginpar{[ p.177 ]}\pstart Qui solus habet immortalitatem ) Solus est fons uitae et ipsa uita, non potest non esse: a semetipso habebat esse, et omnia alia ab illo. Nunquam desinens, sicut nos incoepit. Aeternus def oeere in semetip so non potest, non destrui, contrarium non habet, non est compositus. Simplicissimus et non potest re solui in partes etc. Habet immortalitatem, quam donare possit, cui uoluerit: et a quo eam sperare, consequi, et habere poterimus, et a nullo alio. Cuius regnum, et inquisitio illius regni, angeli, animae, spiritus immortales sunt, per dei ipsius gratiam, non ex se.  \pend\pstart Qui lucem habitat in acessibilem ) Lux est deus, et tene brae non sint ullae. Nthil ignorat omnia no uit, quae sunt, fuerunt, et esse possunt. Praesciuit, prae ordinauit omnia. Et cum fit maxime cognoscibilis, quia perfectissimumuerum, imo ueritas, tum a se solo cognoscitur, et ab unigenito suo. Et a nulla crea tura agnosci potest, nisi quando, et quantum dignatur se ostendere. Voluntarium enim speculun est Incomprehensibilis est omni creaturae, etiam angelicae. Accedere nemo potest, sed admittere dignatur, et trahere ad se amicos et filios, quos dignos habuerit. Habitat autem luoem, quod sibi sufficit, su apte luce, bonitateque beatus, nullius indiget acci\pend
\section*{COMMEN. IN PRIOREM }
\marginpar{[ p.178 ]}\pstart pere nihil potest. Habitat quiete, et beatus suam lu cem. Diffundit tamen se in omnia, et bonitatem par ticipat omnibus, sicut lux splendorem, et uidetur a remotis, et per reflectiones in creaturis agnoscitur quantum uult, et cui nult, et quam diu.  \pend
\phantomsection
\addcontentsline{toc}{subsection}{\textit{Quem nullus hominum uidit. }}
\subsection*{\textit{Quem nullus hominum uidit. }}\pstart Angelos uiderunt patriarchae et prophetae, et fi guras a deo formatas per quos dei mandatum, ange li uoluntatem dei ad homines multiformiter retulerunt. Ipse autem in essentia a nemine uisus est unquam: nec uideri potest, incor poreus, et purissimus spiritus, imo spiritualiter onni sprritui. Verbum qui dem patris et filius in diuinis, in fine temperum humanandus et carnem exuirgine assumpturus, et mundum liberaturus, et amorem dei commendatu rus humano generi: is in ueteri testamento patribus locutus est, aperuit, promisit, terruit, docuit, mo nuit, omniaque peregit, quae deus cumpatribus egisse describitur. Creasse scilicet et ad Adant, cain, Noë Abraam, Isaac, iacob, Mose, Iosue, Iepte, Gedeon, Samuel, Dauid et prophetas. His omnibus deus lo cutus praescribitur, et a sanctis patribus non sine scripturarum autoritate dicitur dei filium sapientiae et uiam egisse.  \pend
\section*{TIMOTHEI EPISTOLAM. }
\marginpar{[ p.179 ]}
\phantomsection
\addcontentsline{toc}{subsection}{\textit{Sed nec uidere potest. }}
\subsection*{\textit{Sed nec uidere potest. }}\pstart Videre homo non potest, nisi lucidum, et colora tum, et accidentibus obnoxium: imo sola ipsa accidentia, nullam autem substantiam, quae nullo sensu cognosci potest a nobis: sed tantum per accidertia. Deus autem substantia pura et spiritualissima est, et oculus noster instar noctuae ad manifestissima naturae occaecatur. Sed spiritu et amore percipitur deus ab homine, dum eum attraxerit, et sese digna bitur inclinare et fauore suo complecti.  \pend
\phantomsection
\addcontentsline{toc}{subsection}{\textit{Cui honor et imperium sempiternum Amen. }}
\subsection*{\textit{Cui honor et imperium sempiternum Amen. }}\pstart Honor professio bonitatis est. Cum itaque solus de us bonus sit, a quo omnia bona, ideo s. li deo honor et gloria. Peautque uchementer, qui uel honorari cupit de donis dei, que a deo suscepit, cui nou refert acepta, quasi non receperit. Vel qui defert honorem ei, cui non debetur, fraudatque ; eum, a quo bonum est, quod honoratur. Talis apud multos sancto rum honor, qui a se nihil, nisi pecata habent, sola dei gratia bona consecutisunt, quem jolam honorari uoluissent.  \pend\pstart Imperiuml Dei solius est, qui imperat omnibus, qui facit, quicquid uult in caelo, terra, et in omnibus  \pend
\section*{COMMEN. IN PRIOREM }
\marginpar{[ p.180 ]}\pstart abyssis. Cuius motu fiuntomnia, ad quem imperium Augusti caesaris uix est eleuatio capitis minimi uermicult. Nemo sibi contra deum uindicet quicquam, cogitet deum praesidem, quae imperat, cuius est minister, ut sit fidelis familiae commissae, iudicem expectet iustum.  \pend\pstart Amen I Veritas, credendo, quae praedicta sunt, et cum reuerentia et timore dei audienda, ut uere profiteamur nos credere de deo, summum honorem, et summum imperium, nos subiectos et miseros.  \pend
\phantomsection
\addcontentsline{toc}{subsection}{\textit{Diuitibus huius saeculi praecipe, non sublime sapere. }}
\subsection*{\textit{Diuitibus huius saeculi praecipe, non sublime sapere. }}\pstart Collecta pro pauperibus, quae semper nobiscum sunt, nunquam intermittenda. Vult autem praecipue diuitibus praecipi eleemosynae largitionem, quae tum parum ualet, nisi ab humili. I deo primo ad humilita tem hortatur, ut sciant, se nihil habere a seipsis, dona dei esse quae habent, non contemnant pauperes, non moribus et uestitu superbiant: agnescant plura deo debere se, qui plura acceperunt dum honorantur, quiete uiuunt, nullius indigent. Dei recordentur, in cuius manu sunt omnia. Fastus solet occupare diuites, gut macht mut.  \pend
\phantomsection
\addcontentsline{toc}{subsection}{\textit{Nec sperare in incerto diuitiarum. }}
\subsection*{\textit{Nec sperare in incerto diuitiarum. }}\pstart Difficile obtinentur, facile amittuntur Caduca sunt spes non illuc reponenda: fures et fata imminent  \pend
\section*{TIMOTHEI EPISTOLAM. }
\marginpar{[ p.181 ]}\pstart nihil certum in eis. Spes non certa, expectatio futu rorun, non incerta sit caducorum. Excmpla supersunt fugatium diuitiarum, gentes sperant in tempo ralia, sed Christianus in deo uiuo speret. Qui nostra semper curat, cuius filii sumus, negligi non possumus, si eidem confidimus, seruatur spe, obtinetur fide. Nemo in mortuis confidat, uel sanctis, de quibus nihil nobis deus commisit, deus solus anchora spei nostrae per Christum, non in opera bona confidendum, de quibus timendum magis, cum mala sint dei iuditiojdiscutienda. Deus autem pater est gratiam, miser icordiam et gloriam pollicitus est, et potest prae stare si ei fidem habuerimus.  \pend
\phantomsection
\addcontentsline{toc}{subsection}{\textit{Et praestet nobis omnia abunde ad fruendum. }}
\subsection*{\textit{Et praestet nobis omnia abunde ad fruendum. }}\pstart Dei dona sunt diuitiae. Ipse donat ut frui cis liceat, et quam diu et affatim. Alioqut etiam multi diuitiis non fruentur nisi deo concedente. Ideo semper dei fauor ambiendus et gratia: hoc autem fit operibus misericordiae. Quales enim erimus pro ximis, talem nobis deum sentiemus. Omnino nihil no stra prudentia, cura, labore, obtinemus: nihil enim proficimus, nisi eo largiente. Timeant semper, ne domino irato, illas possideant: et ut iuste obtinuerint, iuste utantur, iuste dispensent, cogtent sollicite  \pend
\section*{COMMEN. IN PRIOREM }
\marginpar{[ p.182 ]}\pstart generis ratio exigetur. Difficile est Christianum esse diuitem nisi erogando affatim, quae dominus abunde suppeditat: ut non sit nisi dei dispensator etc.  \pend\pstart Improprie dicitur frui, id est, uti. Non enim propter necessitatem nostram et proximi, habendae diuitiae: non ppter se, solo deo fruendum, Non ergo audiamus, nec obsequamur regulis illorum sophistarum, libere uera loquamur: et magis, quid, quam quibus uerbis explicemus, curandum. Nonut abunde fruamur sed abunde praestat, Parce, pie, pro nebis ut endum est diu tiis, non ad luxumuel fastum: sed ad sobrietatem et simplicitatem.  \pend
\phantomsection
\addcontentsline{toc}{subsection}{\textit{Ad fruendam. }}
\subsection*{\textit{Ad fruendam. }}\pstart Bene agere. Hic quinto docet diuites, ut bene agant, non eaopera curare, quae deus non praece pit, ecclesias, monasteria, anniuersaria, praebendas, altaria, titulos non erigere: sed quae deus praecepit, omni petenti tribuere, pauperes pascere, potare, re dimere, iuuare, consolari, mutuum praestare, nihil inde sperare, erogare bona. Non thesaurizare, non ecclesias uisitare, sed pauperes: non audire missas, sed mendicos: non preculas multas numerare, sed precantes exaudire. Non testamentum ordinare, sed in uita erogare, cum sua sint, et cum dei dispensatorem. agit.  \pend
\section*{TIMOTHEI EPISTOLAM. }
\marginpar{[ p.183 ]}
\phantomsection
\addcontentsline{toc}{subsection}{\textit{Diuites fieri in bonis operibus. }}
\subsection*{\textit{Diuites fieri in bonis operibus. }}\pstart Non ut ditescant bonis corporalibus, sed spiritualibus: non tamen ut fiant diuites, ut praesumant de bonis, sed bona opera a deo mandata, obediendo cum det gratia perficiant. Et se seruos inutiles credant et profi teantur quod facere deb ent, faciait, non spe merced is sed instinctu charitatis et fi duciaein deum. Nom inat misericordiae operibus de ditos, diuites, ꝗ soli scilicet habere diuitias censentur coram deo, cum legittime religioseque administrant, iu xta commissionem dei, cuius procuratores consutuuntur. Vnde ad mentem et admodum loquendi sancti Pauli, diuites fieri inbonis operibus, nihil aliud est, qua m'opera mi sericordiae de diuitiarum abundantia affatim exercere, multiplicare opera uerae pietatis: sicut deus abunde praestitit unde pos sint, ne ex illis aliquando penuriam incidant, ut de diuite epulone, cui gutta aquae negata fuit, diuite rebus, sed egeno bonis operibus et gratia dei  \pend\pstart Bona opera quae dicantur, et sint, hactenus a multis ignoratum. Non nunc disuademus bona opera, sed ca quae hactenus falso putabantur bona prae aliis quae uere bona a deo iubentur.  \pend
\section*{COMMEN. IN PRIOREM }
\marginpar{[ p.184 ]}\pstart Diuites autem erunus, non in hoc saeculo, sed n caelo. Illuc enim thesaurizant, hic semper pauperes erimus: plura debentes quam praestare possimus, semper peccatores, semper orantes pro peccatis remittendis, de operibus bonis nihil praesumentes: sed de misericordia dei et fide nostra. Tum glorie mur, sed non in nobis, quia nec fidem praestare possumus, nec opera nisi deo dante.  \pend
\phantomsection
\addcontentsline{toc}{subsection}{\textit{Facile communicare. }}
\subsection*{\textit{Facile communicare. }}\pstart Non difficulter adigi precationibus ad dandum, non expectare precantem, sed inquirere: non tam suscipere, quam cogere hospitem pauperem. Omnia communia putare, et dei, non sua bona ad dei beneplacitum, non grauatim dispensanda, cum sua non sint, sed eius qui magis indiget.  \pend
\phantomsection
\addcontentsline{toc}{subsection}{\textit{Thesaurizare sibi. }}
\subsection*{\textit{Thesaurizare sibi. }}\pstart In caelum reponere diuitias tuto loco et in aeternum usui nostro profuturas alioqui, uel hic eas perdemus, et alibi: sicut diuiti reponenti in multos annos bona contigit. Quicquid enim non in deum retulerimus, propterque eum egerimus, totum per demus. Sic omnia bona opera nostri causa facta, non dei. Sic eleemo synae, non propter deum factae peribunt. Thesaurus noster enim labitur, si non ad dei laudem oculum applicemus, et ad proximi commoda  \pend
\section*{TIMOTHEI EPISTOLAM. }
\marginpar{[ p.185 ]}\pstart magis quam nostra. De suo praemiabimur, non de nostro Suis bonis fruemur, non nostris. Patitur nobis deus impie loqui, dum pie intelligamus.  \pend
\phantomsection
\addcontentsline{toc}{subsection}{\textit{Fundamentum bonum in futurum, ut apprehendant uitam aeternam. }}
\subsection*{\textit{Fundamentum bonum in futurum, ut apprehendant uitam aeternam. }}\pstart Thesaurum inoertum et temporaneum qui bona col locant, et insinum pauperum reponunt: imo Christo commendant, cui donant, habebunt in futuro stabiles diuitias, aeternas, iucundissimas. Dant enim corporalia, et apprehendunt aeterna, uitamsalicet beatam tanquam dignam compensationem: cum nihil suum dederit, quicquid dedit, nec uoluntatem dandi, nisi a deo habuit, multominus dona et opera. Veruntamen patitur se deus suis propriis bonis absolui, donatque unde soluamus, et unum acupit pro mille. Gra tum susapit, quod gratis concessit. His uerbis quasi adie cticiis causam pauperum egit Paulus, et ut admonea tur diuites ad beneficentia, et pauperes, qui ut uelint, non tamen possint, uerecundia liberentur. Non omnibus praeapit, sed tantum diuitibus, qui citra di spendium proprii uictus et suorum donare possunt. Haec consuetudo collectt rum oblationum, utinam reformaretur, cum sit omnium maxime neoessaria: sic tamen, ut reposita bona, pro pauperibus similiter profundantur ad pauperes, non coagmententur in  \pend
\section*{COMMEN. IN PRIOREM }
\marginpar{[ p.186 ]}\pstart ensus ut fit in hospitalibus, quae continue ditescre student, non tam pro pauperum quam diuitumuju.  \pend
\phantomsection
\addcontentsline{toc}{subsection}{\textit{O Timothee depositum custodi. }}
\subsection*{\textit{O Timothee depositum custodi. }}\pstart Deposue rat penes cum Paulus, imo spiritus sanctus graciam prophetiae, sanaeque ; doctrinae. Nam ab adolescentia sacris litteris operam nauauit, et ube rius donum uerbi, pro ecclesiae profectu aoeperat ut depositum, Hoc autem iul etur sic custodire, ut nen perdatur, non deficiat, id autem est liberaraliter erogare, sine internussione docere, talentum creditum non fodere in terram, sed ad mensuram domino augmentandam exponere. Tunc autem auarissime custoditur, cum liberalissime profunditur.  \pend\pstart Depositum dicit, quia mera dei gracia acepit, sine meritis et illapsu et munere diuini spiritus qua si aliunde donatum.  \pend\pstart O Timothee ) Affectum ostendit, et suscitat, hor tatur, et admonet ceu omnino magnum aliquid et ne cessarium et iugiter cogitandum propositurus, etenim depositum illud saluum custodire, hoc est docere iugiter. Maxime necessarium officium est episcopi, non nostrorum temporum quidem, jed uerorum, qui nunc si sunt paucissimi sunt. Deposita alia habent, pecunias scilicet et equos. Adde quoque si placet et scorta.  \pend
\section*{TIMOTHEI EPISTOLAM. }
\marginpar{[ p.187 ]}
\phantomsection
\addcontentsline{toc}{subsection}{\textit{Deuitans prophanas uocum nouitates. }}
\subsection*{\textit{Deuitans prophanas uocum nouitates. }}\pstart Vult disapulum suum non alia dooere subiectos qua didioerat a megistro, non inconsueta uerba inducere uel confingere. Praesensit iam Apellis, Valentinanorum et aliorum figmenta de aeonbus 4.8. et 3o. sin gula singulis nominibus nominam do, quae iam tempore Apoiu lorum a tudaeis et Graecis supstitiosioribus fingebantur, ex Cabalistarum (ut nune uocant) et phi losophorum, maxime Pythagoricorum dictis se recepis se et edoctos iactabant, de quibus Ireneus in primo et secundo libro Suoeesserunt illis nuper Parrhisien sis theologiae professores, aeque portentuosis nominibus, quae nec Latini neque Graed intelligunt, nisi diu edocantur ad quaestum Praemonuit Apoite lus uitare omnes illas nouas uoces quas in dei uerbis, sanctisque ́ scripturis non reoepimus. Quod si seruatum fuisset, non tot schi mata et differentias doctrinarum audiss emus de homusion, transsubs antione, circumira ssione diuinorum, esentialia, uotionalia, et reliqua melta praeter scripturae sacrae autoritatem introducta Taceo tot respectus reales rationis, dependentias formalitates, qdditates etc. de qbus omissis utilibus tantopere digla diati sunt doctores ecclesiae nostrae, magisque ; perdide runt q̃ repererunt ueritatem, saltem spiritum et fidem amise runt et opauere Christiana sanctimontam et exemplum.  \pend
\section*{COMMEN. IN PRIOREM }
\marginpar{[ p.188 ]}
\phantomsection
\addcontentsline{toc}{subsection}{\textit{Et oppositiones falsi nominis scientiae. }}
\subsection*{\textit{Et oppositiones falsi nominis scientiae. }}\pstart Vocatur quidem ab illis ista cognitio, scientia, sed omnino falso et indigno nomine. Est enim reue ra opinio, stulticia, error et haeresis, superbiam, et quaestum dooens non pietatem. Oppositiones sunt non doctrinae, sed subuersiones discentium et doctrinae. Sic enim eos ui demus nihil aliud scribere, nisi oppo sitiones, et argumenta pro et contra, pro opposito contra oppositum. Vide primum inter mendicantes scholasticum Albertum, quem magistrum uocant in magistrum, et inuenies oppositiones et praeter cas nihil, quem ut sunt inre pulchra sectati omnes. Quos alienos nominare possumus, quos hic Paulus prodiderit hoc uerbo oppositiones falsi nominis scientiae. etiam seipsos damnant et saentiam suam qui caligne oppositionum omnia obuoluunt, ut sincera ueritas apud eos nulla superfuerit. Fides quoque et uera pie tas abolita fuerit, quam eorum nullus pro rei dignitate tractauit. Sed in profundum silentium sepelierunt, donec explosa omni corum schola et falsi nomi nis scientia, det gracia, sacrae litterae, et Christi spiri tus sinoeraque fides denuo emergunt, scholasticique futi ules, cum suo sacrileg Aristotele stulti deprehenduntur.  \pend
\section*{TIMOTHEI EPISTOLAM. }
\marginpar{[ p.189 ]}
\phantomsection
\addcontentsline{toc}{subsection}{\textit{Quam quidam promittentes circa fidem exciderunt. }}
\subsection*{\textit{Quam quidam promittentes circa fidem exciderunt. }}\pstart Non quidem exhibentes etiam eam quam promit tunt scientiam, lioet falso nomine promittunt solum, praestare nequeunt, cum contradictiones euadere non docent, et semper disaipulus magistro reclamat, et diuersum, sentit, dooet et scribit. Disputare discunt omnia inuoluere quaestionibus, cauilla mouere ad omniaut sibi appareant, quae nihil sunt. In hoc uno conueniunt omnes, quod fidem ueram nesasse fateri coguntur Biblia nunquam uidisse, didicisse, nec uel in una sui parte intelligere. Digna animaduersione puniti, ut dum sibi aisternas ueteres fodiunt, non ualentes aquas continere, fontem aquae uiuae perdiderunt Dum humana sectantur, diuina amittunt. Et dum ra tione humana idolum colunt, fidem uerbi dei deseru erunt. Et dum Aristotelis doctores se profitentur, amplissimis titulis, Christi discipuli esse desierunt, et utriusque gratiam et fructum amiserunt. Omnem fidei tacturam et superstitionem, et uanitatum incre mentum istis theologis solis debemus.  \pend\pstart Exciderunt circa fidem ) Aberrarunt a sopo totius sacrae scripturae, quae fidem doct, promissi onesque diuinas in Christo, ideo non mirun, si nihil ap posite interpretantur, a scopo errantes et in ianua,  \pend
\section*{COMMEN. IN PRIOREM }
\marginpar{[ p.190 ]}\pstart quo alio delaberentur, nisi ad opera? Et ea denique  non illa quae scriptura commendat, et deo placent, Jed quae contra, et secus, et praeter uerbum dei. Inde traditionum ingluuies et abusionum abyssus, et mare errorum omnium.  \pend
\phantomsection
\addcontentsline{toc}{subsection}{\textit{Gracia tecum. }}
\subsection*{\textit{Gracia tecum. }}\pstart Haec manu Pauli ascripta dicuntur, cum alia per dictatum describi fecerit, iuxta illius temporis consuetudinem et morem Apostoli. Optima imprecatio. Quid molestiae non tollerabimus, si dei gratia ad sit, et in cam speremus? Si deus pro nobis, quis con tra nos? Facessat omnis humana, gratia dum adest diuina. Amen.  \pend\pstart FINIS COMMENTARIORVM IN PRIOREM PAVLI AD TIMOTHEVM EPISTOLAM.  \pend\pstart BASILEA AE EXCVDEBAT HENRICVS PETRVS MENSE MARTUO AN. M. D. XXXIII.  \pend
\endnumbering
\end{pages}
\end{document}
        